% %%%%%%%%%%%%%%%%%%%%%%%%%%%%%%%%%%%%%%%%%%%%%%%%%%%%%%%%%%%%%%%%%%%%%%% %
%                                                                         %
% The Project Gutenberg EBook of Lectures on Elementary Mathematics, by   %
% Joseph Louis Lagrange                                                   %
%                                                                         %
% This eBook is for the use of anyone anywhere at no cost and with        %
% almost no restrictions whatsoever.  You may copy it, give it away or    %
% re-use it under the terms of the Project Gutenberg License included     %
% with this eBook or online at www.gutenberg.org                          %
%                                                                         %
%                                                                         %
% Title: Lectures on Elementary Mathematics                               %
%                                                                         %
% Author: Joseph Louis Lagrange                                           %
%                                                                         %
% Translator: Thomas Joseph McCormack                                     %
%                                                                         %
% Release Date: July 6, 2011 [EBook #36640]                               %
% Most recently updated: June 11, 2021                         %
%                                                                         %
% Language: English                                                       %
%                                                                         %
% Character set encoding: UTF-8                                           %
%                                                                         %
% *** START OF THIS PROJECT GUTENBERG EBOOK LECTURES ON ELEMENTARY MATHEMATICS ***
%                                                                         %
% %%%%%%%%%%%%%%%%%%%%%%%%%%%%%%%%%%%%%%%%%%%%%%%%%%%%%%%%%%%%%%%%%%%%%%% %

\def\ebook{36640}
%%%%%%%%%%%%%%%%%%%%%%%%%%%%%%%%%%%%%%%%%%%%%%%%%%%%%%%%%%%%%%%%%%%%%%
%%                                                                  %%
%% Packages and substitutions:                                      %%
%%                                                                  %%
%% book:     Required.                                              %%
%% inputenc: Latin-1 text encoding. Required.                       %%
%%                                                                  %%
%% babel:    Greek snippets. Required.                              %%
%%                                                                  %%
%% ifthen:   Logical conditionals. Required.                        %%
%%                                                                  %%
%% amsmath:  AMS mathematics enhancements. Required.                %%
%% amssymb:  Additional mathematical symbols. Required.             %%
%%                                                                  %%
%% alltt:    Fixed-width font environment. Required.                %%
%% array:    Enhanced tabular features. Required.                   %%
%%                                                                  %%
%% footmisc: Start footnote numbering on each page. Required.       %%
%%                                                                  %%
%% multicol: Twocolumn environment for index. Required.             %%
%% makeidx:  Indexing. Required.                                    %%
%%                                                                  %%
%% caption:  Caption customization. Required.                       %%
%% graphicx: Standard interface for graphics inclusion. Required.   %%
%%                                                                  %%
%% calc:     Length calculations. Required.                         %%
%%                                                                  %%
%% fancyhdr: Enhanced running headers and footers. Required.        %%
%%                                                                  %%
%% geometry: Enhanced page layout package. Required.                %%
%% hyperref: Hypertext embellishments for pdf output. Required.     %%
%%                                                                  %%
%%                                                                  %%
%% Producer's Comments:                                             %%
%%                                                                  %%
%%   OCR text for this ebook was obtained on June 24, 2011, from    %%
%%   http://www.archive.org/details/lecturesonelemen00lagruoft.     %%
%%                                                                  %%
%%   Minor changes to the original are noted in this file in three  %%
%%   ways:                                                          %%
%%     1. \Typo{}{} for typographical corrections, showing original %%
%%        and replacement text side-by-side.                        %%
%%     2. \Add{} for inconsistent/missing punctuation.              %%
%%     3. [** TN: Note]s for lengthier or stylistic comments.       %%
%%                                                                  %%
%%                                                                  %%
%% Compilation Flags:                                               %%
%%                                                                  %%
%%   The following behavior may be controlled by boolean flags.     %%
%%                                                                  %%
%%   ForPrinting (false by default):                                %%
%%   If true, compile a print-optimized PDF file: Larger text block,%%
%%   two-sided layout, US Letter paper, black hyperlinks. Default:  %%
%%   screen optimized file (one-sided layout, blue hyperlinks).     %%
%%                                                                  %%
%%                                                                  %%
%% PDF pages: 181 (if ForPrinting set to false)                     %%
%% PDF page size: 4.75 x 7"                                         %%
%% PDF bookmarks: created, point to ToC entries                     %%
%% PDF document info: filled in                                     %%
%% Images: 1 jpg, 6 png diagrams                                    %%
%%                                                                  %%
%% Summary of log file:                                             %%
%% * One over-full and two under-full hboxes; no visual issues.     %%
%%                                                                  %%
%% Compile History:                                                 %%
%%                                                                  %%
%% July, 2011: adhere (Andrew D. Hwang)                             %%
%%             texlive2007, GNU/Linux                               %%
%%                                                                  %%
%% Command block:                                                   %%
%%                                                                  %%
%%     pdflatex x2                                                  %%
%%     makeindex                                                    %%
%%     pdflatex x2                                                  %%
%%                                                                  %%
%%                                                                  %%
%% July 2011: pglatex.                                              %%
%%   Compile this project with:                                     %%
%%   pdflatex 36640-t.tex ..... TWO times                           %%
%%   makeindex 36640-t.idx                                          %%
%%   pdflatex 36640-t.tex ..... TWO times                           %%
%%                                                                  %%
%%   pdfTeXk, Version 3.141592-1.40.3 (Web2C 7.5.6)                 %%
%%                                                                  %%
%%%%%%%%%%%%%%%%%%%%%%%%%%%%%%%%%%%%%%%%%%%%%%%%%%%%%%%%%%%%%%%%%%%%%%
\listfiles
\documentclass[12pt]{book}[2005/09/16]

%%%%%%%%%%%%%%%%%%%%%%%%%%%%% PACKAGES %%%%%%%%%%%%%%%%%%%%%%%%%%%%%%%
\usepackage[utf8]{inputenc}[2006/05/05]

\usepackage[greek,english]{babel}

\usepackage{ifthen}[2001/05/26]  %% Logical conditionals

\usepackage{amsmath}[2000/07/18] %% Displayed equations
\usepackage{amssymb}[2002/01/22] %% and additional symbols

\usepackage{alltt}[1997/06/16]   %% boilerplate, credits, license
\usepackage{array}[2005/08/23]   %% extended array/tabular features

\usepackage[perpage,symbol]{footmisc}[2005/03/17]

\usepackage{multicol}[2006/05/18]
\usepackage{makeidx}[2000/03/29]

\usepackage[font=footnotesize,aboveskip=0pt,labelformat=empty]{caption}[2007/01/07]
\usepackage{graphicx}[1999/02/16]%% For diagrams

\usepackage{calc}[2005/08/06]

\usepackage{fancyhdr} %% For running heads

%%%%%%%%%%%%%%%%%%%%%%%%%%%%%%%%%%%%%%%%%%%%%%%%%%%%%%%%%%%%%%%%%
%%%% Interlude:  Set up PRINTING (default) or SCREEN VIEWING %%%%
%%%%%%%%%%%%%%%%%%%%%%%%%%%%%%%%%%%%%%%%%%%%%%%%%%%%%%%%%%%%%%%%%

% ForPrinting=true                     false (default)
% Asymmetric margins                   Symmetric margins
% 1 : 1.62 text block aspect ratio     3 : 4 text block aspect ratio
% Black hyperlinks                     Blue hyperlinks
% Start major marker pages recto       No blank verso pages
%
\newboolean{ForPrinting}

%% UNCOMMENT the next line for a PRINT-OPTIMIZED VERSION of the text %%
%\setboolean{ForPrinting}{true}

%% Initialize values to ForPrinting=false
\newcommand{\Margins}{hmarginratio=1:1}     % Symmetric margins
\newcommand{\HLinkColor}{blue}              % Hyperlink color
\newcommand{\PDFPageLayout}{SinglePage}
\newcommand{\TransNote}{Transcriber's Note}
\newcommand{\TransNoteCommon}{%
  The camera-quality files for this public-domain ebook may be
  downloaded \textit{gratis} at
  \begin{center}
    \texttt{www.gutenberg.org/ebooks/\ebook}.
  \end{center}

  This ebook was produced using OCR text provided by the University of
  Toronto Gerstein Library through the Internet Archive.
  \bigskip

  Minor typographical corrections and presentational changes have been
  made without comment.
  \bigskip
}

\newcommand{\TransNoteText}{%
  \TransNoteCommon

  This PDF file is optimized for screen viewing, but may easily be
  recompiled for printing. Please consult the preamble of the \LaTeX\
  source file for instructions and other particulars.
}
%% Re-set if ForPrinting=true
\ifthenelse{\boolean{ForPrinting}}{%
  \renewcommand{\Margins}{hmarginratio=2:3} % Asymmetric margins
  \renewcommand{\HLinkColor}{black}         % Hyperlink color
  \renewcommand{\PDFPageLayout}{TwoPageRight}
  \renewcommand{\TransNote}{Transcriber's Note}
  \renewcommand{\TransNoteText}{%
    \TransNoteCommon

    This PDF file is optimized for printing, but may easily be
    recompiled for screen viewing. Please consult the preamble
    of the \LaTeX\ source file for instructions and other particulars.
  }
  % Marginal notes omitted in screen version; need these only if ForPrinting
  \setlength{\marginparwidth}{1in}%
  \setlength{\marginparsep}{12pt}%
}{% If ForPrinting=false, don't skip to recto
  \renewcommand{\cleardoublepage}{\clearpage}
}
%%%%%%%%%%%%%%%%%%%%%%%%%%%%%%%%%%%%%%%%%%%%%%%%%%%%%%%%%%%%%%%%%
%%%%  End of PRINTING/SCREEN VIEWING code; back to packages  %%%%
%%%%%%%%%%%%%%%%%%%%%%%%%%%%%%%%%%%%%%%%%%%%%%%%%%%%%%%%%%%%%%%%%

\ifthenelse{\boolean{ForPrinting}}{%
  \setlength{\paperwidth}{8.5in}%
  \setlength{\paperheight}{11in}%
% ~1:1.67
  \usepackage[body={5in,8in},\Margins]{geometry}[2002/07/08]
}{%
  \setlength{\paperwidth}{4.75in}%
  \setlength{\paperheight}{7in}%
  \raggedbottom
% ~3:4
  \usepackage[body={4.5in,6in},\Margins,includeheadfoot]{geometry}[2002/07/08]
}

\providecommand{\ebook}{00000}    % Overridden during white-washing
\usepackage[pdftex,
  hyperref,
  hyperfootnotes=false,
  pdftitle={The Project Gutenberg eBook \#\ebook: Lectures on Elementary Mathematics},
  pdfauthor={Joseph Louis LaGrange},
  pdfkeywords={University of Toronto, The Internet Archive, Thomas J. McCormack, Andrew D. Hwang},
  pdfstartview=Fit,    % default value
  pdfstartpage=1,      % default value
  pdfpagemode=UseNone, % default value
  bookmarks=true,      % default value
  linktocpage=false,   % default value
  pdfpagelayout=\PDFPageLayout,
  pdfdisplaydoctitle,
  pdfpagelabels=true,
  bookmarksopen=true,
  bookmarksopenlevel=0,
  colorlinks=true,
  linkcolor=\HLinkColor]{hyperref}[2007/02/07]


%% Fixed-width environment to format PG boilerplate %%
\newenvironment{PGtext}{%
\begin{alltt}
\fontsize{8.1}{9}\ttfamily\selectfont}%
{\end{alltt}}

%% Miscellaneous global parameters %%
% No hrule in page header
\renewcommand{\headrulewidth}{0pt}

% For extra-loose spacing in catalog and narrow ToC environments
\newcommand{\Loosen}{\spaceskip0.5em plus 0.25em minus 0.25em}

% Globally loosen the spacing
\setlength{\emergencystretch}{1em}

% Crudely add a bit more space after \hlines
\setlength{\extrarowheight}{1pt}

% Scratch pad for length calculations
\newlength{\TmpLen}

%% Running heads %%
\newcommand{\FlushRunningHeads}{\clearpage\fancyhf{}\cleardoublepage}
\newcommand{\InitRunningHeads}{%
  \setlength{\headheight}{15pt}
  \pagestyle{fancy}
  \thispagestyle{empty}
  \ifthenelse{\boolean{ForPrinting}}
             {\fancyhead[RO,LE]{\thepage}}
             {\fancyhead[R]{\thepage}}
}

\newcommand{\SetRunningHeads}[1]{%
  \fancyhead[C]{\textsc{\MakeLowercase{#1}}}
}

\newcommand{\BookMark}[2]{\phantomsection\pdfbookmark[#1]{#2}{#2}}

%% ToC formatting %%
\AtBeginDocument{\renewcommand{\contentsname}%
  {\protect\thispagestyle{empty}%
    \protect\centering\normalfont\large CONTENTS.}}

\newcommand{\TableofContents}{%
  \FlushRunningHeads
  \InitRunningHeads
  \SetRunningHeads{Contents.}
  \BookMark{0}{Contents.}
  \tableofcontents
}

% [** TN: Original ToC has "PAGES" printed at top right of each page; omitted.]
% For internal bookkeeping
\newboolean{ToCNeedDash}  %\ToCNote units are separated by dashes

%\ToCSect{Title}{xref}
\newcommand{\ToCSect}[2]{%
  \smallskip%
  \settowidth{\TmpLen}{9999}%
  \noindent\strut\parbox[b]{\textwidth-\TmpLen}{\small%
    \scshape\hangindent2em #1\dotfill}%
  \makebox[\TmpLen][r]{\pageref{#2}}%
}

% \Lecture, \Appendix macros control group formatting
% Enclosing environment for ToC headings generated by marginal notes
\newenvironment{ToCnarrower}{%
  \begin{list}{}{%
      \setlength{\parskip}{0pt}%
      \setlength{\leftmargin}{3em}%
      \setlength{\parindent}{0pt}%
      \settowidth{\TmpLen}{9999}%
      \setlength{\rightmargin}{\TmpLen}%
    }\item[]\Loosen\ignorespaces%
  }{%
  \end{list}
}

% And the actual marginal note entries
% \ToCNote{Title}{Number}
\newcommand{\ToCNote}[2]{%
  \ifthenelse{\boolean{ToCNeedDash}}{\ --- }{\setboolean{ToCNeedDash}{true}}%
  \hyperref[#2]{#1}%
  \ignorespaces
}

%% Major document divisions %%
\newcommand{\PGBoilerPlate}{%
  \pagenumbering{Alph}
  \pagestyle{empty}
%  \BookMark{-1}{Front Matter.}
  \BookMark{0}{PG Boilerplate.}
}
\newcommand{\FrontMatter}{%
  \cleardoublepage % pagestyle still empty; \Preface calls \pagestyle{fancy}
  \frontmatter
  \BookMark{-1}{Front Matter.}
}
\newcommand{\MainMatter}{%
  \FlushRunningHeads
  \InitRunningHeads
  \mainmatter
  \BookMark{-1}{Main Matter.}
}
\newcommand{\BackMatter}{%
  \FlushRunningHeads
  \InitRunningHeads
  \backmatter
  \BookMark{-1}{Back Matter.}
}
\newcommand{\PGLicense}{%
  \FlushRunningHeads
  \pagenumbering{Roman}
  \InitRunningHeads
  \BookMark{-1}{PG License.}
  \SetRunningHeads{Licensing.}
}

%% Index formatting %%
\newcommand{\FN}[1]{\hyperpage{#1}~footnote}
\newcommand{\EtSeq}[1]{\hyperpage{#1}~et~seq.}
%[** TN: Added word "also"]
\newcommand{\See}[2]{see also~\textit{#1}}

\makeindex
\makeatletter
\renewcommand{\@idxitem}{\par\hangindent 30\p@\global\let\idxbrk\nobreak}
\renewcommand\subitem{\idxbrk\@idxitem \hspace*{12\p@}\let\idxbrk\relax}
\renewcommand{\indexspace}{\par\penalty-3000 \vskip 10pt plus5pt minus3pt\relax}

\renewenvironment{theindex}{%
  \setlength\columnseprule{0.5pt}\setlength\columnsep{18pt}%
  \phantomsection\label{index}
  \addtocontents{toc}{\protect\ToCSect{Index}{index}}
  \InitRunningHeads
  \SetRunningHeads{Index.}
  \BookMark{0}{Index.}
  \begin{multicols}{2}[\SectTitle{Index.}\small]% ** N.B. font size
    \setlength\parindent{0pt}\setlength\parskip{0pt plus 0.3pt}%
    \thispagestyle{empty}\let\item\@idxitem\raggedright%
  }{%
  \end{multicols}\FlushRunningHeads
}
\makeatother

%% Sectional units %%
\newcommand{\SectTitle}[2][\large]{%
  \section*{\centering\normalfont#1\MakeUppercase{#2}}
}
\newcommand{\SectSubtitle}[2][\normalsize]{%
  \subsection*{\centering\normalfont#1\MakeUppercase{#2}}
}

% \Chapter[PDF name]{Number.}{Heading title}
\newcommand{\Lecture}[3][]{%
  \FlushRunningHeads
  \InitRunningHeads
  \ifthenelse{\equal{#1}{}}{%
    \SetRunningHeads{#3}%
  }{%
    \SetRunningHeads{#1}%
  }
  \BookMark{0}{Lecture #2 #3}%
  \label{lecture:#2}%
  \thispagestyle{empty}
  \ifthenelse{\not\equal{#2}{I.}}{% End ToC entry block of previous chapter
    \addtocontents{toc}{\protect\end{ToCnarrower}}%
  }{}
  \addtocontents{toc}{%
    \protect\ToCSect{Lecture #2\protect\quad #3}{lecture:#2}}
  \addtocontents{toc}{%
    \protect\settowidth{\TmpLen}{9999}\protect\addtolength{\TmpLen}{3em}}%
  \addtocontents{toc}{\protect\begin{ToCnarrower}}%
  \SectTitle{Lecture #2}
  \SectSubtitle{#3}
}

\newcommand{\Preface}{%
  \normalsize
  \FlushRunningHeads
  \pagestyle{fancy}
  \InitRunningHeads
  \SetRunningHeads{Preface.}
  \BookMark{0}{Preface.}
  \label{preface}
  \addtocontents{toc}{\protect\ToCSect{Preface}{preface}}
  \SectTitle{Preface.}
}

\newcommand{\BioSketch}[2]{%
  \FlushRunningHeads
  \InitRunningHeads
  \SetRunningHeads{Biographical Sketch.}
  \BookMark{0}{Biographical Sketch.}
  \label{biosketch}
  \addtocontents{toc}{\protect\ToCSect{Biographical Sketch of #1}{biosketch}}
  \SectTitle{#1}
  \SectSubtitle{#2}
}

\newcommand{\Appendix}[1]{%
  \FlushRunningHeads
  \InitRunningHeads
  \SetRunningHeads{Appendix.}
  \BookMark{0}{Appendix.}
  \label{appendix}
  \addtocontents{toc}{\protect\end{ToCnarrower}}% Close chapter subunit block
  \addtocontents{toc}{\protect\ToCSect{Appendix}{appendix}}
  \addtocontents{toc}{\protect\begin{ToCnarrower}}
  \addtocontents{toc}{\protect\ToCNote{#1}{appendix}}
  \addtocontents{toc}{\protect\end{ToCnarrower}}
  \SectTitle{Appendix.}
  \SectSubtitle{#1}
}

\newcommand{\Signature}[2]{%
  \medskip%
  \null\hfill\textsc{#1}\hspace*{\parindent} \\
  \hspace*{\parindent}#2%
}

\newcounter{MNote}
\newcommand{\MNote}[1]{%
  \refstepcounter{MNote}%
  \phantomsection\label{note:\theMNote}%
  \ifthenelse{\boolean{ForPrinting}}{%
    %marginal note
    \marginpar{\raggedright\footnotesize#1}%
  }{}% Nothing
  \addtocontents{toc}{\protect\ToCNote{#1}{note:\theMNote}}%
  \ignorespaces%
}

%% Illustrations %%
\newcommand{\Frontispiece}{%
  \ifthenelse{\boolean{ForPrinting}}{%
    \cleardoublepage % Place verso, opposite the title page
    \null
    \newpage
  }{}% Else do nothing
  \BookMark{0}{Frontispiece.}
  \null\vfill
  \begin{figure}[hp!]
    \centering
    \ifthenelse{\boolean{ForPrinting}}{%
      \includegraphics[width=\textwidth]{./images/lagrange.jpg}
    }{%
      \includegraphics[width=0.9\textwidth]{./images/lagrange.jpg}
    }
  \end{figure}
  \vfill
  \cleardoublepage
}
% \Figure{Number}{width}
\newcommand{\Figure}[2]{%
  \begin{figure}[hbt!]
    \centering
    \includegraphics[width=#2]{./images/fig#1.png}
    \caption{Fig.~#1.}
  \end{figure}\ignorespaces%
}

%% Book Catalogs %%
\newcommand{\CatalogSmallFont}{%
  \ifthenelse{\boolean{ForPrinting}}{\footnotesize}{\scriptsize}%
}
% Catalog at front of book
\newcommand{\FrontCatalog}[1]{%
  \newpage
  \thispagestyle{empty}
  \SectTitle{#1}
}
\newcommand{\Book}[1]{%
  \medskip\par\noindent\CatalogSmallFont\Loosen\hangindent 2em#1%
}

% and at back
\newcommand{\Catalog}{%
  \FlushRunningHeads
  \InitRunningHeads
  \fancyhf{}
  \BookMark{0}{Catalogue.}
  \begin{center}
    \Large CATALOGUE OF PUBLICATIONS \\[12pt]
    \footnotesize OF THE \\[12pt]
    \large OPEN COURT PUBLISHING CO.
  \end{center}
  \tb
}
\newenvironment{Author}[1]{\medskip\par\noindent #1}{}
\newcommand{\Title}[3][4\parindent]{%
\par\footnotesize\hangindent3\parindent#2%

\ifthenelse{\not\equal{#3}{}}{%
    \hspace*{\parindent}\CatalogSmallFont\hangindent#1 #3\par%
  }{}
}

\newcommand{\Item}[1]{\makebox[1em][r]{#1}\ \hangindent4em}

%% Corrections. %%
\newcommand{\Typo}[2]{#2}
\newcommand{\Add}[1]{\Typo{}{#1}}

%% Page separators and cross-references %%
\newcommand{\PageSep}[1]{\ignorespaces}

\newcommand{\PgLabel}[1]{\phantomsection\label{pg#1}}
\newcommand{\PgRef}[1]{\hyperref[pg#1]{p.~\pageref*{pg#1}}}
\newcommand{\PgRange}[2]{%
  \ifthenelse{\equal{\pageref{pg#1}}{\pageref{pg#2}}}{%
    \hyperref[pg#1]{p.~\pageref*{pg#1}}%
  }{%
    pp.~\hyperref[pg#1]{\pageref*{pg#1}}--\hyperref[pg#2]{\pageref*{pg#2}}%
  }%
}

%% Miscellaneous textual formatting %%
\newcommand{\First}[1]{\textsc{\large #1}}
\newcommand{\ieme}{\textsuperscript{me}}

% Decorative breaks
\newcommand{\tb}[1][0.75in]{\begin{center}\rule{#1}{0.5pt}\end{center}}
\newcommand{\stars}{%
\begin{center}
  \makebox[1in][c]{
    \raisebox{-0.5ex}{*}\hfill\raisebox{0.5ex}{*}\hfill\raisebox{-0.5ex}{*}%
  }
\end{center}
}

%% Miscellaneous mathematical formatting %%
\DeclareMathSizes{12}{11}{8}{7}
\DeclareInputMath{183}{\cdot}

\newcommand{\PadTo}[3][c]{%
  \settowidth{\TmpLen}{\ensuremath{#2}}%
  \makebox[\TmpLen][#1]{\ensuremath{#3}}%
}

\newcommand{\Tag}[1]{%
  \tag*{\ensuremath{#1}}
}

% Square roots of matching height
\newcommand{\mysqrt}[1]{\sqrt{\vphantom{b}#1}}
\newcommand{\sqrta}{\mysqrt{a}}
\newcommand{\sqrtc}{\mysqrt{c}}

% Multiplication row for table on page 30
\newcommand{\MultRow}[2]{#1\,\smash{\rule[-5pt]{0.5pt}{15pt}}&#2}

%%%%%%%%%%%%%%%%%%%%%%%% START OF DOCUMENT %%%%%%%%%%%%%%%%%%%%%%%%%%
\begin{document}
%% PG BOILERPLATE %%
\PGBoilerPlate
\begin{center}
\begin{minipage}{\textwidth}
\small
\begin{PGtext}
The Project Gutenberg EBook of Lectures on Elementary Mathematics, by 
Joseph Louis Lagrange

This eBook is for the use of anyone anywhere at no cost and with
almost no restrictions whatsoever.  You may copy it, give it away or
re-use it under the terms of the Project Gutenberg License included
with this eBook or online at www.gutenberg.org


Title: Lectures on Elementary Mathematics

Author: Joseph Louis Lagrange

Translator: Thomas Joseph McCormack

Release Date: July 6, 2011 [EBook #36640]
Most recently updated: June 11, 2021

Language: English

Character set encoding: UTF-8     

*** START OF THIS PROJECT GUTENBERG EBOOK LECTURES ON ELEMENTARY MATHEMATICS ***
\end{PGtext}
\end{minipage}
\end{center}
\newpage
%% Credits and transcriber's note %%
\begin{center}
\begin{minipage}{\textwidth}
\begin{PGtext}
Produced by Andrew D. Hwang.
\end{PGtext}
\end{minipage}
\end{center}
\vfill

\begin{minipage}{0.85\textwidth}
\small
\BookMark{0}{Transcriber's Note.}
\subsection*{\centering\normalfont\scshape%
\normalsize\MakeLowercase{\TransNote}}%

\raggedright
\TransNoteText
\end{minipage}
%%%%%%%%%%%%%%%%%%%%%%%%%%% FRONT MATTER %%%%%%%%%%%%%%%%%%%%%%%%%%
\FrontMatter
\null\vfill
\noindent {\Large ON ELEMENTARY MATHEMATICS}
\vfill
\PageSep{}
\FrontCatalog{IN THE SAME SERIES.}

\tb

\Book{ON CONTINUITY AND IRRATIONAL NUMBERS, and
ON THE NATURE AND MEANING OF NUMBERS\@.
By R.~\textsc{Dedekind}. From the German by \textit{W.~W. Beman}.
Pages,~115. Cloth, 75~cents net (3s.~6d.~net).}

\Book{GEOMETRIC EXERCISES IN PAPER-FOLDING\@. By \textsc{T.~Sundara Row}.
Edited and revised by \textit{W.~W. Beman} and
\textit{D.~E. Smith}. With many half-tone engravings from photographs
of actual exercises, and a package of papers for
folding. Pages, circa~200. Cloth, \$1.00\Typo{.}{} net (4s.~6d.~net).
(In Preparation.)}

\Book{ON THE STUDY AND DIFFICULTIES OF MATHEMATICS\@.
By \textsc{Augustus De~Morgan}. Reprint edition\Typo{`}{}
with portrait and bibliographies. Pp.,~288. Cloth, \$1.25
net (4s.~6d.~net).}

\Book{LECTURES ON ELEMENTARY MATHEMATICS\@. By
\textsc{Joseph Louis Lagrange}. From the French by \textit{Thomas~J.
McCormack}, With portrait and biography. Pages,~172.
Cloth, \$1.00 net (4s.~6d.~net).}

\Book{ELEMENTARY ILLUSTRATIONS OF THE DIFFERENTIAL
AND INTEGRAL CALCULUS\@. By \textsc{Augustus De~Morgan}.
Reprint edition. With a bibliography of text-books
of the Calculus. Pp.,~144. Price, \$1.00 net (4s.~6d.~net).}

\Book{MATHEMATICAL ESSAYS AND RECREATIONS\@. By
\textsc{Prof.\ Hermann Schubert}, of Hamburg, Germany. From
the German by \textit{T.~J. McCormack}. Essays on Number\Typo{.}{,}
The Magic Square, The Fourth Dimension, The Squaring
of the Circle. Pages,~149. Price, Cloth, 75c.~net (3s.~net).}

\Book{A BRIEF HISTORY OF ELEMENTARY MATHEMATICS\@.
By \textsc{Dr.\ Karl Fink}, of Tübingen. From the German by \textit{W.~W.
Beman} and \textit{D.~E. Smith}. Pp.~333. Cloth, \$1.50 net
(5s.~6d.~net).}

\tb
\vfill

\noindent\makebox[\textwidth][s]{\large THE OPEN COURT PUBLISHING COMPANY}
\begin{center}
\footnotesize 324 DEARBORN ST., CHICAGO. \\
\normalsize LONDON: Kegan Paul, Trench, Trübner \& Co.
\end{center}
\PageSep{i}
%[Blank page]
\PageSep{ii}
\Frontispiece
\PageSep{iii}
\begin{center}
\Large LECTURES\\[24pt]
\footnotesize ON\\[24pt]
\LARGE ELEMENTARY MATHEMATICS
\vfill

\footnotesize BY\\[18pt]
\large JOSEPH LOUIS LAGRANGE
\vfill

\footnotesize FROM THE FRENCH BY\\[18pt]
\normalsize THOMAS J. McCORMACK
\vfill\vfill

\small SECOND EDITION
\vfill\vfill

\large CHICAGO \\
\normalsize THE OPEN COURT PUBLISHING COMPANY \\[12pt]
\footnotesize LONDON AGENTS \\
\textsc{Kegan Paul, Trench, Trübner \& Co., Ltd.} \\
1901
\end{center}
\newpage
\PageSep{iv}
\null\vfill
\begin{center}
\footnotesize TRANSLATION COPYRIGHTED \\
BY \\
\small\textsc{The Open Court Publishing Co.} \\
1898.
\end{center}
\vfill
\PageSep{v}


\Preface

\First{The} present work, which is a translation of the \textit{Leçons élémentaires
sur les \Typo{mathematiques}{mathématiques}} of Joseph Louis Lagrange,
\index{Lagrange, J. L.}%
the greatest of modern analysts, and which is to be found in Volume~VII.
of the new edition of his collected works, consists of a
series of lectures delivered in the year 1795 at the \textit{\Typo{Ecole}{École} Normale},---an
institution which was the direct outcome of the French Revolution
and which gave the first impulse to modern practical
ideals of education. With Lagrange, at this institution, were associated,
as professors of mathematics. Monge and Laplace, and we
\index{Laplace}%
\index{Monge}%
owe to the same historical event the final form of the famous \textit{Géométrie
descriptive}, as well as a second course of lectures on arithmetic
and algebra, introductory to these of Lagrange, by Laplace.

With the exception of a German translation by Niedermüller
\index{Ecole@{\Typo{Ecole}{École} Normale}}%
(Leipsic, 1880), the lectures of Lagrange have never been published
in separate form; originally they appeared in a fragmentary
shape in the \textit{Séances des \Typo{Ecoles}{Écoles} Normales}, as they had been reported
by the stenographers, and were subsequently reprinted in
the journal of the Polytechnic School. From references in them
\index{Polytechnic School}%
to subjects afterwards to be treated it is to be inferred that a fuller
development of higher algebra was intended,---an intention which
the brief existence of the \textit{\Typo{Ecole}{École} Normale} defeated. With very few
exceptions, we have left the expositions in their historical form,
having only referred in an Appendix to a point in the early history
of algebra.

The originality, elegance, and symmetrical character of these
lectures have been pointed out by \Typo{DeMorgan}{De~Morgan}, and notably by Dühring,
\index{DeMorgan@{\Typo{DeMorgan}{De Morgan}}}%
\index{Duhring@{Dühring, E.}}%
who places them in the front rank of elementary expositions,
as an exemplar of their kind. Coming, as they do, from one of
the greatest mathematicians of modern times, and with all the excellencies
which such a source implies, unique in their character
\PageSep{vi}
as a \emph{reading-book} in mathematics, and interwoven with historical
and philosophical remarks of great helpfulness, they cannot fail
to have a beneficent and stimulating influence\Typo{,}{.}

The thanks of the translator of the present volume are due to
Professor Henry~B. Fine, of Princeton, N.~J., for having read the
proofs.

\Signature{Thomas J. McCormack.}
{\textsc{La Salle, Illinois}, August~1, 1898.}
\PageSep{vii}


\BioSketch{Joseph Louis Lagrange.}
{Biographical Sketch.}
\index{Economy of thought}%
\index{Lagrange, J. L.|EtSeq}%
\index{Short-mind symbols|EtSeq}%
\index{Stenophrenic symbols|EtSeq}%
\index{Symbols|EtSeq}%

\First{A great} part of the progress of formal thought, where it is
not hampered by outward causes, has been due to the invention
of what we may call \emph{stenophrenic}, or \emph{short-mind}, symbols.
These, of which all written language and scientific notations are
examples, disengage the mind from the consideration of ponderous
and circuitous mechanical operations and economise its energies
for the performance of new and unaccomplished tasks of thought.
And the advancement of those sciences has been most notable
which have made the most extensive use of these short-mind symbols.
Here mathematics and chemistry stand pre-eminent. The
\index{Greeks, mathematics of the}%
\index{Mathematics!evolution of}%
ancient Greeks, with all their mathematical endowment as a race,
and even admitting that their powers were more visualistic than
analytic, were yet so impeded by their lack of short-mind symbols
as to have made scarcely any progress whatever in analysis. Their
arithmetic was a species of geometry. They did not possess the
sign for zero, and also did not make use of position as an indicator
of value. Even later, when the germs of the indeterminate analysis
were disseminated in Europe by Diophantus, progress ceased
here in the science, doubtless from this very cause. The historical
\index{Science!development of|EtSeq}%
calculations of Archimedes, his approximation to the value of~$\pi$,~etc,
owing to this lack of appropriate arithmetical and algebraical
symbols, entailed enormous and incredible labors, which, if
they had been avoided, would, with his genius, indubitably have
led to great discoveries.
\PageSep{viii}

Subsequently, at the close of the Middle Ages, when the so-called
Arabic figures became established throughout Europe with
the symbol~$0$ and the principle of local value, immediate progress
was made in the art of reckoning. The problems which arose
gave rise to questions of increasing complexity and led up to the
general solutions of equations of the third and fourth degree by
the Italian mathematicians of the sixteenth century. Yet even
these discoveries were made in somewhat the same manner as
problems in mental arithmetic are now solved in common schools;
for the present signs of plus, minus, and equality, the radical and
exponential signs, and especially the systematic use of letters for
denoting general quantities in algebra, had not yet become universal.
The last step was definitively due to the French mathematician
Vieta (1540--1603), and the mighty advancement of analysis
\index{Vieta}%
resulting therefrom can hardly be measured or imagined. The
trammels were here removed from algebraic thought, and it ever
afterwards pursued its way unincumbered in development as if impelled
by some intrinsic and irresistible potency. Then followed
the introduction of exponents by Descartes, the representation of
\index{Descartes}%
geometrical magnitudes by algebraical symbols, the extension of
the theory of exponents to fractional and negative numbers by
Wallis (1616--1703), and other symbolic artifices, which rendered
\index{Wallis}%
the language of analysis as economic, unequivocal, and appropriate
as the needs of the science appeared to demand. In the famous
dispute regarding the invention of the infinitesimal calculus, while
not denying and even granting for the nonce the priority of Newton
\index{Newton, his problem}%
in the matter, some writers have gone so far as to regard Leibnitz's
\index{Leibnitz}%
introduction of the integral symbol~$\int$ as alone a sufficient substantiation
of his claims to originality and independence, so far as the
power of the new science was concerned.

For the \emph{development} of science all such short-mind symbols
are of paramount importance, and seem to carry within themselves
the germ of a perpetual mental motion which needs no outward
power for its unfoldment. Euler's well-known saying that his
\index{Euler}%
\PageSep{ix}
pencil seemed to surpass him in intelligence finds its explanation
here, and will be understood by all who have experienced the uncanny
feeling attending the rapid development of algebraical formulæ,
where the urned thought of centuries, so to speak, rolls from
one's finger's ends.

But it should never be forgotten that the mighty stenophrenic
engine of which we here speak, like all machinery, affords us rather
a mastery over nature than an insight into it; and for some, unfortunately,
the higher symbols of mathematics are merely brambles
that hide the living springs of reality. Many of the greatest
discoveries of science,---for example, those of Galileo, Huygens,
\index{Galileo}%
\index{Huygens}%
and Newton,---were made without the mechanism which afterwards
becomes so indispensable for their development and application.
Galileo's reasoning anent the summation of the impulses imparted
to a falling stone is virtual integration; and Newton's mechanical
discoveries were made by the man who invented, but evidently did
not use to that end, the doctrine of fluxions.
\stars

We have been following here, briefly and roughly, a line of
progressive abstraction and generalisation which even in its beginning
was, psychologically speaking, at an exalted height, but in the
course of centuries had been carried to points of literally ethereal
refinement and altitude. In that long succession of inquirers by
whom this result was effected, the process reached, we may say,
its culmination and purest expression in Joseph Louis Lagrange,
born in Turin, Italy, the 30th~of January,~1736, died in Paris, April~10,
1813. Lagrange's power over symbols has, perhaps, never been
paralleled either before his day or since. It is amusing to hear his
biographers relate that in early life he evinced no aptitude for
mathematics, but seemed to have been given over entirely to the
pursuits of pure literature; for at fifteen we find him teaching
mathematics in an artillery school in Turin, and at nineteen he
had made the greatest discovery in mathematical science since that
of the infinitesimal calculus, namely, the creation of the algorism
\PageSep{x}
\index{Variations, calculus of}%
and method of the Calculus of Variations. ``Your analytical solution
of the isoperimetrical problem,'' writes Euler, then the prince
\index{Euler}%
of European mathematicians, to him, ``leaves nothing to be desired
in this department of inquiry, and I am delighted beyond measure
that it has been your lot to carry to the highest pitch of perfection,
a theory, which since its inception I have been almost the only one
to cultivate.'' But the exact nature of a ``variation'' even Euler
did not grasp, and even as late as~1810 in the English treatise of
Woodhouse on this subject we read regarding a certain new sign
\index{Woodhouse}%
introduced, that M.~Lagrange's ``power over symbols is so unbounded
that the possession of it seems to have made him capricious.''

Lagrange himself was conscious of his wonderful capacities in
this direction. His was a time when geometry, as he himself
phrased it, had become a dead language, the abstractions of analysis
were being pushed to their highest pitch, and he felt that with
his achievements its possibilities within certain limits were being
rapidly exhausted. The saying is attributed to him that chairs of
mathematics, so far as creation was concerned, and unless new
fields were opened up, would soon be as rare at universities as
chairs of Arabic. In both research and exposition, he totally reversed
the methods of his predecessors. They had proceeded in
their exposition from special cases by a species of induction; his
eye was always directed to the highest and most general points of
view; and it was by his suppression of details and neglect of minor,
unimportant considerations that he swept the whole field of analysis
with a generality of insight and power never excelled, adding
to his originality and profundity a conciseness, elegance, and lucidity
which have made him the model of mathematical writers.
\stars

Lagrange came of an old French family of Touraine, France,
said to have been allied to that of Descartes. At the age of twenty-six
he found himself at the zenith of European fame. But his
reputation had been purchased at a great cost. Although of ordinary
\PageSep{xi}
height and well proportioned, he had by his ecstatic devotion
to study,---periods always accompanied by an irregular pulse and
high febrile \Typo{excitatian}{excitation},---almost ruined his health. At this age,
accordingly, he was seized with a hypochondriacal affection and
with bilious disorders, which accompanied him \Typo{thronghout}{throughout} his life,
and which were only allayed by his great abstemiousness and careful
regimen. He was bled twenty-nine times, an infliction which
alone would have affected the most robust constitution. Through
his great care for his health he gave much attention to medicine.
He was, in fact, conversant with all the sciences, although knowing
his \textit{forte} he rarely expressed an opinion on anything unconnected
with mathematics.

When Euler left Berlin for St.~Petersburg in~1766 he and
D'Alembert induced Frederick the Great to make Lagrange president
of the Academy of Sciences at Berlin. Lagrange accepted
the position and lived in Berlin twenty years, where he wrote some
of his greatest works. He was a great favorite of the Berlin people,
and enjoyed the profoundest respect of Frederick the Great,
although the latter seems to have preferred the noisy reputation of
Maupertuis, Lamettrie, and Voltaire to the unobtrusive fame and
personality of the man whose achievements were destined to shed
more lasting light on his reign than those of any of his more strident
literary predecessors: Lagrange was, as he himself said, \textit{philosophe
sans crier}.

The climate of Prussia agreed with the mathematician. He
refused the most seductive offers of foreign courts and princes, and
it was not until the death of Frederick and the intellectual reaction
of the Prussian court that he returned to Paris, where his career
broke forth in renewed splendor. He published in~1788 his great
\textit{Mécanique analytique}, that ``scientific poem'' of Sir William
Rowan Hamilton, which gave the quietus to mechanics as then
formulated, and having been made during the Revolution Professor
of Mathematics at the new \textit{\Typo{Ecole}{École} Normale} and the \textit{\Typo{Ecole}{École} Polytechnique},
\index{Ecole@{\Typo{Ecole}{École} Normale}}%
\index{Polytechnic School}%
he entered with Laplace and Monge upon the activity
\index{Laplace}%
\index{Monge}%
\PageSep{xii}
which made these schools for generations to come exemplars of
practical scientific education, systematising by his lectures there,
and putting into definitive form, the science of mathematical analysis
of which he had developed the extremest capacities. Lagrange's
activity at Paris was interrupted only once by a brief period
of melancholy aversion for mathematics, a lull which he
devoted to the adolescent science of chemistry and to philosophical
studies; but he afterwards resumed his old love with increased ardor
and assiduity. His significance for thought generally is far
beyond what we have space to insist upon. Not least of all, theology,
which had invariably mingled itself with the researches of his
predecessors, was with him forever divorced from a legitimate influence
of science.

The honors of the world sat ill upon Lagrange: \textit{la magnificence
le gênait}, he said; but he lived at a time when proffered
things were usually accepted, not refused. He was loaded with
personal favors and official distinctions by Napoleon who called
\index{Napoleon}%
him \textit{la haute pyramide des sciences mathématiques}, was made a
Senator, a Count of the Empire, a Grand Officer of the Legion of
Honor, and, just before his death, received the grand cross of the
Order of Reunion. He never feared death, which he termed \textit{une
dernière fonction, ni pénible ni désagréable}, much less the disapproval
of the great. He remained in Paris during the Revolution
when \textit{savants} were decidedly in disfavor, but was suspected
of aspiring to no throne but that of mathematics. When Lavoisier
\index{Lavoisier}%
was executed he said: ``It took them but a moment to lay low that
head; yet a hundred years will not suffice perhaps to produce its
like again.''

Lagrange would never allow his portrait to be painted, maintaining
that a man's works and not his personality deserved preservation.
The frontispiece to the present work is from a steel
engraving based on a sketch obtained by stealth at a meeting of
the Institute. His genius was excelled only by the purity and
nobleness of his character, in which the world never even sought
\PageSep{xiii}
to find a blot, and by the exalted Pythagorean simplicity of his
life. He was twice married, and by his wonderful care of his person
lived to the advanced age of seventy-seven years, not one of
which had been misspent. His life was the veriest incarnation of
the scientific spirit; he lived for nothing else. He left his weak
body, which retained its intellectual powers to the very last, as an
offering upon the altar of science, happily made when his work
had been done; but to the world he bequeathed his ``ever-living''
thoughts now recently resurgent in a new and monumental edition
of his works (published by Gauthier-Villars, Paris). \textit{Ma vie est
là!} he said, pointing to his brain the day before his death.

\Signature{Thomas J. McCormack.}{}
\PageSep{xiv}
%[Blank page]
\PageSep{xv}
\TableofContents
\iffalse
%[** TN: Used marginal notes to generate entries; entries in original ToC
% don't obviously match the book's units.]
CONTENTS.

PAGES

Preface

Biographical Sketch of Joseph Louis LaGrange.

Lecture I. On Arithmetic, and in Particular Fractions
and Logarithms. 1-23

Systems of Numeration. Fractions. Greatest Common
Divisor. Continued Fractions. Theory of
Powers, Proportions, and Progressions. Involution
and Evolution. Rule of Three. Interest. Annuities.
Logarithms.

Lecture II. On the Operations of Arithmetic . . . 24-53

Arithmetic and Geometry. New Method of Subtraction.
Abridged and Approximate Multiplication.
Decimals. Property of the Number 9.
Tests of Divisibility. Theory of Remainders.
Checks on Multiplication and Division. Evolution.
Rule of Three. Theory and Practice. Probability
of Life. Alligation or the Rule of Mixtures.

Lecture III. On Algebra, Particularly the Resolution
of Equations of the Third and Fourth Degree 54-95

Origin of Greek Algebra. Diophantus. Indeterminate
Analysis. Equations of the Second Degree.
Translations of Diophantus. Algebra Among the
Arabs. History of Algebra in Italy, France, and
Germany. History of Equations of the Third and
Fourth Degree and of the Irreducible Case. Theory
of Equations. Discussion of Cubic Equations.
Discussion of the Irreducible Case. The Theory
\PageSep{xvi}
of Roots. Extraction of the Square and Cube Roots
of Two Imaginary Binomials. Theory of Imaginary
Expressions. Trisection of an Angle. Method
of Indeterminates. Discussion of Biquadratic Equations.

Lecture IV. On the Resolution of Numerical Equations ... 96-126

Algebraical Resolution of Equations. Numerical
Resolution of Equations. Position of the Roots.
Representation of Equations by Curves. Graphic
Resolution of Equations. Character of the Roots of
Equations. Limits of the Roots of Numerical Equations.
Separation of the Roots. Method of Substitutions.
The Equation of Differences. Method of
Elimination. Constructions and Instruments for
Solving Equations.

Lecture V. On the Employment of Curves in the Solution
of Problems 127-149

Application of Geometry to Algebra. Resolution of
Problems by Curves. The Problem of Two Lights.
Variable Quantities Minimal Values. Analysis
of Biquadratic Equations Conformably to the Problem
of the Two Lights. Advantages of the Method
of Curves The Curve of Errors. \textit{Regula falsi.}
Solution of Problems by the Curve of Errors.
Problem of the Circle and Inscribed Polygon.
Problem of the Observer and Three Objects. Parabolic
Curves. Newton's Problem. Interpolation
of Intermediate Terms in Series of Observations,
Experiments, etc.

Appendix . 151

Note on the Origin of Algebra.
\fi
\PageSep{1}
\MainMatter
\index{Numerical equations|See{Equations}}%


\Lecture[On Arithmetic.]{I.}{On Arithmetic, and in Particular Fractions
and Logarithms.}

\First{Arithmetic} is divided into two parts. The first
is based on the decimal system of notation and
\MNote{Systems of numeration\Add{.}}
on the manner of arranging numeral characters to express
numbers. This first part comprises the four
common operations of addition, subtraction, multiplication,
and division,---operations which, as you
know, would be different if a different system were
adopted, but, which it would not be difficult to transform
from one system to another, if a change of systems
were desirable.

The second part is independent of the system of
\index{Numeration, systems of}%
numeration. It is based on the consideration of quantities
and on the general properties of numbers. The
theory of fractions, the theory of powers and of roots,
the theory of arithmetical and geometrical progressions,
and, lastly, the theory of logarithms, fall under
this head. I purpose to advance, here, some remarks
on the different branches of this part of arithmetic.
\PageSep{2}

It may be regarded as \emph{universal arithmetic}, having
\index{Arithmetic!universal|EtSeq}%
an intimate affinity to algebra. For, if instead of
\index{Algebra!definition of}%
particularising the quantities considered, if instead of
assigning them numerically, we treat them in quite a
general way, designating them by letters, we have
algebra.

You know what a fraction is. The notion of a
\index{Fractions|EtSeq}%
\index{Ratios, constant|EtSeq}%
\MNote{Fractions.}
fraction is slightly more composite than that of whole
numbers. In whole numbers we consider simply a
quantity repeated. To reach the notion of a fraction
it is necessary to consider the quantity divided into a
certain number of parts. Fractions represent in general
ratios, and serve to express one quantity by means
of another. In general, nothing measurable can be
measured except by fractions expressing the result of
the measurement, unless the measure be contained an
exact number of times in the thing to be measured.

You also know how a fraction can be reduced to
\index{Divisor, greatest common|EtSeq}%
its lowest terms. When the numerator and the denominator
are both divisible by the same number,
their greatest common divisor can be found by a very
ingenious method which we owe to Euclid. This
\index{Euclid}%
method is exceedingly simple and lucid, but it may
be rendered even more palpable to the eye by the following
consideration. Suppose, for example, that you
have a given length, and that you wish to measure it.
The unit of measure is given, and you wish to know
how many times it is contained in the length. You
first lay off your measure as many times as you can on
\PageSep{3}
the given length, and that gives you a certain whole
number of measures. If there is no remainder your
operation is finished. But if there be a remainder,
\MNote{Greatest common divisor.}
that remainder is still to be evaluated. If the measure
is divided into equal parts, for example, into ten,
twelve, or more equal parts, the natural procedure is
to use one of these parts as a new measure and to see
how many times it is contained in the remainder.
You will then have for the value of your remainder,
a fraction of which the numerator is the number of
parts contained in the remainder and the denominator
the total number of parts into which the given measure
is divided.

I will suppose, now, that your measure is not so
divided but that you still wish to determine the ratio
of the proposed length to the length which you have
adopted as your measure. The following is the procedure
which most naturally suggests itself.

If you have a remainder, since that is less than the
\index{Fractions!continued|EtSeq}%
measure, naturally you will seek to find how many
times your remainder is contained in this measure.
Let us say two times, and that a remainder is still
left. Lay this remainder on the preceding remainder.
Since it is necessarily smaller, it will still be contained
a certain number of times in the preceding remainder,
say three times, and there will be another remainder
or there will not; and so on. In these different remainders
you will have what is called a \emph{continued fraction}.
For example, you have found that the measure
\PageSep{4}
is contained three times in the proposed length. You
have, to start with, the number \emph{three}. Then you have
\MNote{Continued fractions.}
found that your first remainder is contained twice in
your measure. You will have the fraction \emph{one} divided
by \emph{two}. But this last denominator is not complete,
for it was supposed there was still a remainder. That
remainder will give another and similar fraction, which
is to be added to the last denominator, and which by
our supposition is \emph{one} divided by \emph{three}. And so with
the rest. You will then have the fraction
\[
3 + \cfrac{1}{2 + \cfrac{1}{3 + \ddots}}
\]
as the expression of your ratio between the proposed
length and the adopted measure.

Fractions of this form are called \emph{continued fractions},
and can be reduced to ordinary fractions by the common
rules. Thus, if we stop at the first fraction, i.e.,
if we consider only the first remainder and neglect the
second, we shall have $3 + \frac{1}{2}$, which is equal to~$\frac{7}{2}$. Considering
only the first and the second remainders, we
stop at the second fraction, and shall have $3 + \dfrac{1}{2 + \frac{1}{3}}$.
Now $2 + \frac{1}{3} = \frac{7}{3}$. We shall have therefore $3 + \frac{3}{7}$, which
is equal to~$\frac{24}{7}$. And so on with the rest. If we arrive
in the course of the operation at a remainder which is
contained exactly in the preceding remainder, the
operation is terminated, and we shall have in the continued
\PageSep{5}
fraction a common fraction that is the exact
value of the length to be measured, in terms of the
length which served as our measure. If the operation
\MNote{Terminating continued fractions.}
is not thus terminated, it can be continued to infinity,
and we shall have only fractions which approach more
and more nearly to the true value.

If we now compare this procedure with that employed
for finding the greatest common divisor of two
numbers, we shall see that it is virtually the same
thing; the difference being that in finding the greatest
common divisor we devote our attention solely to
the different remainders, of which the last is the divisor
sought, whereas by employing the successive
quotients, as we have done above, we obtain fractions
which constantly approach nearer and nearer to the
fraction formed by the two numbers given, and of
which the last is that fraction itself reduced to its
lowest terms.

As the theory of continued fractions is little known,
but is yet of great utility in the solution of important
numerical questions, I shall enter here somewhat
more fully into the formation and properties of these
fractions. And, first, let us suppose that the quotients
found, whether by the mechanical operation, or by
the method for finding the greatest common divisor,
are, as above, $3$,~$2$, $3$, $5$, $7$,~$3$. The following is a rule
by which we can write down at once the convergent
fractions which result from these quotients, without
developing the continued fraction.
\PageSep{6}

The first quotient, supposed divided by unity,
will give the first fraction, which will be too small,
\MNote{Converging fractions.}
\index{Fractions!converging}%
namely,~$\frac{3}{1}$. Then, multiplying the numerator and denominator
of this fraction by the second quotient and
adding unity to the numerator, we shall have the second
fraction,~$\frac{7}{2}$, which will be too large. Multiplying
in like manner the numerator and denominator of this
fraction by the third quotient, and adding to the numerator
the numerator of the preceding fraction, and
to the denominator the denominator of the preceding
fraction, we shall have the third fraction, which will
be too small. Thus, the third quotient being~$3$, we
have for our numerator $(7 × 3 = 21) + 3 = 24$, and for
our denominator $(2 × 3 = 6) + 1 = 7$. The third convergent,
therefore, is~$\frac{24}{7}$. We proceed in the same
manner for the fourth convergent. The fourth quotient
being~$5$, we say $24$~times~$5$ is~$120$, and this plus~$7$,
the numerator of the fraction preceding, is~$127$;
similarly, $7$~times~$5$ is~$35$, and this plus~$2$ is~$37$. The
new fraction, therefore, is~$\frac{127}{37}$. And so with the rest.

In this manner, by employing the six quotients $3$,~$2$,
$3$, $5$, $7$,~$3$ we obtain the six fractions
\[
\frac{3}{1},\quad
\frac{7}{2},\quad
\frac{24}{7},\quad
\frac{127}{37},\quad
\frac{913}{266},\quad
\frac{2866}{835},
\]
of which the last, supposing the operation to be completed
at the sixth quotient~$3$, will be the required
value of the length measured, or the fraction itself
reduced to its lowest terms.

The fractions which precede the last are alternately
\PageSep{7}
smaller and larger than the last, and have the advantage
of approaching more and more nearly to its value
in such wise that no other fraction can approach it
\MNote{Convergents.}
\index{Convergents}%
more nearly except its denominator be larger than the
product of the denominator of the fraction in question
and the denominator of the fraction following. For
example, the fraction~$\frac{24}{7}$ is less than the true value
which is that of the fraction~$\frac{2866}{835}$, but it approaches
to it more nearly than any other fraction does whose
denominator is not greater than the product of~$7$ by~$37$,
that is,~$259$. Thus, any fraction expressed in large
numbers may be reduced to a series of fractions expressed
in smaller numbers and which approach as
near to it as possible in value.

The demonstration of the foregoing properties is
deduced from the nature of continued fractions, and
from the fact that if we seek the difference between
one of the convergent fractions and that next adjacent
to it we shall obtain a fraction of which the numerator
is always unity and the denominator the product of
the two denominators; a consequence which follows
\textit{\Typo{a}{à}~priori} from the very law of formation of these fractions.
Thus the difference between $\frac{7}{2}$~and~$\frac{3}{1}$ is~$\frac{1}{2}$, in
excess; between $\frac{24}{7}$~and~$\frac{7}{2}$, $\frac{1}{14}$,~in defect; between $\frac{127}{37}$
and~$\frac{24}{7}$, $\frac{1}{259}$,~in excess; and so on. The result being,
that by employing this series of differences we can
express in another and very simple manner the fractions
with which we are here concerned, by means of
a second series of fractions of which the numerators
\PageSep{8}
are all unity and the denominators successively the
products of every two adjacent denominators. Instead
\MNote{A second method of expression.}
of the fractions written above, we have thus the
series:
\[
\frac{3}{1} + \frac{1}{1 × 2}
  - \frac{1}{2 × 7}
  + \frac{1}{7 × 37}
  - \frac{1}{37 × 266}
  + \frac{1}{266 × 835}.
\]

The first term, as we see, is the first fraction, the
first and second together give the second fraction~$\frac{7}{2}$,
the first, the second, and the third give the third fraction~$\frac{24}{7}$,
and so on with the rest; the result being that
the series entire is equivalent to the last fraction.

There is still another way, less known but in some
respects more simple, of treating the same question---which
leads directly to a series similar to the preceding.
Reverting to the previous example, after having
found that the measure goes three times into the length
to be measured and that after the first remainder has
been applied to the measure there is left a new remainder,
instead of comparing this second remainder
with the preceding, as we did above, we may compare
it with the measure itself. Thus, supposing it goes
into the latter seven times with a remainder, we again
compare this last remainder with the measure, and so
on, until we arrive, if possible, at a remainder which
is an aliquot part of the measure,---which will terminate
the operation. In the contrary event, if the
measure and the length to be measured are incommensurable,
the process may be continued to infinity.
\PageSep{9}
We shall have then, as the expression of the length
measured, the series
\MNote{A third method of expression.}
\[
3 + \frac{1}{2} - \frac{1}{2 × 7} + \ldots.
\]

It is clear that this method is also applicable to
ordinary fractions. We constantly retain the denominator
of the fraction as the dividend, and take the different
remainders successively as divisors. Thus, the
fraction~$\frac{2866}{835}$ gives the quotients $3$,~$2$, $7$, $18$, $19$, $46$,
$119$, $417$\Typo{}{,}~$835$; from which we obtain the series
\[
3 + \frac{1}{2} - \frac{1}{2 × 7}
  + \frac{1}{2 × 7 × 18}
  - \frac{1}{2 × 7 × 18 × 19} + \ldots;
\]
and as these partial fractions rapidly diminish, we
shall have, by combining them successively, the simple
fractions,
\[
\frac{7}{2},\quad
\frac{48}{2 × 7},\quad
\frac{865}{2 × 7 × 18}, \ldots,
\]
which will constantly approach nearer and nearer to
the true value sought, and the error will be less than
the first of the partial fractions neglected.

Our remarks on the foregoing methods of evaluating
fractions should not be construed as signifying
that the employment of decimal fractions is not nearly
\index{Decimal!fractions}%
\index{Fractions!decimal}%
always preferable for expressing the values of fractions
to whatever degree of exactness we wish. But cases
occur where it is necessary that these values should
be expressed by as few figures as possible. For example,
if it were required to construct a planetarium,
\index{Planetarium}%
\PageSep{10}
since the ratios of the revolutions of the planets to one
another are expressed by very large numbers, it would
\MNote{Origin of continued fractions.}
\index{Fractions!origin of continued}%
be necessary, in order not to multiply unduly the
number of the teeth on the wheels, to avail ourselves
of smaller numbers, but at the same time so to select
them that their ratios should approach as nearly as
possible to the actual ratios. It was, in fact, this very
question that prompted Huygens, in his search for its
\index{Huygens}%
solution, to resort to continued fractions and that so
gave birth to the theory of these fractions. Afterwards,
in the elaboration of this theory, it was found
adapted to the solution of other important questions,
and this is the reason, since it is not found in elementary
works, that I have deemed it necessary to go
somewhat into detail in expounding its principles.

We will now pass to the theory of powers, proportions,
and progressions.

As you already know, a number multiplied by itself
\index{Powers|EtSeq}%
gives its square, and multiplied again by itself
gives its cube, and so on. In geometry we do not go
beyond the cube, because no body can have more than
three dimensions. But in algebra and arithmetic we
may go as far as we please. And here the theory of
the extraction of roots takes its origin. For, although
every number can be raised to its square and to its
cube and so forth, it is not true reciprocally that every
number is an exact square or an exact cube. The
number~$2$, for example, is not a square; for the square
of~$1$ is~$1$, and the square of~$2$ is four; and there being
\PageSep{11}
no other whole numbers between these two, it is impossible
to find a whole number which multiplied by
itself will give~$2$. It cannot be found in fractions, for
\MNote{Involution and evolution.}
\index{Evolution}%
\index{Involution and evolution}%
if you take a fraction reduced to its lowest terms, the
square of that fraction will again be a fraction reduced
to its lowest terms, and consequently cannot be equal
to the whole number~$2$. But though we cannot obtain
the square root of~$2$ exactly, we can yet approach to it
as nearly as we please, particularly by decimal fractions.
By following the common rules for the extraction
of square roots, cube roots, and so forth, the process
may be extended to infinity, and the true values
of the roots may be approximated to any degree of
exactitude we wish.

But I shall not enter into details here. The theory
of powers has given rise to that of progressions, before
entering on which a word is necessary on proportions.

Every fraction expresses a ratio. Having two equal
\index{Proportion|EtSeq}%
\index{Ratios, constant|EtSeq}%
fractions, therefore, we have two equal ratios; and
the numbers constituting the fractions or the ratios
form what is called a \emph{proportion}. Thus the equality
of the ratios $2$~to~$4$ and $3$~to~$6$ gives the proportion
$2 : 4 :: 3 : 6$, because $4$~is the double of~$2$ as $6$~is the
double of~$3$. Many of the rules of arithmetic depend
on the theory of proportions. First, it is the foundation
of the famous \emph{rule of three}, which is so extensively
\index{Rule!three@of three|EtSeq}%
used. You know that when the first three terms of a
proportion are given, to obtain the fourth you have
\PageSep{12}
only to multiply the last two together and divide the
product by the first. Various special rules have also
\MNote{Proportions\Add{.}}
been conceived and have found a place in the books
on arithmetic; but they are all reducible to the rule
of three and may be neglected if we once thoroughly
grasp the conditions of the problem. There are direct,
inverse, simple, and compound rules of three, rules of
partnership, of mixtures, and so forth. In all cases
it is only necessary to consider carefully the conditions
of the problem and to arrange the terms of the
proportion correspondingly.

I shall not enter into further details here. There
\index{Progressions, theory of}%
is, however, another theory which is useful on numerous
occasions,---namely, the \emph{theory of progressions}.
When you have several numbers that bear the same
proportion to one another, and which follow one another
in such a manner that the second is to the first
as the third is to the second, as the fourth is to the
third, and so forth, these numbers form a progression.
I shall begin with an observation.

The books of arithmetic and algebra ordinarily distinguish
between two kinds of progression, arithmetical
and geometrical, corresponding to the proportions
called arithmetical and geometrical. But the appellation
proportion appears to me extremely inappropriate
as applied to \emph{arithmetical proportion}. And as it
\index{Arithmetical proportion}%
is one of the objects of the \textit{École Normale} to rectify
\index{Ecole@{\Typo{Ecole}{École} Normale}}%
the language of science, the present slight digression
will not be considered irrelevant.
\PageSep{13}

I take it, then, that the idea of proportion is already
well established by usage and that it corresponds solely
to what is called \emph{geometrical proportion}. When we
\index{Geometrical!proportion}%
\MNote{Arithmetical and geometrical proportions.}
speak of the proportion of the parts of a man's body,
of the proportion of the parts of an edifice,~etc.; when
we say that a plan should be reduced proportionately
in size,~etc.; in fact, when we say generally that one
thing is proportional to another, we understand by
proportion equality of ratios only, as in geometrical
proportion, and never equality of differences as in
arithmetical proportion. Therefore, instead of saying
\index{Equi-different numbers}%
that the numbers, $3$,~$5$, $7$,~$9$, are in arithmetical proportion,
because the difference between $5$~and~$3$ is the
same as that between $9$~and~$7$, I deem it desirable that
some other term should be employed, so as to avoid
all ambiguity. We might, for instance, call such numbers
\emph{equi-different}, reserving the name of \emph{proportionals}
for numbers that are in geometrical proportion, as $2$,~$4$,
$6$,~$8$,~etc.

As for the rest, I cannot see why the proportion
called \emph{arithmetical} is any more arithmetical than that
which is called geometrical, nor why the latter is more
geometrical than the former. On the contrary, the
primitive idea of geometrical proportion is based on
arithmetic, for the notion of ratios springs essentially
from the consideration of numbers.

Still, in waiting for these inappropriate designations
to be changed, I shall continue to make use of
them, as a matter of simplicity and convenience.
\PageSep{14}

The theory of arithmetical progressions presents
few difficulties. Arithmetical progressions consist of
\MNote{Progressions.}
\index{Progressions, theory of}%
quantities which increase or diminish constantly by
the same amount. But the theory of geometrical progressions
is more difficult and more important, as a
large number of interesting questions depend upon it---for
example, all problems of compound interest, all
problems that relate to discount, and many others of
like nature.

In general, quantities in geometrical proportion
are produced, when a quantity increases and the force
generating the increase, so to speak, is proportional
to that quantity. It has been observed that in countries
where the means of subsistence are easy of acquisition,
as in the first American colonies, the population
is doubled at the expiration of twenty years; if
it is doubled at the end of twenty years it will be quadrupled
at the end of forty, octupled at the end of sixty,
and so on; the result being, as we see, a geometrical
progression, corresponding to intervals of time in
arithmetical progression. It is the same with compound
interest. If a given sum of money produces,
at the expiration of a certain time, a certain sum, at
the end of double that time, the original sum will have
produced an equivalent additional sum, and in addition
the sum produced in the first space of time will,
in its proportion, likewise have produced during the
second space of time a certain sum; and so with the
rest. The original sum is commonly called the \emph{principal},
\PageSep{15}
the sum produced the \emph{interest}, and the constant
\index{Interest}%
ratio of the principal to the interest per annum, the
\emph{rate}. Thus, the rate \emph{twenty} signifies that the interest
\MNote{Compound interest.}
is the twentieth part of the principal,---a rate which
is commonly called $5$~\emph{per cent.}, $5$~being the twentieth
part of~$100$. On this basis, the principal, at the end
of one year, will have increased by its one-twentieth
part; consequently, it will have been augmented in
the ratio of $21$~to~$20$. At the end of two years, it will
have been increased again in the same ratio, that is in
the ratio of $\frac{21}{20}$~multiplied by~$\frac{21}{20}$; at the end of three
years, in the ratio of $\frac{21}{20}$~multiplied twice by itself; and
so on. In this manner we shall find that at the end of
fifteen years it will almost have doubled itself, and that
at the end of fifty-three years it will have increased
tenfold. Conversely, then, since a sum paid now will
be doubled at the end of fifteen years, it is clear that
a sum not payable till after the expiration of fifteen
years is now worth only one-half its amount. This
is what is termed the \emph{present value} of a sum payable
\index{Present value}%
at the end of a certain time; and it is plain, that to
find that value, it is only necessary to divide the sum
promised by the fraction~$\frac{21}{20}$, or to multiply it by the
fraction~$\frac{20}{21}$, as many times as there are years for the
sum to run. In this way we shall find that a sum
payable at the end of fifty-three years, is worth at
present only one-tenth. From this it is evident what
little advantage is to be derived from surrendering the
absolute ownership of a sum of money in order to obtain
\PageSep{16}
the enjoyment of it for a period of only fifty
years, say; seeing that we gain by such a transaction
\MNote{Present values and annuities.}
\index{Annuities}%
only one-tenth in actual use, whilst we lose the ownership
of the property forever.

In \emph{annuities}, the consideration of interest is combined
with that of the probability of life; and as
every one is prone to believe that he will live very
long, and as, on the other hand, one is apt to under-*estimate
the value of property which must be abandoned
on death, a peculiar temptation arises, when
one is without children, to invest one's fortune, wholly
or in part, in annuities. Nevertheless, when put to
the test of rigorous calculation, annuities are not
found to offer sufficient advantages to induce people
to sacrifice for them the ownership of the original
capital. Accordingly, whenever it has been attempted
to create annuities sufficiently attractive to induce individuals
to invest in them, it has been necessary to
offer them on terms which are onerous to the company.

But we shall have more to say on this subject when
we expound the theory of annuities, which is a branch
of the calculus of probabilities.

I shall conclude the present lecture with a word
\index{Logarithms|EtSeq}%
on \emph{logarithms}. The simplest idea which we can form
of the theory of logarithms, as they are found in the
ordinary tables, is that of conceiving all numbers
as powers of~$10$; the exponents of these powers,
then, will be the logarithms of the numbers. From
\PageSep{17}
this it is evident that the multiplication and division
of two numbers is reducible to the addition and subtraction
of their respective exponents, that is, of their
\MNote{Logarithms\Add{.}}
logarithms. And, consequently, involution and the
extraction of roots are reducible to multiplication and
division, which is of immense advantage in arithmetic
and renders logarithms of priceless value in that science.

But in the period when logarithms were invented,
mathematicians were not in possession of the theory
of powers. They did not know that the root of a number
could be represented by a fractional power. The
following was the way in which they approached the
problem.

The primitive idea was that of two corresponding
progressions, one arithmetical, and the other geometrical.
In this way the general notion of a logarithm
was reached. But the means for finding the logarithms
of all numbers were still lacking. As the numbers
follow one another in arithmetical progression, it
was requisite, in order that they might all be found
among the terms of a geometrical progression, so to
establish that progression that its successive terms
should differ by extremely small quantities from one
another; and, to prove the possibility of expressing
all numbers in this way, Napier, the inventor, first
\index{Napier|EtSeq}%
considered them as expressed by lines and parts of
lines, and these lines he considered as generated by
\PageSep{18}
the continuous motion of a point, which was quite
natural.

\MNote{Napier (1550--1617).}
He considered, accordingly, two lines, the first of
which was generated by the motion of a point describing
in equal times spaces in geometrical progression,
and the other generated by a point which described
spaces that increased as the times and consequently
formed an arithmetical progression corresponding to
the geometrical progression. And he supposed, for
the sake of simplicity, that the initial velocities of
these two points were equal. This gave him the logarithms,
at first called \emph{natural}, and afterwards \emph{hyperbolical},
when it was discovered that they could be expressed
as parts of the area included between a
hyperbola and its asymptotes. By this method it is
clear that to find the logarithm of any given number,
it is only necessary to take a part on the first line
equal to the given number, and to seek the part on
the second line which shall have been described in
the same interval of time as the part on the first.

Conformably to this idea, if we take as the two
first terms of our geometrical progression the numbers
with very small differences $1$~and~$1.0000001$, and as
those of our arithmetical progression $0$~and $0.0000001$,
and if we seek successively, by the known rules, all
the following terms of the two progressions, we shall
find that the number~$2$ expressed approximately to the
eighth place of decimals is the $6931472$th~term of the
geometrical progression, that is, that the logarithm of~$2$
\PageSep{19}
is~$0.6931472$. The number~$10$ will be found to be the
$23025851$th~term of the same progression; therefore,
the logarithm of~$10$ is~$2.3025851$, and so with the rest.
\MNote{Origin of logarithms\Add{.}}
\index{Logarithms!origin of}%
But Napier, having to determine only the logarithms
of numbers less than unity for the purposes of trigonometry,
where the sines and cosines of angles are
expressed as fractions of the radius, considered a decreasing
geometrical progression of which the first
two terms were $1$~and~$0.9999999$; and of this progression
he determined the succeeding terms by enormous
computations. On this last hypothesis, the logarithm
which we have just found for~$2$ becomes that of the
number~$\frac{1}{5}$ or~$0.5$, and that of the number~$10$ becomes
that of the number~$\frac{1}{10}$ or~$0.1$; as is readily apparent
from the nature of the two progressions.

Napier's work appeared in~1614. Its utility was
felt at once. But it was also immediately seen that it
would conform better to the decimal system of our
arithmetic, and would be simpler, if the logarithm of~$10$
were made unity, conformably to which that of~$100$
would be~$2$, and so with the rest. To that end, instead
of taking as the first two terms of our geometrical
progression the numbers $1$~and~$\Typo{0.0000001}{1.0000001}$, we should
have to take the numbers $1$~and~$1.0000002302$, retaining
$0$~and~$0.0000001$ as the corresponding terms of the
arithmetical progression. Whence it will be seen,
that, while the point which is supposed to generate by
its motion the geometrical line, or the numbers, is
describing the very small portion~$0.0000002302\dots$,
\PageSep{20}
the other point, the office of which is to generate
simultaneously the arithmetical line, will have described
\MNote{Briggs (1556--1631). Vlacq.}
\index{Briggs}%
\index{Vlacq}%
the portion~$0.0000001$; and that therefore the
spaces described in the same time by the two points
at the beginning of their motion, that is to say, their
initial velocities, instead of being equal, as in the
preceding system, will be in the proportion of the
numbers $2.302\dots$~to~$1$, where it will be remarked
that the number~$2.302\dots$ is exactly the number
which in the original system of natural logarithms
stood for the logarithm of~$10$,---a result demonstrable
\textit{à~priori}, as we shall see when we come to apply
the formulæ of algebra to the theory of logarithms.
Briggs, a contemporary of Napier, is the author of this
change in the system of logarithms, as he is also of
the tables of logarithms now in common use. A portion
\index{Logarithms!tables of}%
of these was calculated by Briggs himself, and
the remainder by Vlacq, a Dutchman.

These tables appeared at Gouda, in~1628. They
contain the logarithms of all numbers from~$1$ to~$100000$
to ten decimal places, and are now extremely rare.
But it was afterwards discovered that for ordinary purposes
seven decimals were sufficient, and the logarithms
are found in this form in the tables which are
used to-day. Briggs and Vlacq employed a number
of highly ingenious artifices for facilitating their work.
The device which offered itself most naturally and
which is still one of the simplest, consists in taking
the numbers $1$,~$10$, $100$,~$\dots$, of which the logarithms
\PageSep{21}
are $0$,~$1$,~$2$,~$\dots$, and in interpolating between the successive
terms of these two series as many corresponding
terms as we desire, in the first series by geometrical
\MNote{Computation of logarithms.}
mean proportionals and in the second by
arithmetical means. In this manner, when we have
arrived at a term of the first series approaching, to the
eighth decimal place, the number whose logarithm
we seek, the corresponding term of the other series
will be, to the eighth decimal place approximately,
the logarithm of that number. Thus, to obtain the
logarithm of~$2$, since $2$~lies between $1$~and~$10$, we seek
first by the extraction of the square root of~$10$, the
geometrical mean between $1$~and~$10$, which we find to
be~$3.16227766$, while the corresponding arithmetical
mean between $0$~and~$1$ is~$\frac{1}{2}$ or~$0.50000000$; we are
assured thus that this last number is the logarithm of
the first. Again, as $2$~lies between $1$~and~$3.16227766$,
the number just found, we seek in the same manner
the geometrical mean between these two numbers,
and find the number~$1.77827941$. As before, taking
the arithmetical mean between $0$~and~$5.0000000$, we
shall have for the logarithm of~$1.77827941$ the number~$0.25000000$.
Again, $2$~lying between $1.77827941$
and~$3.16227766$, it will be necessary, for still further
approximation, to find the geometrical mean between
these two, and likewise the arithmetical mean between
their logarithms. And so on. In this manner,
by a large number of similar operations, we find that
the logarithm of~$2$ is~$0.3010300$, that of~$3$ is~$0.4771213$,
\PageSep{22}
and so on, not carrying the degree of exactness beyond
the seventh decimal place. But the preceding
\MNote{Value of the history of science.}
\index{Science!history of}%
calculation is necessary only for prime numbers; because
the logarithms of numbers which are the product
of two or several others, are found by simply
taking the sum of the logarithms of their factors.

As for the rest, since the calculation of logarithms
is now a thing of the past, except in isolated instances,
it may be thought that the details into which we have
here entered are devoid of value. We may, however,
justly be curious to know the trying and tortuous
paths which the great inventors have trodden, the different
\index{Inventors, great}%
steps which they have taken to attain their goal,
and the extent to which we are indebted to these veritable
benefactors of the human race. Such knowledge,
moreover, is not matter of idle curiosity. It can
afford us guidance in similar inquiries and sheds an
increased light on the subjects with which we are
employed.

Logarithms are an instrument universally employed
in the sciences, and in the arts depending on calculation.
The following, for example, is a very evident
application of their use.

Persons not entirely unacquainted with music know
\index{Music}%
that the different notes of the octave are expressed by
numbers which give the divisions of a stretched cord
producing those notes. Thus, the principal note being
denoted by~$1$, its octave will be denoted by~$\frac{1}{2}$,
its fifth by~$\frac{2}{3}$, its third by~$\frac{4}{5}$, its fourth by~$\frac{3}{4}$, its second
\PageSep{23}
by~$\frac{8}{9}$, and so on. The distance of one of these notes
from that next adjacent to it is called an \emph{interval}, and
is measured, not by the difference, but by the ratio of
the numbers expressing the two sounds. Thus, the
interval between the fourth and fifth, which is called
the \emph{major tone}, is regarded as sensibly double of that
between the third and fourth, which is called the \emph{semi-major}.
In fact, the first being expressed by~$\frac{8}{9}$, the
second by~$\frac{15}{16}$, it can be easily proved that the first
does not differ by much from the square of the second.
Now, it is clear that this conception of intervals, on
\MNote{Musical temperament.}
\index{Temperament, theory of}%
which the whole theory of temperament is founded,
conducts us naturally to logarithms. For if we express
the value of the different notes by the logarithms
of the lengths of the cords answering to them,
then the interval of one note from another will be
expressed by the simple difference of values of the
two notes; and if it were required to divide the octave
into twelve equal semi-tones, which would give the
temperament that is simplest and most exact, we
should simply have to divide the logarithm of one
half, the value of the octave, into twelve equal parts.
\PageSep{24}


\Lecture{II.}{On the Operations of Arithmetic.}
\index{Arithmetic!operations of|EtSeq}%

\First{An ancient} writer once remarked that arithmetic
and geometry were \emph{the wings of mathematics}.
\index{Geometry}%
\index{Mathematics!wings of}%
\MNote{Arithmetic and geometry.}
I believe we can say, without metaphor, that
these two sciences are the foundation and essence of
all the sciences that treat of magnitude. But not
only are they the foundation, they are also, so to
speak, the capstone of these sciences. For, whenever
we have reached a result, in order to make use of it,
it is requisite that it be translated into numbers or
into lines; to translate it into numbers, arithmetic is
necessary; to translate it into lines, we must have
recourse to geometry.

The importance of arithmetic, accordingly, leads
me to the further discussion of that subject to-day,
although we have begun algebra. I shall take up its
several parts, and shall offer new observations, which
will serve to supplement what I have already expounded
to you. I shall employ, moreover, the geometrical
\index{Geometrical!calculus}%
calculus, wherever that is necessary for giving
\PageSep{25}
greater generality to the demonstrations and
methods.

First, then, as regards addition, there is nothing
to be added to what has already been said. Addition
is an operation so simple in character that its conception
is a matter of course. But with regard to subtraction,
\MNote{New method of subtraction\Add{.}}
\index{Subtraction, new method of|EtSeq}%
there is another manner of performing that
operation which is frequently more advantageous than
the common method, particularly for those familiar
with it. It consists in converting the subtraction into
addition by taking the complement of every figure of
the number which is to be subtracted, first with respect
to~$10$ and afterwards with respect to~$9$. Suppose,
for example, that the number~$2635$ is to be subtracted
from the number~$7853$. Instead of saying $5$~from~$13$
\begin{figure}[hbt!]
\centering
$\begin{array}{r}
7853 \\
2635 \\
\hline
5218
\end{array}$
\end{figure}
leaves~$8$; $3$~from~$4$ leaves~$1$; $6$~from~$8$ leaves~$2$;
and $2$~from~$7$ leaves~$5$, giving a total remainder of~$5218$,---I
say: $5$~the complement of~$5$ with respect to~$10$
added to~$3$ gives~$8$,---I write down~$8$; $6$~the complement
of~$3$ with respect to~$9$ added to~$5$ gives~$11$,---I
write down~$1$ and carry~$1$; $3$~the complement of~$6$
with respect to~$9$, plus~$9$, by reason of the $1$~carried,
gives~$12$,---I put down~$2$ and carry~$1$; lastly, $7$~the
complement of~$2$ with respect to~$9$ plus~$8$, on account
of the $1$~carried, gives~$15$,---I put down~$5$ and this time
carry nothing, for the operation is completed, and the
\PageSep{26}
last~$10$ which was borrowed in the course of the operation
must be rejected. In this manner we obtain the
same remainder as above,~$5218$.

The foregoing method is extremely convenient
\MNote{Subtraction by complements.}
\index{Complements, subtraction by}%
when the numbers are large; for in the common
method of subtraction, where borrowing is necessary
in subtracting single numbers from one another, mistakes
are frequently made, whereas in the method
with which we are here concerned we never borrow
but simply carry, the subtraction being converted into
addition. With regard to the complements they are
discoverable at the merest glance, for every one knows
that $3$~is the complement of~$7$ with respect to~$10$, $4$~the
complement of~$5$ with respect to~$9$,~etc. And as
to the reason of the method, it too is quite palpable.
The different complements taken together form the
total complement of the number to be subtracted
either with respect to~$10$ or~$100$ or~$1000$, etc., according
as the number has $1$,~$2$,~$3$~$\dots$ figures; so that the
operation performed is virtually equivalent to first
adding $10$,~$100$, $1000$~$\dots$ to the minuend and then
taking the subtrahend from the minuend as so augmented.
Whence it is likewise apparent why the~$10$
of the sum found by the last partial addition must be
rejected.

As to multiplication, there are various abridged
\index{Multiplication!abridged methods of|EtSeq}%
methods possible, based on the decimal system of
numbers. In multiplying by~$10$, for example, we have,
as we know, simply to add a cipher; in multiplying
\PageSep{27}
by~$100$ we add two ciphers; by~$1000$, three ciphers,~etc.
Consequently, to multiply by any aliquot part of~$10$,
for example~$5$, we have simply to multiply by~$10$
\MNote{Abridged multiplication.}
and then divide by~$2$; to multiply by~$25$ we multiply
by~$100$ and divide by~$4$, and so on for all the products
of~$5$.

When decimal numbers are to be multiplied by
\index{Decimal!numbers|EtSeq}%
decimal numbers, the general rule is to consider the
two numbers as integers and when the operation is
finished to mark off from the right to the left as many
places in the product as there are decimal places in
the multiplier and the multiplicand together. But in
practice this rule is frequently attended with the inconvenience
of unnecessarily lengthening the operation,
for when we have numbers containing decimals
these numbers are ordinarily exact only to a certain
number of places, so that it is necessary to retain in
the product only the decimal places of an equivalent
order. For example, if the multiplicand and the multiplier
each contain two places of decimals and are exact
only to two decimal places, we should have in the
product by the ordinary method four decimal places,
the two last of which we should have to reject as useless
and inexact. I shall give you now a method for
obtaining in the product only just so many decimal
places as you desire.

I observe first that in the ordinary method of multiplying
we begin with the units of the multiplier which
we multiply with the units of the multiplicand, and so
\PageSep{28}
continue from the right to the left. But there is nothing
compelling us to begin at the right of the multiplier.
\MNote{Inverted multiplication.}
\index{Multiplication!inverted}%
We may equally well begin at the left. And
I cannot in truth understand why the latter method
should not be preferred, since it possesses the advantage
of giving at once the figures having the greatest
value, and since, in the majority of cases where large
numbers are multiplied together, it is just these last
and highest places that concern us most; we frequently,
in fact, perform multiplications only to find
what these last figures are. And herein, be it parenthetically
remarked, consists one of the great advantages
in calculating by logarithms, which always
\index{Logarithms!advantages in calculating by}%
give, be it in multiplication or division, in involution
or evolution, the figures in the descending order of
their value, beginning with the highest and proceeding
from the left to the right.

By performing multiplication in this manner, no
difference is caused in the total product. The sole
distinction is, that by the new method the first line,
the first partial product, is that which in the ordinary
method is last, and the second partial product is that
which in the ordinary method is next to the last, and
so with the rest.

Where whole numbers are concerned and the exact
product is required, it is indifferent which method we
employ. But when decimal places are involved the
prime essential is to have the figures of the whole
numbers first in the product and to descend afterwards
\PageSep{29}
successively to the figures of the decimal parts,
instead of, as in the ordinary method, beginning with
the last decimal places and successively ascending to
the figures forming the whole numbers.

In applying this method practically, we write the
multiplier underneath the multiplicand so that the
units' figure of the multiplier falls beneath the last
\MNote{Approximate multiplication.}
\index{Multiplication!approximate}%
figure of the multiplicand. We then begin with the
last left-hand figure of the multiplier which we multiply
as in the ordinary method by all the figures of the
multiplicand, beginning with the last to the right and
proceeding successively to the left; observing that the
first figure of the product is to be placed underneath
the figure with which we are multiplying, while the
others follow in their successive order to the left. We
proceed in the same manner with the second figure of
the multiplier, likewise placing beneath this figure the
first figure of the product, and so on with the rest.
The place of the decimal point in these different products
will be the same as in the multiplicand, that is
to say, the units of the products will all fall in the
same vertical line with those of the multiplicand and
consequently those of the sum of all the products or
of the total product will also fall in that line. In this
manner it is an easy matter to calculate only as many
decimal places as we wish. I give below an example
of this method in which the multiplicand is~$437.25$
and the multiplier~$27.34$:
\PageSep{30}
\MNote{The new method exemplified.}
\[
\begin{array}{r@{\,}l}
437\PadTo[l]{\,}{.} & 25 \\
            & 27.34 \\
\hline
\MultRow{8745}{0} \\
\MultRow{3060}{75} \\
\MultRow{131}{17\phantom{.}5} \\
\MultRow{17}{49\phantom{.}00} \\
\hline
\MultRow{11954}{41\phantom{.}50}
\end{array}
\]

I have written all the decimals in the product, but
\index{Decimals!multiplication of}%
it is easy to see how we may omit calculating the decimals
which we wish to neglect. The vertical line is
used to mark more distinctly the place of the decimal
point.

The preceding rule appears to me simpler and
more natural than that which is attributed to Oughtred
\index{Oughtred}%
and which consists in writing the multiplier underneath
the multiplicand in the reverse order.

There is one more point, finally, to be remarked
in connexion with the multiplication of numbers containing
\index{Multiplication!decimals@of decimals}%
decimals, and that is that we may alter the
place of the decimal point of either number at will.
For seeing that moving the decimal point from the
right to the left in one of the numbers is equivalent to
dividing the number by~$10$, by~$100$, or by~$1000\dots$, and
that moving the decimal point back in the other number
the same number of places from the left to the
right is tantamount to multiplying that number by~$10$,
$100$, or~$1000$,~$\dots$, it follows that we may push the
decimal point forward in one of the numbers as many
places as we please provided we move it back in the
other number the same number of places, without in
\PageSep{31}
any wise altering the product. In this way we can
always so arrange it that one of the two numbers shall
contain no decimals---which simplifies the question.

Division is susceptible of a like simplification, for
\index{Decimals!division of}%
\index{Division!decimals@of decimals}%
since the quotient is not altered by multiplying or dividing
\MNote{Division of decimals.}
the dividend and the divisor by the same number,
it follows that in division we may move the decimal
point of both numbers forwards or backwards as
many places as we please, provided we move it the
same distance in each case. Consequently, we can
always reduce the divisor to a whole number---which
facilitates infinitely the operation for the reason that
when there are decimal places in the dividend only,
we may proceed with the division by the common
method and neglect all places giving decimals of a
lower order than those we desire to take account of.

You know the remarkable property of the number~$9$,
\index{Nine!property of the number|EtSeq}%
whereby if a number be divisible by~$9$ the sum of
its digits is also divisible by~$9$. This property enables
us to tell at once, not only whether a number is divisible
by~$9$ but also what is its remainder from such division.
For we have only to take the sum of its digits
and to divide that sum by~$9$, when the remainder will
be the same as that of the original number divided
by~$9$.

The demonstration of the foregoing proposition is
not difficult. It reposes upon the fact that the numbers
$10$~less~$1$, $100$~less~$1$, $1000$~less~$1$,~$\dots$ are all divisible
\PageSep{32}
by~$9$,---which seeing that the resulting numbers
are $9$,~$99$, $999$,~$\dots$ is quite obvious.

If, now, you subtract from a given number the
sum of all its digits, you will have as your remainder
\MNote{Property of the number~$9$.}
the tens' digit multiplied by~$9$, the hundreds' digit
multiplied by~$99$, the thousands' digit multiplied by~$999$,
and so on,---a remainder which is plainly divisible
by~$9$. Consequently, if the sum of the digits is
divisible by~$9$, the original number itself will be so
divisible, and if it is not divisible by~$9$ the original
number likewise will not be divisible thereby. But
the remainder in the one case will be the same as in
the other.

In the case of the number~$9$, it is evident immediately
that $10$~less~$1$, $100$~less~$1$,~$\dots$ are divisible by~$9$;
but algebra demonstrates that the property in
question holds good for every number~$a$. For it can
be shown that
\[
a - 1,\quad a^{2} - 1,\quad a^{3} - 1,\quad a^{4} - 1, \dots
\]
are all quantities divisible by~$a - 1$, actual division
giving the quotients
\[
1,\quad a + 1,\quad a^{2} + a + 1,\quad a^{3} + a^{2} + a + 1, \dots.
\]

The conclusion is therefore obvious that the aforesaid
property of the number~$9$ holds good in our decimal
system of arithmetic because $9$~is $10$~less~$1$, and
that in any other system founded upon the progression
$a$,~$a^{2}$,~$a^{3}$,~$\dots$ the number~$a - 1$ would enjoy the
same property. Thus in the duodecimal system it
\index{Duodecimal system}%
\PageSep{33}
would be the number~$11$; and in this system every
number, the sum of whose digits was divisible by~$11$,
would also itself be divisible by that number.

The foregoing property of the number~$9$, now, admits
\index{Nine!property of the number generalised}%
of generalisation, as the following consideration
\MNote{Property of the number~$9$ generalised.}
will show. Since every number in our system is represented
by the sum of certain terms of the progression
$1$,~$10$, $100$, $1000$,~$\dots$, each multiplied by one of
the nine digits $1$,~$2$, $3$, $4$,~$\dots$\Add{,}~$9$, it is easy to see that
the remainder resulting from the division of any number
by a given divisor will be equal to the sum of the
remainders resulting from the division of the terms $1$,
$10$, $100$, $1000$,~$\dots$ by that divisor, each multiplied by
the digit showing how many times the corresponding
term has been taken. Hence, generally, if the given
divisor be called~$D$, and if $m$,~$n$,~$p$,~$\dots$ be the remainders
of the division of the numbers $1$, $10$, $100$, $1000$
by~$D$, the remainder from the division of any number
whatever~$N$, of which the characters proceeding from
the right to the left are $a$,~$b$,~$c$,~$\dots$, by~$D$ will obviously
be equal to
\[
ma + nb + pc + \dots.
\]
Accordingly, if for a given divisor~$D$ we know the remainders
$m$,~$n$,~$p$,~$\dots$, which depend solely upon that
divisor and which are always the same for the same
divisor, we have only to write the remainders underneath
the original number, proceeding from the right
to the left, and then to find the different products of
\PageSep{34}
each digit of the number by the digit which is underneath
it. The sum of all these products will be the
\MNote{Theory of remainders\Add{.}}
\index{Remainders!theory of|EtSeq}%
total remainder resulting from the division of the proposed
number by the same divisor~$D$. And if the sum
found is greater than~$D$, we can proceed in the same
manner to seek its remainder from division by~$D$, and
so on until we arrive finally at a remainder which is
less than~$D$, which will be the true remainder sought.
It follows from this that the proposed number cannot
be exactly divisible by the given divisor unless the
last remainder found by this method is zero.

The remainders resulting from the division of the
terms $1$, $10$, $100$,~$\dots$\Add{,} $1000$, by~$9$ are always unity.
\index{Division!nine@by \textit{nine}}%
Hence, the sum of the digits of any number whatever
is the remainder resulting from the division of that
number by~$9$. The remainders resulting from the division
of the same terms by~$8$ are $1$,~$2$, $4$, $0$, $0$, $0$,~$\dots$.
\index{Division!eight@by \textit{eight}}%
We shall obtain, accordingly, the remainder resulting
from dividing any number by~$8$, by taking the sum
of the first digit to the right, the second digit next
thereto to the left multiplied by~$2$, and the third digit
multiplied by~$4$.

The remainders resulting from the divisions of the
\index{Division!seven@by \textit{seven}|EtSeq}%
terms $1$, $10$, $100$, $1000$,~$\dots$ by~$7$ are $1$, $3$, $2$, $6$, $4$, $5$,
$1$, $3$,~$\dots$, where the same remainders continually recur
in the same order. If I have, now, the number
$13527541$ to be divided by~$7$, I write it thus with the
above remainders underneath it:
\PageSep{35}
\index{Seven, tests of divisibility by}%
\MNote{Test of divisibility by~$7$.}
\[
\begin{array}{@{\,}*{2}{r@{}}r@{\,}}
13527&5&41 \\
31546&2&31 \\
\hline
&&  1 \\
&& 12 \\
&& 10 \\
&& 42 \\
&&  8 \\
&& 25 \\
&&  3 \\
&&  3 \\
\cline{2-3}
& 1&04 \\
& 2&31 \\
\cline{2-3}
&&  4 \\
&&  0 \\
&&  2 \\
\cline{3-3}
&&  6
\end{array}
\]

Taking the partial products and adding them, I
obtain~$104$, which would be the remainder from the
division of the given number by~$7$, were it not greater
than the divisor. I accordingly repeat the operation
with this remainder, and find for my second remainder~$6$,
which is the real remainder in question.

I have still to remark with regard to the preceding
remainders and the multiplications which result from
them, that they may be simplified by introducing negative
remainders in the place of remainders which are
greater than half the divisor, and to accomplish this
we have simply to subtract the divisor from each of
such remainders. We obtain thus, instead of the remainders
$6$,~$5$,~$4$, the following:
\[
-1,\quad -2,\quad -3.
\]
\PageSep{36}
The remainders for the divisor~$7$, accordingly, are
\[
1,\quad 3,\quad 2,\quad -1,\quad -3,\quad -2,\quad 1,\quad 3, \dots
\]
and so on to infinity.

\MNote{Negative remainders\Add{.}}
\index{Remainders!negative|EtSeq}%
The preceding example, then, takes the following
form:
\[
\begin{array}{@{\,}*{3}{r@{}}r@{\,}}
135&27&5&41 \\
31\underline{2}&\underline{31}&2&31 \\
\hline
   & 7& & 1 \\
   & 6& &12 \\
   &10& &10 \\
\cline{2-2}
   &23& & 3 \\
   &  & & 3 \\
\cline{4-4}
   &  & &29 \\
\multicolumn{2}{r}{\llap{\text{subtract}}} & &23 \\
\cline{4-4}
   &  & & 6
\end{array}
\]

I have placed a bar beneath the digits which are
to be taken negatively, and I have subtracted the sum
of the products of these numbers by those above them
from the sum of the other products.

The whole question, therefore, resolves itself into
finding for every divisor the remainders resulting from
dividing $1$, $10$, $100$, $1000$\Add{,~$\dots$} by that divisor. This can be
readily done by actual division; but it can be accomplished
more simply by the following consideration.
If $r$~be the remainder from the division of~$10$ by a
given divisor, $r^{2}$~will be the remainder from the division
of~$100$, the square of~$10$, by that divisor; and
consequently it will be necessary merely to subtract
the given divisor from~$r^{2}$ as many times as is requisite
to obtain a positive or negative remainder less than
\PageSep{37}
half of that divisor. Let $s$ be that remainder; we shall
then only have to multiply $s$~by~$r$, the remainder from
the division of~$10$, to obtain the remainder from the
division of~$1000$ by the given divisor, because $1000$~is
$100 × 10$, and so~on.

For example, dividing $10$ by~$7$ we have a remainder
of~$3$; hence, the remainder from dividing $100$ by~$7$
will be~$9$, or, subtracting from~$9$ the given divisor~$7$,~$2$.
The remainder from dividing $1000$ by~$7$, then, will
be the product of~$2$ by $3$~or~$6$, or, subtracting the divisor,~$7$,~$-1$.
Again, the remainder from dividing
%[** TN: Removed comma in 10,000 for consistency]
$\Typo{10,000}{10000}$ by~$7$ will be the product of $-1$~and~$3$, or~$-3$,
and so~on.

Let us now take the divisor~$11$. The remainder
\index{Eleven, the number, test of divisibility by}%
from dividing~$1$ by~$11$ is~$1$, from dividing~$10$ by~$11$ is~$10$,
\MNote{Test of divisibility by~$11$.}
or, subtracting the divisor,~$-1$. The remainder
from dividing~$100$ by~$11$, then, will be the square of~$-1$,
or~$1$; from dividing $1000$ by~$11$ it will be $1$~multiplied
by~$-1$ or\Add{ }$-1$~again, and so on forever, the remainders
forming the infinite series
\[
1,\quad -1,\quad 1,\quad -1,\quad 1,\quad -1,\dots\Add{.}
\]

Hence results the remarkable property of the number~$11$,
that if the digits of any number be alternately
added and subtracted, that is to say, if we take the
sum of the first, the third, and the fifth, etc., and subtract
from it the sum of the second, the fourth, the
sixth, etc., we shall obtain the remainder which results
from dividing that number by the number~$11$.
\PageSep{38}

The preceding theory of remainders is fraught
\index{Remainders!theory of}%
with remarkable consequences, and has given rise to
\MNote{Theory of remainders\Add{.}}
many ingenious and difficult investigations. We can
demonstrate, for example, that if the divisor is a prime
number, the remainders of any progression $1$, $a$, $a^{2}$,
$a^{3}$, $a^{4}$,~$\dots$ form periods which will recur continually
to infinity, and all of which, like the first, begin with
unity; in such wise that when unity reappears among
the remainders we may continue them to infinity by
simply repeating the remainders which precede. It
has also been demonstrated that these periods can
only contain a number of terms which is equal to the
divisor less~$1$ or to an aliquot part of the divisor less~$1$.
But we have not yet been able to determine \textit{à~priori}
this number for any divisor whatever.

As to the utility of this method for finding the remainder
\index{Theory of remainders, utility of the}%
resulting from dividing a given number by a
given divisor, it is frequently very useful when one
has several numbers to divide by the same number,
and it is required to prepare a table of the remainders.
While as to division by $9$~and~$11$, since that is very
simple, it can be employed as a check upon multiplication
and division. Having found the remainders
from dividing the multiplicand and the multiplier by
either of these numbers it is simply necessary to take
the product of the two remainders so resulting, from
which, after subtracting the divisor as many times as
is requisite, we shall obtain the remainder from dividing
their product by the given divisor,---a remainder
\PageSep{39}
which should agree with the remainder obtained
from treating the actual product in this manner. And
since in division the dividend less the remainder should
\MNote{checks on multiplication and division.}
\index{Checks on multiplication and division}%
be equal to the product of the divisor and the quotient,
the same check may also be applied here to advantage.

The supposition which I have just made that the
product of the remainders from dividing two numbers
by the same divisor is equal to the remainder from
dividing the product of these numbers by the same
divisor is easily proved, and I here give a general
demonstration of it.

Let $M$~and~$N$ be two numbers, $D$~the divisor, $p$~and~$q$
the quotients, and $r$,~$s$ the two remainders. We
shall plainly have
\[
M = pD + r,\quad
N = qD + s,
\]
from which by multiplying we obtain
\[
MN = pqD^{2} + spD + rqD + rs;
\]
where it will be seen that all the terms are divisible
by~$D$ with the exception of the last,~$rs$, whence it follows
that $rs$~will be the remainder from dividing~$MN$
by~$D$. It is further evident that if any multiple whatever
of~$D$, as~$mD$, be subtracted from~$rs$, the result
$rs - mD$ will also be the remainder from dividing~$MN$
by~$D$. For, putting the value of~$MN$ in the following
form:
\[
pqD^{2} + spD + rqD + mD + rs - mD,
\]
it is obvious that the remaining terms are all divisible
\PageSep{40}
by~$D$. And this remainder $rs - mD$ can always be
made less than~$D$, or, by employing negative remainders,
less even than~$\dfrac{D}{2}$.

This is all that I have to say upon multiplication
\MNote{Evolution.}
\index{Evolution}%
and division. I shall not speak of the \emph{extraction of
roots}. The rule is quite simple for square roots; it
leads directly to its goal; trials are unnecessary. As
to cube and higher roots, the occasion rarely arises
for extracting them, and when it does arise the extraction
can be performed with great facility by means
of logarithms, where the degree of exactitude can be
\index{Logarithms}%
carried to as many decimal places as the logarithms
themselves have decimal places. Thus, with seven-place
logarithms we can extract roots having seven
figures, and with the large tables where the logarithms
have been calculated to ten decimal places we
can obtain even ten figures of the result.

One of the most important operations in arithmetic
\index{Rule!three@of three|EtSeq}%
is the so-called \emph{rule of three}, which consists in
finding the fourth term of a proportion of which the
first three terms are given.

In the ordinary text-books of arithmetic this rule
has been unnecessarily complicated, having been divided
into simple, direct, inverse, and compound rules
of three. In general it is sufficient to comprehend the
conditions of the problem thoroughly, for the common
rule of three is always applicable where a quantity increases
or diminishes in the same proportion as another.
\PageSep{41}
For example, the price of things augments in
proportion to the quantity of the things, so that the
quantity of the thing being doubled, the price also
\MNote{Rule of three.}
will be doubled, and so on. Similarly, the amount of
work done increases proportionally to the number of
persons employed. Again, things may increase simultaneously
in two different proportions. For example,
the quantity of work done increases with the
number of the persons employed, and also with the
time during which they are employed. Further, there
are things that decrease as others increase.

Now all this may be embraced in a single, simple
proposition. If a quantity increases both in the ratio
in which one or several other quantities increase and
in that in which one or several other quantities decrease,
it is the same thing as saying that the proposed
quantity increases proportionally to the product of the
quantities which increase with it, divided by the product
of the quantities which simultaneously decrease.
For example, since the quantity of work done increases
proportionally with the number of laborers
\index{Laborers, work of}%
and with the time during which they work and since
it diminishes in proportion as the work becomes more
difficult, we may say that the result is proportional to
the number of laborers multiplied by the number
measuring the time during which they labor, divided
by the number which measures or expresses the difficulty
of the work.

The further fact should not be lost sight of that
\PageSep{42}
the rule of three is properly applicable only to things
which increase in a constant ratio. For example, it is
\index{Ratios, constant}%
\MNote{Applicability of the rule of three.}
assumed that if a man does a certain amount of work
in one day, two men will do twice that amount in one
day, three men three times that amount, four men
four times that amount,~etc. In reality this is not the
case, but in the rule of proportion it is assumed to be
such, since otherwise we should not be able to employ
it.

When the law of augmentation or diminution varies,
the rule of three is not applicable, and the ordinary
methods of arithmetic are found wanting. We
must then have recourse to algebra.

A cask of a certain capacity empties itself in a certain
\index{Efflux, law of}%
time. If we were to conclude from this that a
cask of double that capacity would empty itself in
double the time, we should be mistaken, for it will
empty itself in a much shorter time. The law of efflux
does not follow a constant ratio but a variable
ratio which diminishes with the quantity of liquid remaining
in the cask.

We know from mechanics that the spaces traversed
\index{Falling stone, spaces traversed by a}%
by a body in uniform motion bear a constant ratio to
the times elapsed. If we travel one mile in one hour,
in two hours we shall travel two miles. But the spaces
traversed by a falling stone are not in a fixed ratio to
the time. If it falls sixteen feet in the first second, it
will fall forty-eight feet in the second second.

The rule of three is applicable when the ratios are
\PageSep{43}
constant only. And in the majority of affairs of ordinary
life constant ratios are the rule. In general, the
price is always proportional to the quantity, so that if
\MNote{Theory and practice.}
\index{Practice, theory and}%
\index{Theory and practice}%
a given thing has a certain value, two such things will
have twice that value, three three times that value,
four four times that value,~etc. It is the same with
the product of labor relatively to the number of laborers
and to the duration of the labor. Nevertheless,
cases occur in which we may be easily led into error.
If two horses, for example, can pull a load of a certain
\index{Horses}%
weight, it is natural to suppose that four horses
could pull a load of double that weight, six horses a
load of three times that weight. Yet, strictly speaking,
such is not the case. For the inference is based
upon the assumption that the four horses pull alike in
amount and direction, which in practice can scarcely
ever be the case. It so happens that we are frequently
led in our reckonings to results which diverge widely
from reality. But the fault is not the fault of mathematics;
\index{Mathematics!exactness of}%
for mathematics always gives back to us exactly
what we have put into it. The ratio was constant
according to the supposition. The result is founded
upon that supposition. If the supposition is false the
result is necessarily false. Whenever it has been attempted
to charge mathematics with inexactitude, the
accusers have simply attributed to mathematics the
error of the calculator. False or inexact data having
been employed by him, the result also has been necessarily
false or inexact.
\PageSep{44}

Among the other rules of arithmetic there is one
called \emph{alligation} which deserves special consideration
\index{Alligation!generally|EtSeq}%
\MNote{Alligation.}
from the numerous applications which it has. Although
alligation is mainly used with reference to the
mingling of metals by fusion, it is yet applied generally
\index{Metals, mingling of, by fusion}%
to mixtures of any number of articles of different
values which are to be compounded into a whole of a
like number of parts having a mean value. The rule
\index{Mixtures, rule of|EtSeq}%
\index{Rule!mixtures@of mixtures|EtSeq}%
of alligation, or mixtures, accordingly, has two parts.

In one we seek the mean and common value of
each part of the mixture, having given the number
of the parts and the particular value of each. In the
second, having given the total number of the parts
and their mean value, we seek the composition of the
mixture itself, or the proportional number of parts of
each ingredient which must be mixed or alligated together.

Let us suppose, for example, that we have several
\index{Grain, of different prices}%
bushels of grain of different prices, and that we are
desirous of knowing the mean price. The mean price
must be such that if each bushel were of that price the
total price of all the bushels together would still be
the same. Whence it is easy to see that to find the
mean price in the present case we have first simply to
find the total price and to divide it by the number of
bushels.

In general if we multiply the number of things of
each kind by the value of the unit of that kind and
then divide the sum of all these products by the total
\PageSep{45}
number of things, we shall have the mean value, because
that value multiplied by the number of the
things will again give the total value of all the things
taken together.

This mean or average value as it is called, is of
\index{Mean values|EtSeq}%
\index{Values!mean|EtSeq}%
great utility in almost all the affairs of life. Whenever
\MNote{Mean values.}
we arrive at a number of different results, we
always like to reduce them to a mean or average expression
which will yield the same total result.

You will see when you come to the calculus of
\index{Probabilities, calculus of|EtSeq}%
probabilities that this science is almost entirely based
upon the principle we are discussing.

The registration of births and deaths has rendered
\index{Average life|EtSeq}%
\index{Life insurance|EtSeq}%
\index{Mortality, tables of}%
possible the construction of so-called \emph{tables of mortality}
which show what proportion of a given number of
children born at the same time or in the same year
survive at the end of one year, two years, three years,~etc.
So that we may ask upon this basis what is the
mean or average value of the life of a person at any
given age. If we look up in the tables the number of
people living at a certain age, and then add to this
the number of persons living at all subsequent ages,
it is clear that this sum will give the total number of
years which all living persons of the age in question
have still to live. Consequently, it is only necessary
to divide this sum by the number of living persons of
a certain age in order to obtain the average duration
of life of such persons, or better, the number of years
which each person must live that the total number of
\PageSep{46}
years lived by all shall be the same and that each
person shall have lived an equal number. It has been
\MNote{Probability of life.}
\index{Life, probability of}%
found in this manner by taking the mean of the results
of different tables of mortality, that for an infant
one year old the average duration of life is about
$40$~years; for a child ten years old it is still $40$~years;
for~$20$ it is~$34$; for~$30$ it is~$26$; for~$40$ it is~$23$; for~$50$
it is~$17$; for~$60$ it is~$12$; for~$70$,~$8$; and for~$80$,~$5$.

To take another example, a number of different
experiments are made. Three experiments have given~$4$
\index{Experiments!average of}%
as a result; two experiments have given~$5$; and one
has given~$6$. To find the mean we multiply~$4$ by~$3$, $5$~by~$2$,
and $1$~by~$6$, add the products which gives~$28$,
and divide~$28$ by the number of experiments or~$6$,
which gives~$4\frac{2}{3}$ as the mean result of all the experiments.

But it will be apparent that this result can be regarded
as exact only upon the condition of our having
supposed that the experiments were all conducted with
equal precision. But it is impossible that such could
have been the case, and it is consequently imperative
to take account of these inequalities, a requirement
which would demand a far more complicated calculus
than that which we have employed, and one which is
now engaging the attention of mathematicians.

The foregoing is the substance of the first part of
the rule of alligation; the second part is the opposite
of the first. Given the mean value, to find how much
\PageSep{47}
must be taken of each ingredient to produce the required
mean value.

The problems of the first class are always determinate,
because, as we have just seen, the number of
\MNote{Alternate alligation.}
\index{Alligation!alternate}%
units of each ingredient has simply to be multiplied
by the value of each ingredient and the sum of all
these products divided by the number of the ingredients.

The problems of the second class, on the other
\index{Analysis!indeterminate|EtSeq}%
\index{Indeterminate analysis|EtSeq}%
hand, are always indeterminate. But the condition
that only positive whole numbers shall be admitted
in the result serves to limit the number of the solutions.

Suppose that we have two kinds of things, that
the value of the unit of one kind is~$a$, and that of the
unit of the second is~$b$, and that it is required to find
how many units of the first kind and how many units
of the second must be taken to form a mixture or
whole of which the mean value shall be~$m$.

Call $x$~the number of units of the first kind that
must enter into the mixture, and $y$~the number of units
of the second kind. It is clear that $ax$~will be the
value of the $x$~units of the first kind, and $by$~the value
of the $y$~units of the second. Hence $ax + by$ will be
the total value of the mixture. But the mean value
of the mixture being by supposition~$m$, the sum~$x + y$
of the units of the mixture multiplied by~$m$, the mean
value of each unit, must give the same total value.
We shall have, therefore, the equation
\PageSep{48}
\[
ax + by = mx + my.
\]
Transposing to one side the terms multiplied by~$x$
and to the other the terms multiplied by~$y$, we obtain:
\MNote{Two ingredients.}
\index{Ingredients}%
\[
(a - m)x = (m - b)y,
\]
and dividing by~$a - m$ we get
\[
x = \frac{(m - b)y}{a - m},
\]
whence it appears that the number~$y$ may be taken at
pleasure, for whatever be the value given to~$y$, there
will always be a corresponding value of~$x$ which will
satisfy the problem.

Such is the general solution which algebra gives.
But if the condition be added that the two numbers $x$~and~$y$
shall be integers, then $y$~may not be taken at
pleasure. In order to see how we can satisfy this last
condition in the simplest manner, let us divide the
last equation by~$y$, and we shall have
\[
\frac{x}{y} = \frac{m - b}{a - m}.
\]
For $x$~and~$y$ both to be positive, it is necessary that
the quantities
\[
m - b \quad\text{and}\quad a - m
\]
should both have the same sign; that is to say, if $a$~is
greater or less than~$m$, then conversely $b$~must be less
or greater than~$m$; or again, $m$~must lie between $a$~and~$b$,
which is evident from the condition of the
problem. Suppose $a$, then, to be the greater and $b$~the
\PageSep{49}
smaller of the two prices. It remains to find the
value of the fraction
\MNote{Rule of mixtures.}
\index{Mixtures, rule of}%
\[
\frac{m - b}{a - m},
\]
which if necessary is to be reduced to its lowest terms.
Let~$\dfrac{B}{A}$ be that fraction reduced to its lowest terms. It
is clear that the simplest solution will be that in which
\[
x = B \quad\text{and}\quad y = A.
\]
But since a fraction is not altered by multiplying its
numerator and denominator by the same number, it
is clear that we may also take $x = nB$ and $y = nA$, $n$~being
any number whatever, provided it is an integer,
for by supposition $x$~and~$y$ must be integers. And it
is easy to prove that these expressions of $x$~and~$y$ are
the only ones which will resolve the proposed problem.
According to the ordinary rule of mixtures, $x$,
the quantity of the dearer ingredient, is made equal
to~$m - b$, the excess of the average price above the
lower price, and $y$~the quantity of the cheaper ingredient
is made equal to~$a - m$, the excess of the higher
price above the average price,---a rule which is contained
directly in the general solution above given.

Suppose, now, that instead of two kinds of things,
we have three kinds, the values of which beginning
with the highest are $a$,~$b$, and~$c$. Let $x$,~$y$,~$z$ be the
quantities which must be taken of each to form a mixture
or compound having the mean value~$m$. The
sum of the values of the three quantities $x$,~$y$,~$z$ will
then be
\[
ax + by + cz.
\]
\PageSep{50}
But this total value must be the same as that produced
if all the individual values were~$m$, in which
\MNote{Three ingredients.}
case the total value is obviously
\[
mx + my + mz.
\]
The following equation, therefore, must be satisfied:
\[
ax + by + cz = mx + my + mz,
\]
or, more simply,
\[
(a - m)x + (b - m)y + (c - m)z = 0.
\]
Since there are three unknown quantities in this equation,
two of them may be taken at pleasure. But if
the condition is that they shall be expressed by positive
integers, it is to be observed first that the numbers
\[
a - m \quad\text{and}\quad m - c
\]
are necessarily positive; so that putting the equation
in the form
\[
(a - m)x - (m - c)z = (m - b)y,
\]
the question resolves itself into finding two multiples
of the given numbers
\[
a - m \quad\text{and}\quad m - c
\]
whose difference shall be equal to~$(m - b)y$.

This question is always resolvable in whole numbers
whatever the given numbers be of which we seek
the multiples, and whatever be the difference between
these multiples. As it is sufficiently remarkable in itself
and may be of utility in many emergencies, we
shall give here a general solution of it, derived from
the properties of continued fractions.
\index{Continued fractions, solution of alligation by|EtSeq}%
\PageSep{51}

Let $M$~and~$N$ be two whole numbers. Of these
numbers two multiples $xM$,~$zN$ are sought whose difference
is given and equal to~$D$. The following equation
\MNote{General solution.}
will then have to be satisfied
\[
xM - zN = D,
\]
where $x$~and~$z$ by supposition are whole numbers. In
the first place, it is plain that if $M$~and~$N$ are not
prime to each other, the number~$D$ is divisible by the
greatest common divisor of $M$~and~$N$; and the division
having been performed, we should have a similar
equation in which the numbers $M$~and~$N$ are prime
to each other, so that we are at liberty always to suppose
them reduced to that condition. I now observe
that if we know the solution of the equation for the
case in which the number~$D$ is equal to $+1$~or~$-1$,
we can deduce the solution of it for any value whatever
of~$D$. For example, suppose that we know two
multiples of $M$~and~$N$, say $pM$~and~$qN$, the difference
of which $pM - qN$ is equal to~$±1$. Then obviously
we shall merely have to multiply both these multiples
by the number~$D$ to obtain a difference equal to~$±D$.
For, multiplying the preceding equation by~$D$, we
have
\[
pDM - qDN = ±D;
\]
and subtracting the latter equation from the original
equation
\[
xM - zN = D,
\]
or adding it, according as the term~$D$ has the sign
$+$~or~$-$ before it, we obtain
\PageSep{52}
\[
(x \mp pD)M - (z \mp qD)N = 0,
\]
which gives at once, as we saw above in the rule for
the mixture of two different ingredients,
\MNote{Development.}
\[
(x \mp pD) = nN,\quad
(z \mp qD) = nM,
\]
$n$~being any number whatever. So that we have generally
\[
x = nN ± pD \quad\text{and}\quad z = nM ± qD
\]
where $n$~is any whole number, positive or negative.
It remains merely to find two numbers $p$~and~$q$ such
that
\[
pM - qN = ±1.
\]
Now this question is easily resolvable by continued
fractions. For we have seen in treating of these fractions
that if the fraction~$\dfrac{M}{N}$ be reduced to a continued
fraction, and all the successive fractions approximating
to its value be calculated, the last of these successive
fractions being the fraction~$\dfrac{M}{N}$ itself, then the series
of fractions so reached is such that the difference
between any two consecutive fractions is always equal
to a fraction of which the numerator is unity and the
denominator the product of the two denominators.
For example, designating by~$\dfrac{K}{L}$ the fraction which
immediately precedes the last fraction~$\dfrac{M}{N}$ we obtain
necessarily
\[
LM - KN = 1 \quad\text{or}\quad -1,
\]
according as $\dfrac{M}{N}$~is greater or less than~$\dfrac{K}{L}$, in other
\PageSep{53}
words, according as the place occupied by the last
fraction~$\dfrac{M}{N}$ in the series of fractions successively approximating
to its value is even or odd; for, the first
\MNote{Resolution by continued fractions.}
fraction of the approximating series is always smaller,
the second larger, the third smaller,~etc., than the
original fraction which is identical with the last fraction
of the series. Making, therefore,
\[
p = L \quad\text{and}\quad q = K,
\]
the problem of the two multiples will be resolved in
all its generality.

It is now clear that in order to apply the foregoing
solution to the initial question regarding alligation we
have simply to put
\[
M = a - m,\quad N = m - c, \quad\text{and}\quad D = (m - b)y;
\]
so that the number~$y$ remains undetermined and may
be taken at pleasure, as may also the number~$N$ which
appears in the expressions for $x$~and~$z$.
\PageSep{54}


\Lecture[On Algebra.]{III.}{On Algebra, Particularly the Resolution of
Equations of the Third and
Fourth Degree.}
\index{Algebra!history of|EtSeq}%
\index{Diophantus|EtSeq}%
\index{Geometers, ancient|EtSeq}%
\index{Greeks, mathematics of the|EtSeq}%
\index{Romans, mathematics of the}%
\PgLabel{54}

\First{Algebra} is a science almost entirely due to the
moderns. I say almost entirely, for we have
\MNote{Algebra among the ancients.}
one treatise from the Greeks, that of Diophantus, who
flourished in the third\footnote
  {The period is uncertain. Some say in the fourth century. See Cantor,
  \index{Cantor|FN}%
  \textit{Geschichte der Mathematik}, 2nd.~ed., Vol.~I., p.~434.---\textit{Trans.}}
century of the Christian era.
This work is the only one which we owe to the ancients
in this branch of mathematics. When I speak
of the ancients I speak of the Greeks only, for the
Romans have left nothing in the sciences, and to all
appearances did nothing.

Diophantus may be regarded as the inventor of
algebra.\footnote
  {On this point, see \textit{Appendix}, \PgRef{151}.---\textit{Trans.}}
From a word in his preface, or rather in his
letter of dedication, (for the ancient geometers were
wont to address their productions to certain of their
friends, a practice exemplified in the prefaces of Apollonius
\index{Apollonius}%
and Archimedes), from a word in his preface, I
\index{Archimedes}%
say, we learn that he was the first to occupy himself
\PageSep{55}
with that branch of arithmetic which has since been
called algebra.

His work contains the first elements of this science.
He employed to express the unknown quantity a Greek
\index{Unknown quantity}%
\MNote{Diophantus\Add{.}}
letter which corresponds to our~$st$\footnote
  {According to a recent conjecture, the character in question is an abbreviation
  of~\textgreek{ar} the first letters of \textgreek{>arijm'os}, \textit{number}, the appellation technically
  applied by Diophantus to the unknown quantity.---\textit{Trans.}}
and which has
been replaced in the translations by~$N$. To express
the known quantities he employed numbers solely, for
algebra was long destined to be restricted entirely to
the solution of numerical problems. We find, however,
that in setting up his equations consonantly with
the conditions of the problem he uses the known and
the unknown quantities alike. And herein consists
\index{Algebra!essence of}%
virtually the essence of algebra, which is to employ
unknown quantities, to calculate with them as we do
with known quantities, and to form from them one
or several equations from which the value of the unknown
quantities can be determined. Although the
work of Diophantus contains indeterminate problems
\index{Analysis!indeterminate}%
\index{Indeterminate analysis}%
almost exclusively, the solution of which he seeks in
rational numbers,---problems which have been designated
after him \emph{Diophantine problems},---we nevertheless
\index{Diophantine problems}%
find in his work the solution of a number of determinate
problems of the first degree, and even of such
as involve several unknown quantities. In the latter
case, however, the author invariably has recourse to
particular artifices for reducing the problem to a single
unknown quantity,---which is not difficult. He gives,
\PageSep{56}
also, the solution of \emph{equations of the second degree}, but
\index{Equations!second@of the second degree}%
is careful so to arrange them that they never assume
the affected form containing the square and the first
power of the unknown quantity.

He proposed, for example, the following question
\MNote{Equations of the second degree.}
which involves the general theory of equations of the
second degree:

\textit{To find two numbers the sum and the product of which
are given.}
\index{Sum and difference, of two numbers}%

If we call the sum~$a$ and the product~$b$ we have at
once, by the theory of equations, the equation
\[
x^{2} - ax + b = 0.
\]

Diophantus resolves this problem in the following
manner. The sum of the two numbers being given,
he seeks their difference, and takes the latter as the
unknown quantity. He then expresses the two numbers
in terms of their sum and difference,---the one
by half the sum plus half the difference, the other by
half the sum less half the difference,---and he has
then simply to satisfy the other condition by equating
their product to the given number. Calling the given
sum~$a$, the unknown difference~$x$, one of the numbers
will be~$\dfrac{a + x}{2}$ and the other will be~$\dfrac{a - x}{2}$. Multiplying
these together we have~$\dfrac{a^{2} - x^{2}}{4}$. The term containing~$x$
is here eliminated, and equating the quantity
last obtained to the given product, we have the
simple equation
\[
\frac{a^{2} - x^{2}}{4} = b,
\]
\PageSep{57}
from which we obtain
\[
x^{2} = a^{2} - 4b,
\]
and from the latter
\[
x = \sqrt{a^{2} - 4b}.
\]

Diophantus resolves several other problems of this
class. By appropriately treating the sum or difference
\MNote{Other problems solved by Diophantus.}
as the unknown quantity he always arrives at an
equation in which he has only to extract a square root
to reach the solution of his problem.

But in the books which have come down to us
(for the entire work of Diophantus has not been preserved)
this author does not proceed beyond equations
of the second degree, and we do not know if he
or any of his successors (for no other work on this
subject has been handed down from antiquity) ever
pushed their researches beyond this point.

I have still to remark in connexion with the work
\index{Signs $+$ and $-$}%
of Diophantus that he enunciated the principle that
$+$~and~$-$ give~$-$ in multiplication, and $-$~and~$-$,~$+$,
in the form of a definition. But I am of opinion that
this is an error of the copyists, since he is more likely
to have considered it as an axiom, as did Euclid some
\index{Euclid}%
of the principles of geometry. However that may be,
it will be seen that Diophantus regarded the rule of
the signs as a self-evident principle not in need of demonstration.

The work of Diophantus is of incalculable value
from its containing the first germs of a science which
because of the enormous progress which it has since
\PageSep{58}
made constitutes one of the chiefest glories of the human
intellect. Diophantus was not known in Europe
\MNote{Translations of Diophantus\Add{.}}
until the end of the sixteenth century, the first translation
having been a wretched one by Xylander made
\index{Xylander}%
in~1575 and based upon a manuscript found about the
middle of the sixteenth century in the Vatican library,
\index{Vatican library}%
where it had probably been carried from Greece when
the Turks took possession of Constantinople.
\index{Constantinople}%
\index{Turks}%

Bachet de Méziriac, one of the earliest members
\index{Bachet de Méziriac}%
\index{Meziriac@Méziriac, Bachet de}%
of the French Academy, and a tolerably good mathematician
for his time, subsequently published~(1621)
a new translation of the work of Diophantus accompanied
by lengthy commentaries, now superfluous.
Bachet's translation was afterwards reprinted with observations
and notes by Fermat, one of the most celebrated
\index{Fermat}%
mathematicians of France, who flourished
\index{France}%
about the middle of the seventeenth century, and of
whom we shall have occasion to speak in the sequel
for the important discoveries which he has made in
analysis. Fermat's edition bears the date of~1670.\footnote
  {There have since been published a new critical edition of the text by
  M.~Paul Tannery (Leipsic, 1893), and two German translations, one by O.~Schulz
  \index{Tannery, M. Paul|FN}%
  \index{Wertheim, G.|FN}%
  (Berlin, 1822) and one by G.~Wertheim (Leipsic, 1890). Fermat's notes
  on Diophantus have been republished in Vol.~I. of the new edition of Fermat's
  works (Paris, Gauthier-Villars et Fils, 1891).---\textit{Trans.}}

It is much to be desired that good translations
\index{Geometers, ancient}%
should be made, not only of the work of Diophantus,
but also of the small number of other mathematical
works which the Greeks have left us.\footnote
  {Since Lagrange's time this want has been partly supplied. Not to mention
  Euclid, we have, for example, of Archimedes the German translation of
  \index{Archimedes|FN}%
  Nizze (Stralsund, 1824) and the French translation of Peyrard (Paris, 1807); of
  \index{Nizze|FN}%
  \index{Peyrard}%
  Apollonius, several translations; also modern translations of Hero, Ptolemy,
  \index{Apollonius}%
  \index{Geometers, ancient}%
  \index{Hero}%
  \index{Pappus}%
  \index{Proclus}%
  \index{Ptolemy}%
  \index{Theon}%
  Pappus, Theon, Proclus, and several others.}
\PageSep{59}

Prior to the discovery and publication of Diophantus,
however, algebra had already found its way into
\index{Algebra!name@the name of}%
\index{Algebra!among the Arabs|EtSeq}%
Europe. Towards the end of the fifteenth century
there appeared in Venice a work by an Italian Franciscan
monk named Lucas Paciolus on arithmetic and
\index{Paciolus, Lucas}%
geometry in which the elementary rules of algebra
were stated. This book was published (1494) in the
\MNote{Algebra among the Arabs.}
\index{Arabs!Algebra among the|EtSeq}%
early days of the invention of printing, and the fact
\index{Printing, invention of}%
that the name of \emph{algebra} was given to the new science
shows clearly that it came from the Arabs. It is true
that the signification of this Arabic word is still disputed,
but we shall not stop to discuss such matters,
for they are foreign to our purpose. Let it suffice
that the word has become the name for a science that
is universally known, and that there is not the slightest
ambiguity concerning its meaning, since up to the
present time it has never been employed to designate
anything else.

We do not know whether the Arabs invented algebra
\PgLabel{59}
themselves or whether they took it from the
Greeks.\footnote
  {See Appendix, \PgRef{152}.}
There is reason to believe that they possessed
the work of Diophantus, for when the ages of
barbarism and ignorance which followed their first
conquests had passed by, they began to devote themselves
to the sciences and to translate into Arabic all
the Greek works which treated of scientific subjects.
It is reasonable to suppose, therefore, that they also
\PageSep{60}
translated the work of Diophantus and that the same
work stimulated them to push their inquiries farther
in this science.

Be that as it may, the Europeans, having received
\MNote{Algebra in Europe.}
\index{Algebra!Europe@in Europe}%
\index{Europe, algebra in}%
algebra from the Arabs, were in possession of it one
hundred years before the work of Diophantus was
known to them. They made, however, no progress
beyond equations of the first and second degree. In
\index{Equations!third@of the third degree}%
the work of Paciolus, which we mentioned above, the
\index{Paciolus, Lucas}%
general resolution of equations of the second degree,
such as we now have it, was not given. We find in
this work simply rules, expressed in bad Latin verses,
for resolving each particular case according to the
different combinations of the signs of the terms of
equation, and even these rules applied only to the
case where the roots were real and positive. Negative
\index{Negative roots}%
\index{Roots!negative}%
roots were still regarded as meaningless and superfluous.
It was geometry really that suggested to us the
\index{Geometry}%
use of negative quantities, and herein consists one of
the greatest advantages that have resulted from the
application of algebra to geometry,---a step which we
owe to Descartes.
\index{Descartes}%
\PgLabel{60}

In the subsequent period the resolution of \emph{equations
of the third degree} was investigated and the discovery
for a particular case ultimately made by a mathematician
\index{Ferrous, Scipio|EtSeq}%
of Bologna named Scipio Ferreus (1515).\footnote
  {The date is uncertain. Tartaglia gives 1506, Cardan 1515. Cantor prefers
  \index{Cantor|FN}%
  \index{Cardan}%
  \index{Tartaglia}%
  the latter.---\textit{Trans.}}
Two
other Italian mathematicians, Tartaglia and Cardan,
\PageSep{61}
subsequently perfected the solution of Ferreus and
rendered it general for all equations of the third degree.
At this period, Italy, which was the cradle of
\index{Italy, cradle of algebra in Europe}%
\MNote{Tartaglia (1500--1559). Cardan (1501--1576).}
\index{Cardan}%
\index{Tartaglia}%
algebra in Europe, was still almost the sole cultivator
of the science, and it was not until about the middle
of the sixteenth century that treatises on algebra began
to appear in France, Germany, and other countries.
\index{France}%
\index{Germany}%
The works of Peletier and Buteo were the first
\index{Buteo}%
\index{Peletier}%
which France produced in this science, the treatise of
the former having been printed in~1554 and that of
the latter in~1559.

Tartaglia expounded his solution in bad Italian
verses in a work treating of divers questions and inventions
printed in~1546, a work which enjoys the
distinction of being one of the first to treat of modern
fortifications by bastions.

About the same time (1545) Cardan published his
treatise \textit{Ars Magna}, or \textit{Algebra}, in which he left
scarcely anything to be desired in the resolution of
equations of the third degree. Cardan was the first to
perceive that equations had several roots and to distinguish
them into positive and negative. But he is
particularly known for having first remarked the so-called
\emph{irreducible case} in which the expression of the
\index{Irreducible case}%
real roots appears in an imaginary form. Cardan convinced
himself from several special cases in which the
equation had rational divisors that the imaginary form
did not prevent the roots from having a real value.
But it remained to be proved that not only were the
\PageSep{62}
roots real in the irreducible case, but that it was impossible
for all three together to be real except in that
case. This proof was afterwards supplied by Vieta,
\index{Vieta}%
and particularly by Albert Girard, from considerations
\index{Girard, Albert}%
touching the trisection of an angle.
\index{Angle, trisection of an}%
\index{Trisection of an angle}%

We shall revert later on to the \emph{irreducible case of
equations of the third degree}, not solely because it presents
\MNote{The irreducible case.}
a new form of algebraical expressions which
have found extensive application in analysis, but because
it is constantly giving rise to unprofitable inquiries
with a view to reducing the imaginary form to
a real form and because it thus presents in algebra a
problem which may be placed upon the same footing
with the famous problems of the duplication of the
\index{Problems!solution@for solution}%
cube and the squaring of the circle in geometry.
\index{Circle!squaring of the}%
\index{Cube, duplication of the}%
\index{Squaring of the circle}%

The mathematicians of the period under discussion
\index{Academies, rise of}%
were wont to propound to one another problems
for solution. These problems were in the nature of
public challenges and served to excite and to maintain
in the minds of thinkers that fermentation which
is necessary for the pursuit of science. The challenges
in question were continued down to the beginning of
the eighteenth century by the foremost mathematicians
of Europe, and really did not cease until the rise
of the Academies which fulfilled the same end in a
manner even more conducive to the progress of science,
partly by the union of the knowledge of their
various members, partly by the intercourse which they
maintained between them, and not least by the publication
\PageSep{63}
of their memoirs, which served to disseminate
the new discoveries and observations among all persons
interested in science.

The challenges of which we speak supplied in a
\index{Academies, rise of}%
measure the lack of Academies, which were not yet
\MNote{Biquadratic equations.}
\index{Biquadratic equations}%
\index{Equations!fourth@of the fourth degree}%
in existence, and we owe to these passages at arms
many important discoveries in analysis. Such was
the resolution of \emph{equations of the fourth degree}, which
was propounded in the following problem.

%[** TN: Next paragraph centered in the original]
\textit{To find three numbers in continued proportion of which
the sum is~$10$, and the product of the first two~$6$.}

Generalising and calling the sum of the three numbers~$a$,
the product of the first two~$b$, and the first two
numbers themselves $x$,~$y$, we shall have, first, $xy = b$.
Owing to the continued proportion, the third number
will then be expressed by~$\dfrac{y^{2}}{x}$, so that the remaining
condition will give
\[
x + y + \frac{y^{2}}{x} = a.
\]
From the first equation we obtain $x = \dfrac{b}{y}$, which substituted
in the second gives
\[
\frac{b}{y} + y + \frac{y^{2}}{b} = a\Typo{,}{.}
\]
Removing the fractions and arranging the terms, we
get finally
\[
y^{4} + by^{2} - aby + b^{2} = 0,
\]
an equation of the fourth degree with the second term
missing.

According to Bombelli, of whom we shall speak
\index{Bombelli}%
\PageSep{64}
again, Louis Ferrari of Bologna resolved the problem
\index{Ferrari, Louis}%
by a highly ingenious method, which consists in
\MNote{Ferrari (1522-1565). Bombelli.}
\index{Bombelli}%
dividing the equation into two parts both of which
permit of the extraction of the square root. To do
this it is necessary to add to the two numbers quantities
whose determination depends on an equation of
the third degree, so that the resolution of equations
\index{Equations!fifth@of the fifth degree}%
of the fourth degree depends upon the resolution of
equations of the third and is therefore subject to the
same drawbacks of the irreducible case.

The \textit{Algebra} of Bombelli was printed in Bologna
\index{Algebra!Italy@in Italy}%
in~1579\footnote
  {This was the second edition. The first edition appeared in Venice in~1572.---\textit{Trans.}}
in the Italian language. It contains not only
the discovery of Ferrari but also divers other important
remarks on equations of the second and third
degree and particularly on the theory of radicals by
means of which the author succeeded in several cases
in extracting the imaginary cube roots of the two
binomials of the formula of the third degree in the irreducible
case, so finding a perfectly real result and
furnishing thus the most direct proof possible of the
reality of this species of expressions.

Such is a succinct history of the first progress of
algebra in Italy. The solution of equations of the
\index{Italy, cradle of algebra in Europe}%
third and fourth degree was quickly accomplished.
But the successive efforts of mathematicians for over
two centuries have not succeeded in surmounting the
difficulties of the equation of the fifth degree.
\PageSep{65}

Yet these efforts are far from having been in vain.
They have given rise to the many beautiful theorems
which we possess on the formation of equations, on
\MNote{Theory of equations.}
\index{Equations!theory of}%
the character and signs of the roots, on the transformation
of a given equation into others of which the
roots may be formed at pleasure from the roots of the
given equation, and finally, to the beautiful considerations
concerning the metaphysics of the resolution
of equations from which the most direct method of
arriving at their solution, when possible, has resulted.
All this has been presented to you in previous lectures
and would leave nothing to be desired if it were
but applicable to the resolution of equations of higher
degree.

Vieta and Descartes in France, Harriot in England,
\index{Descartes}%
\index{Harriot}%
\index{Vieta}%
and Hudde in Holland, were the first after the
\index{Hudde}%
Italians whom we have just mentioned to perfect the
theory of equations, and since their time there is
scarcely a mathematician of note that has not applied
himself to its investigation, so that in its present state
this theory is the result of so many different inquiries
that it is difficult in the extreme to assign the author
of each of the numerous discoveries which constitute it.

I promised to revert to the irreducible case. To
\index{Irreducible case}%
this end it will be necessary to recall the method
which seems to have led to the original resolution of
equations of the third degree and which is still employed
in the majority of the treatises on algebra.
\PageSep{66}
Let us consider the general equation of the third degree
deprived of its second term, which can always be
removed; in a word, let us consider the equation
\MNote{Equations of the third degree.}
\index{Equations!third@of the third degree}%
\[
x^{3} + px + q = 0.
\]
Suppose
\[
x = y + z,
\]
where $y$~and~$z$ are two new unknown quantities, of
which one consequently may be taken at pleasure and
determined as we think most convenient. Substituting
this value for~$x$, we obtain \emph{the transformed equation}
\[
y^{3} + 3y^{2}z + 3yz^{2} + z^{3} + p(y + z) + q = 0.
\]
Factoring the two terms $3y^{2}z + 3yz^{2}$ we get
\[
3yz(y + z),
\]
and the transformed equation may be written as follows:
\[
y^{3} + z^{3} + (3yz + p)(y + z) + q = 0.
\]
Putting the factor multiplying $y + z$ equal to zero,---which
is permissible owing to the two undetermined
quantities involved,---we shall have the two equations
\[
3yz + p = 0\Typo{.}{}
\]
and
\[
y^{3} + z^{3} + q = 0\Typo{.}{,}
\]
from which $y$~and~$z$ can be determined. The means
which most naturally suggests itself to this end is to
take from the first equation the value of~$z$,
\[
z = -\frac{p}{3y},
\]
and to substitute it in the second equation, removing
the fractions by multiplication. So proceeding, we
\PageSep{67}
obtain the following equation of the sixth degree in~$y$,
called \emph{the reduced equation},
\MNote{The reduced equation.}
\[
y^{6} + qy^{3} - \frac{p^{3}}{27} = 0,
\]
which, since it contains two powers only of the unknown
quantity, of which one is the square of the
other, is resolvable after the manner of equations of
the second degree and gives immediately
\[
y^{3} = -\frac{q}{2} + \sqrt{\frac{q^{2}}{4} + \frac{p^{3}}{27}},
\]
from which, by extracting the cube root, we get
\[
y = \sqrt[3]{-\frac{q}{2} + \sqrt{\frac{q^{2}}{4} + \frac{p^{3}}{27}}},
\]
and finally,
\[
x = y + z = y - \frac{p}{3y}\Add{.}
\]
This expression for~$x$ may be simplified by remarking
that the product of~$y$ by the radical
\[
 \sqrt[3]{-\frac{q}{2} - \sqrt{\frac{q^{2}}{4} + \frac{p^{3}}{27}}}\Add{,}
\]
supposing all the quantities under the sign to be multiplied
together, is
\[
\sqrt[3]{-\frac{p^{3}}{27}} = -\frac{p}{3}.
\]
The term $\dfrac{p}{3y}$, accordingly, takes the form
\[
-\sqrt[3]{-\frac{q}{2} - \sqrt{\frac{q^{2}}{4} + \frac{p^{3}}{27}}},
\]
and we have
\[
x = \sqrt[3]{-\frac{q}{2} + \sqrt{\frac{q^{2}}{4} + \frac{p^{3}}{27}}}
  + \sqrt[3]{-\frac{q}{2} - \sqrt{\frac{q^{2}}{4} + \frac{p^{3}}{27}}},
\]
\PageSep{68}
an expression in which the square root underneath the
cubic radical occurs in both its plus and minus forms
and where consequently there can, on this score, be
no occasion for ambiguity.

This last expression is known as the \emph{Rule of Cardan},
\index{Cardan}%
\index{Rule!Cardan's}%
\MNote{Cardan's rule.}
and there has hitherto been no method devised
for the resolution of equations of the third degree
which does not lead to it. Since cubic radicals naturally
present but a single value, it was long thought
that Cardan's rule could give but one of the roots of
the equation, and that in order to find the two others
we must have recourse to the original equation and divide
it by~$x - a$, $a$~being the first root found. The
resulting quotient being an equation of the second degree
may be resolved in the usual manner. The division
in question is not only always possible, but it is
also very easy to perform. For in the case we are
considering the equation being
\[
x^{3} + px + q = 0,
\]
if $a$~is one of the roots we shall have
\[
a^{3} + pa + q = 0,
\]
which subtracted from the preceding will give
\[
x^{3} - a^{3} + p(x - a) = 0,
\]
a quantity divisible by~$x - a$ and having as its resulting
quotient
\[
x^{2} + ax + a^{2} + p = 0;
\]
so that the new equation which is to be resolved for
finding the two other roots will be
\PageSep{69}
\[
x^{2} + ax + a^{2} + p = 0,
\]
from which we have at once
\[
x = -\frac{a}{2} ± \sqrt{-p - \frac{3a^{2}}{4}}.
\]

I see by the \textit{Algebra} of Clairaut, printed in~1746,
\index{Clairaut}%
and by D'Alem\-bert's article on the \emph{Irreducible Case} in
\index{Irreducible case}%
\MNote{The generality of algebra.}
the first \textit{Encyclopædia} that the idea referred to prevailed
even in that period. But it would be the height
of injustice to algebra to accuse it of not yielding results
\index{Algebra!generality@the generality of}%
which were possessed of all the generality of
which the question was susceptible. The sole requisite
is to be able to read the peculiar hand-writing
\index{Algebra!hand-writing of}%
\index{Hand-writing of algebra}%
of algebra, and we shall then be able to see in it everything
which by its nature it can be made to contain.
In the case which we are considering it was forgotten
that every cube root may have three values, as every
square root has two. For the extraction of the cube
root of~$a$ for example is merely equivalent to the resolution
of the equation of the third degree $x^{3} - a = 0$.
Making $x = y\sqrt[3]{a}$, this last equation passes into the
simpler form $y^{3} - 1 = 0$, which has the root $y = 1$.
Then dividing by~$y - 1$ we have
\[
y^{2} + y + 1 = 0,
\]
from which we deduce directly the two other roots
\[
y = \frac{-1 ± \sqrt{-3}}{2}.
\]
These three roots, accordingly, are the three cube
roots of unity, and they may be made to give the three
cube roots of any other quantity~$a$ by multiplying
\PageSep{70}
them by the ordinary cube root of that quantity. It
is the same with roots of the fourth, the fifth, and all
the following degrees. For brevity, let us designate
the two roots
\MNote{The three cube roots of a quantity.}
\index{Cube roots of a quantity, the three}%
\[
\frac{-1 + \sqrt{-3}}{2} \quad\text{and}\quad \frac{-1 - \sqrt{-3}}{2}\Typo{,}{}
\]
by $m$~and~$n$. It will be seen that they are imaginary,
although their cube is real and equal to~$1$, as we may
readily convince ourselves by raising them to the
third power. We have, therefore, for the three cube
roots of~$a$,
\[
\sqrt[3]{a},\quad m\sqrt[3]{a},\quad n\sqrt[3]{a}.
\]

Now, in the resolution of the equation of the third
degree above considered, on coming to the reduced
expression $y^{3} = A$, where for brevity we suppose
\[
A = -\frac{q}{2} + \sqrt{\frac{q^{2}}{4} + \frac{p^{3}}{27}},
\]
we deduced the following result only:
\[
y = \sqrt[3]{A}.
\]
But from what we have just seen, it is clear that we
shall have not only
\[
y = \sqrt[3]{A},
\]
but also
\[
y = m\sqrt[3]{A} \quad\text{and}\quad y = n\sqrt[3]{A}.
\]
The root~$x$ of the equation of the third degree which
we found equal to
\[
y - \frac{p}{3y},
\]
will therefore have the three following values
\PageSep{71}
\[
\sqrt[3]{A} - \frac{p}{3\sqrt[3]{A}},\quad
m\sqrt[3]{A} - \frac{p}{3m\sqrt[3]{A}},\quad
n\sqrt[3]{A} - \frac{p}{3n\sqrt[3]{A}},
\]
\MNote{The roots of equations of the third degree.}
\index{Roots!equations@of equations of the third degree}%
\index{Third degree, equations of the}%
which will be the three roots of the equation proposed.
But making
\[
B = -\frac{q}{2} - \sqrt{\frac{q^{2}}{4} + \frac{p^{3}}{27}},
\]
it is clear that
\[
AB = -\frac{p^{3}}{27},
\]
whence
\[
\sqrt[3]{A} × \sqrt[3]{B} = -\frac{p}{3}.
\]
Substituting $\sqrt[3]{B}$ for $-\dfrac{p}{3\sqrt[3]{A}}$, and remarking that
$mn = 1$, and that consequently
\[
\frac{1}{m} = n,\quad \frac{1}{n} = m,
\]
the three roots which we are considering will be expressed
as follows:
%[** TN: Set on two lines in the original]
\[
x = \sqrt[3]{A} + \sqrt[3]{B},\quad
x = m\sqrt[3]{A} + n\sqrt[3]{B},\quad
x = n\sqrt[3]{A} + m\sqrt[3]{B}.
\]

We see, accordingly, that when properly understood
the ordinary method gives the three roots directly,
and gives three only. I have deemed it necessary
to enter upon these slight details for the reason
that if on the one hand the method was long taxed
with being able to give but one root, on the other
hand when it was seen that it really gave three it was
thought that it should have given six, owing to the
\PageSep{72}
false employment of all the possible combinations of
the three cubic roots of unity, viz., $1$,~$m$,~$n$, with the
\index{Unity, three cubic roots of}%
two cubic radicals $\sqrt[3]{A}$~and~$\sqrt[3]{B}$.

We could have arrived directly at the results which
\MNote{A direct method of reaching the roots.}
we have just found by remarking that the two equations
\[
y^{3} + z^{3} + q = 0 \quad\text{and}\quad 3yz + p = 0
\]
give
\[
y^{3} + z^{3} = -q \quad\text{and}\quad y^{3}z^{3} = -\frac{p^{3}}{27};
\]
where it will be seen at once that $y^{3}$~and~$z^{3}$ are the
roots of an equation of the second degree of which
the second term is~$q$ and the third~$-\dfrac{p^{3}}{27}$. This equation,
which is called \emph{the reduced equation}, will accordingly
have the form
\[
u^{2} + qu - \frac{p^{3}}{27} = 0;
\]
and calling $A$~and~$B$ its two roots we shall have immediately
\[
y = \sqrt[3]{A},\quad z = \sqrt[3]{B},
\]
where it will be observed that $A$~and~$B$ have the same
values that they had in the previous discussion. Now,
from what has gone before, we shall likewise have
\[
y = m\sqrt[3]{A} \quad\text{or}\quad y = n\sqrt[3]{A},
\]
and the same will also hold good for~$z$. But the equation
\[
zy = -\frac{p}{3},
\]
of which we have employed the cube only, limits these
\PageSep{73}
values and it is easy to see that the restriction requires
the three corresponding values of~$z$ to be
\[
\sqrt[3]{B},\quad m\sqrt[3]{B},\quad n\sqrt[3]{B};
\]
whence follow for the value of~$x$, which is equal to~$y + z$,
the same three values which we found above.

As to the form of these values it is apparent, first,
that so long as $A$~and~$B$ are real quantities, one only
\MNote{The form of the roots\Add{.}}
of them can be real, for $m$~and~$n$ are imaginary. They
can consequently all three be real only in the case
where the roots $A$~and~$B$ of the reduced equation are
imaginary, that is, when the quantity
\[
\frac{q^{2}}{4} + \frac{p^{3}}{27}
\]
beneath the radical sign is negative, which happens
only when $p$~is negative and greater than
\[
3\sqrt[3]{\frac{q^{2}}{4}}.
\]
And this is the so-called \emph{irreducible case}.
\index{Irreducible case}%

Since in this event
\[
\frac{q^{2}}{4} + \frac{p^{3}}{27}
\]
is a negative quantity, let us suppose it equal to~$-g^{2}$,
$g$~being any real quantity whatever. Then making,
for the sake of simplicity,
\[
-\frac{q}{2} = f,
\]
the two roots $A$~and~$B$ of the reduced equation assume
the form
\[
A = f + g\sqrt{-1},\quad B = f - g\sqrt{-1}.
\]
\PageSep{74}

Now I say that if $\sqrt[3]{A} + \sqrt[3]{B}$, which is one of the
\MNote{The reality of the roots\Add{.}}
\index{Roots!reality@the reality of the}%
roots of the equation of the third degree, is real, then
the two other roots, expressed by
\[
m\sqrt[3]{A} + n\sqrt[3]{B} \quad\text{and}\quad n\sqrt[3]{A} + m\sqrt[3]{B},
\]
will also be real. Put
\[
\sqrt[3]{A} = t,\quad \sqrt[3]{B} = u;
\]
we shall have
\[
t + u = h,
\]
where $h$~by hypothesis is a real quantity. Now,
\[
tu = \sqrt[3]{AB} \quad\text{and}\quad AB = f^{2} + g^{2},
\]
therefore
\[
tu = \sqrt[3]{f^{2} + g^{2}};
\]
squaring the equation $t + u = h$ we have
\[
t^{2} + 2tu + u^{2} = h^{2};
\]
from which subtracting~$4tu$ we obtain
\[
(t - u)^{2} = h^{2} - 4\sqrt[3]{f^{2} + g^{2}}.
\]
I observe that this quantity must necessarily be negative,
for if it were positive and equal to~$k^{2}$ we should
have
\[
(t - u)^{2} = k^{2},
\]
whence
\[
t - u = k.
\]
Then since
\[
t + u = h,
\]
it would follow that
\[
t = \frac{h + k}{2} \quad\text{and}\quad u = \frac{h - k}{2},
\]
\PageSep{75}
both of which are real quantities. But then $t^{3}$~and~$u^{3}$
would also be real quantities, which is contrary to
our hypothesis, since these quantities are equal to $A$~and~$B$,
both of which are imaginary.

The quantity
\[
h^{2} - 4\sqrt[3]{f^{2} + g^{2}}
\]
therefore, is necessarily negative. Let us suppose it
equal to~$-k^{2}$; we shall have then
\[
(t - u)^{2} = -k^{2},
\]
and extracting the square root
\[
t - u = k\sqrt{-1};
\]
\MNote{The form of the two cubic radicals.}
\index{Cubic radicals}%
\index{Radicals, cubic}%
whence
\[
t = \frac{h + k\sqrt{-1}}{2} = \sqrt[3]{A},\quad
u = \frac{h - k\sqrt{-1}}{2} = \sqrt[3]{B}.
\]

Such necessarily will be the form of the two cubic
radicals
\[
\sqrt[3]{f + g\sqrt{-1}} \quad\text{and}\quad \sqrt[3]{f - g\sqrt{-1}},
\]
a form at which we can arrive directly by expanding
these roots according to the Newtonian theorem into
series. But since proofs by series are apt to leave
some doubt in the mind, I have sought to render the
preceding discussion entirely independent of them.

If, therefore,
\[
\sqrt[3]{A} + \sqrt[3]{B} = h,
\]
we shall have
\[
\sqrt[3]{A} = \frac{h + k\sqrt{-1}}{2} \quad\text{and}\quad
\sqrt[3]{B} = \frac{h - k\sqrt{-1}}{2}.
\]
Now we have found above that
\[
m = \frac{-1 + \sqrt{-3}}{2},\quad n = \frac{-1 - \sqrt{-3}}{2};
\]
\PageSep{76}
wherefore, multiplying these quantities together, we
have
\begin{align*}
m\sqrt[3]{A} + n\sqrt[3]{B} &= \frac{-h + k\sqrt{-3}}{2} \\
\intertext{and}
n\sqrt[3]{A} + m\sqrt[3]{B} &= \frac{-h - k\sqrt{-3}}{2},
\end{align*}
which are real quantities. Consequently, if the root~$h$
\MNote{Condition of the reality of the roots.}
\index{Reality of roots}%
\index{Roots!reality@the reality of the}%
is real, the two other roots also will be real in the
irreducible case and they will be real in that case only.

But the invariable difficulty is, to demonstrate directly
that
\[
\sqrt[3]{f + g\sqrt{-1}} + \sqrt[3]{f - g\sqrt{-1}},
\]
which we have supposed equal to~$h$, is always a real
quantity whatever be the values of $f$~and~$g$. In particular
cases the demonstration can be effected by the
extraction of the cube root, when that is possible. For
example, if $f = 2$, $g = 11$, we shall find that the cube
root of~$2 + 11\sqrt{-1}$ will be~$2 + \sqrt{-1}$, and similarly
that the cube root of~$2 - 11\sqrt{-1}$ will be~$2 - \sqrt{-1}$,
and the sum of the radicals will be~$4$. An infinite
number of examples of this class may be constructed
and it was through the consideration of such instances
that Bombelli became convinced of the reality of the
imaginary expression in the formula for the irreducible
case. But forasmuch as the extraction of cube roots
is in general possible only by means of series, we cannot
arrive in this way at a general and direct demonstration
of the proposition under consideration.
\PageSep{77}

It is otherwise with square roots and with all roots
of which the exponents are powers of~$2$. For example,
\MNote{Extraction of the square roots of two imaginary binomials.}
\index{Binomials, extraction of the square roots of two imaginary}%
\index{Imaginary binomials, square roots of}%
if we have the expression
\[
\sqrt{f + g\sqrt{-1}} + \sqrt{f - g\sqrt{-1}},
\]
composed of two imaginary radicals, its square will be
\[
2f + 2\sqrt{f^{2} + g^{2}},
\]
a quantity which is necessarily positive. Extracting
the square root, so as to obtain the equivalent expression,
we have
\[
\sqrt{2f + 2\sqrt{f^{2} + g^{2}}},
\]
for the real value of the imaginary quantity we started
with. But if instead of the sum we had had the difference
between the two proposed imaginary radicals
we should then have obtained for its square the following
expression
\[
2f - 2\sqrt{f^{2} + g^{2}},
\]
a quantity which is necessarily negative; and, taking
the square root of the latter, we should have obtained
the simple imaginary expression
\[
\sqrt{2f - 2\sqrt{f^{2} + g^{2}}}.
\]

Further, if the quantity
\[
\sqrt[4]{f + g\sqrt{-1}} + \sqrt[4]{f - g\sqrt{-1}}
\]
were given, we should have, by squaring, the form
\begin{multline*}
%[** TN: Moved equality sign to second line]
\sqrt{f + g\sqrt{-1}} + \sqrt{f - g\sqrt{-1}} + 2\sqrt[4]{f^{2} + g^{2}} \\
= \sqrt{2f + 2\sqrt{f^{2} + g^{2}}} + 2\sqrt[4]{f^{2} + g^{2}},
\end{multline*}
a real and positive quantity. Extracting the square
\PageSep{78}
root of this expression we should obtain a real value
for the original quantity; and so on for all the other
remaining even roots. But if we should attempt to
apply the preceding method to cubic radicals we
should be led again to equations of the third degree
in the irreducible case.

For example, let
\MNote{Extraction of the cube roots of two imaginary binomials.}
\[
\sqrt[3]{f + g\sqrt{-1}} + \sqrt[3]{f - g\sqrt{-1}} = x.
\]
Cubing, we get
\[
2f + 3\sqrt[3]{f^{2} + g^{2}}\left(
\sqrt[3]{f + g\sqrt{-1}} + \sqrt[3]{f - g\sqrt{-1}}
\right) = x^{3};
\]
that is
\[
2f + 3x\sqrt[3]{f^{2} + g^{2}} = x^{3},
\]
or, with the terms properly arranged,
\[
x^{3} - 3x\sqrt[3]{f^{2} + g^{2}} - 2f = 0,
\]
the general formula of the irreducible case, for
\[
\frac{1}{4}(2f)^{2} + \frac{1}{27}\bigl(-3\sqrt[3]{f^{2} + g^{2}}\bigr)^{3}
  = -g^{2}.
\]
If $g = 0$ we shall have $x = 2\sqrt[3]{f}$. The sole \textit{desideratum},
therefore, is to demonstrate that if $g$~have any value
whatever, $x$~has a corresponding real value. Now the
second last equation gives
\[
\sqrt[3]{f^{2} + g^{2}} = \frac{x^{3} - 2f}{3x}\Add{,}
\]
and cubing we get
\[
f^{2} + g^{2} = \frac{x^{9} - 6x^{6}f + 12x^{3}f^{2} - 8f^{3}}{27x^{3}},
\]
whence
\[
g^{2} = \frac{x^{9} - 6x^{6}f - 15x^{3}f^{2} - 8f^{3}}{27x^{3}},
\]
\PageSep{79}
an equation which may be written as follows
\[
g^{2} = \frac{(x^{3} - 8f)(x^{3} + f)^{2}}{27x^{3}},
\]
or, better, thus:
\[
g^{2} = \frac{1}{27}\left(1 - \frac{8f}{x^{3}}\right)(x^{3} + f)^{2}.
\]

It is plain from the last expression that $g$~is zero
when $x^{3} = 8f$; further, that $g$~constantly and uninterruptedly
\MNote{General theory of the reality of the roots\Add{.}}
\index{Roots!reality@the reality of the}%
increases as $x$~increases; for the factor
$(x^{3} + f)^{2}$ augments constantly, and the other factor
$1 - \dfrac{8f}{x^{3}}$ also keeps increasing, seeing that as the denominator~$x^{3}$
increases the negative part~$\dfrac{8f}{x^{3}}$, which is
originally equal to~$1$, keeps constantly growing less
than~$1$. Therefore, if the value of~$x^{3}$ be increased by
insensible degrees from~$8f$ to infinity, the value of~$g^{2}$
will also augment by insensible and corresponding
degrees from zero to infinity. And therefore, reciprocally,
to every value of~$g^{2}$ from zero to infinity there
must correspond some value of~$x^{3}$ lying between the
limits of~$8f$ and infinity, and since this is so whatever
be the value of~$f$ we may legitimately conclude that,
be the values of $f$~and~$g$ what they may, the corresponding
value of~$x^{3}$ and consequently also of~$x$ is
always real.

But how is this value of~$x$ to be assigned? It would
\index{Imaginary expressions|EtSeq}%
seem that it can be represented only by an imaginary
expression or by a series which is the development of
an imaginary expression. Are we to regard this class
of imaginary expressions, which correspond to real
\PageSep{80}
values, as constituting a new species of algebraical expressions
which although they are not, like other expressions,
\MNote{Imaginary expressions\Add{.}}
susceptible of being numerically evaluated
in the form in which they exist, yet possess the indisputable
advantage---and this is the chief requisite---that
they can be employed in the operations of algebra
exactly as if they did not contain imaginary expressions.
They further enjoy the advantage of having a
wide range of usefulness in geometrical constructions,
as we shall see in the theory of angular sections, so
\index{Angular sections, theory of}%
that they can always be exactly represented by lines;
while as to their numerical value, we can always find
it approximately and to any degree of exactness that
we desire, by the approximate resolution of the equation
on which they depend, or by the use of the common
trigonometrical tables.

It is demonstrated in geometry that if in a circle
having the radius~$r$ an arc be taken of which the chord
is~$c$, and that if the chord of the third part of that arc
be called~$x$, we shall have for the determination of~$x$
the following equation of the third degree
\[
x^{3} - 3r^{2}x + r^{2}c = 0,
\]
an equation which leads to the irreducible case since
$c$~is always necessarily less than~$2r$, and which, owing
to the two undetermined quantities $r$~and~$c$, may be
taken as the type of all equations of this class. For,
if we compare it with the general equation
\[
x^{3} + px + q = 0,
\]
we shall have
\PageSep{81}
\[
r = \sqrt{-\frac{p}{3}} \quad\text{and}\quad c = -\frac{3q}{p}
\]
so that by trisecting the arc corresponding to the
chord~$c$ in a circle of the radius~$r$ we shall obtain at
\MNote{Trisection of an angle.}
\index{Angle, trisection of an}%
\index{Trisection of an angle}%
once the value of a root~$x$, which will be the chord of
the third part of that arc. Now, from the nature of a
circle the same chord~$c$ corresponds not only to the
arc~$s$ but (calling the entire circumference~$u$) also to
the arcs
\[
u - s,\quad 2u + s,\quad 3u - s, \dots\Add{.}
\]
Also the arcs
\[
u + s,\quad 2u - s,\quad 3u + s, \dots
\]
have the same chord, but taken negatively, for on
completing a full circumference the chords become
zero and then negative, and they do not become positive
again until the completion of the second circumference,
as you may readily see. Therefore, the values
of~$x$ are not only the chord of the arc~$\dfrac{s}{3}$ but also
the chords of the arcs
\[
\frac{u - s}{3},\quad \frac{2u + s}{3},
\]
and these chords will be the three roots of the equation
proposed. If we were to take the succeeding arcs
which have the same chord~$c$ we should be led simply
to the same roots, for the arc~$3u - s$ would give the
chord of~$\dfrac{3u - s}{3}$, that is, of~$u - \dfrac{s}{3}$, which we have already
seen is the same as that of~$\dfrac{s}{3}$, and so with the
rest.
\PageSep{82}

Since in the irreducible case the coefficient~$p$ is
\index{Irreducible case}%
necessarily negative, the value of the given chord~$c$
\MNote{Trigonometrical solution.}
will be positive or negative according as $q$~is positive
or negative. In the first case, we take for~$s$ the arc
subtended by the positive chord $c = -\dfrac{3q}{p}$. The second
case is reducible to the first by making $x$~negative,
whereby the sign of the last term is changed; so
that if again we take for~$s$ an arc subtended by the
positive chord~$\dfrac{3q}{p}$, we shall have simply to change
the sign of the three roots.

Although the preceding discussion may be deemed
sufficient to dispel all doubts concerning the nature
of the roots of equations of the third degree, we propose
\index{Equations!third@of the third degree}%
\index{Third degree, equations of the}%
adding to it a few reflexions concerning the
method by which the roots are found. The method
which we have propounded in the foregoing and which
is commonly called \emph{Cardan's method}, although it seems
\index{Cardan}%
to me that we owe it to Hudde, has been frequently
\index{Hudde}%
criticised, and will doubtless always be criticised, for
giving the roots in the irreducible case in an imaginary
form, solely because a supposition is here made which
is contradictory to the nature of the equation. For
the very gist of the method consists in its supposing
\index{Undetermined quantities}%
the unknown quantity equal to two undetermined
quantities $y + z$, in order to enable us afterwards to
separate the resulting equation
\[
y^{3} + z^{3} + (3yz + p)(y + z) + q = 0
\]
into the two following:
\PageSep{83}
\[
3yz + p = 0 \quad\text{and}\quad y^{3} + z^{3} + q = 0.
\]
Now, throwing the first of these into the form
\MNote{The method of indeterminates.}
\index{Indeterminates, the method of}%
\[
y^{3}z^{3} = -\frac{p^{3}}{27}
\]
it is plain that the question reduces itself to finding
two numbers $y^{3}$~and~$z^{3}$ of which the sum is~$-q$ and
the product~$-\dfrac{p^{3}}{27}$, which is impossible unless the
square of half the sum exceed the product, for the
difference between these two quantities is equal to the
square of half the difference of the numbers sought.

The natural conclusion was that it was not at all
astonishing that we should reach imaginary expressions
\index{Imaginary expressions}%
when proceeding from a supposition which it
was impossible to express in numbers, and so some
writers have been induced to believe that by adopting
a different course the expression in question could be
avoided and the roots all obtained in their real form.
\index{Reality of roots}%
\index{Roots!reality@the reality of the}%

Since pretty much the same objection can be advanced
against the other methods which have since
been found and which are all more or less based upon
the method of indeterminates, that is, the introduction
of certain arbitrary quantities to be determined
so as to satisfy the conditions of the problem,---we
propose to consider the question of the reality of the
roots by itself and independently of any supposition
whatever. Let us take again the equation
\[
x^{3} + px + q = 0;
\]
and let us suppose that its three roots are $a$,~$b$,~$c$.
\PageSep{84}

By the theory of equations the left-hand side of
\index{Equations!theory of}%
the preceding expression is the product of three quantities
\MNote{An independent consideration.}
\[
x - a,\quad x - b,\quad x - c,
\]
which, multiplied together, give
\[
x^{3} - (a + b + c)x^{2} + (ab + ac + bc)x - abc;
\]
and comparing the corresponding terms, we have
\[
a + b + c = 0,\quad
ab + ac + bc = p,\quad
abc = -q.
\]
As the degree of the equation is odd we may be certain,
as you doubtless already know and in any event
will clearly see from the lecture which is to follow,
that it has necessarily one real root. Let that root
be~$c$. The first of the three equations which we have
just found will then give
\[
c = -a - b,
\]
whence it is plain that $a + b$ is also necessarily a real
quantity. Substituting the last value of~$c$ in the second
and third equations, we have
\[
ab - a^{2} - ab - ab - b^{2} = p,\quad -ab(a + b) = -q,
\]
or
\[
a^{2} + ab + b^{2} = -p,\quad ab(a + b) = q,
\]
from which are to be found $a$~and~$b$. The last equation
gives $ab = \dfrac{q}{a + b}$ from which I conclude that $ab$
also is necessarily a real quantity. Let us consider
now the quantity $\dfrac{q^{2}}{4} + \dfrac{p^{3}}{27}$ or, clearing of fractions, the
quantity $27q^{2} + 4p^{3}$, upon the sign of which the irreducible
case depends. Substituting in this for $p$~and~$q$
their value as given above in terms of $a$~and~$b$,
\PageSep{85}
we shall find that when the necessary reductions are
made the quantity in question is equal to the square of
\MNote{New view of the reality of the roots.}
\index{Reality of roots}%
\index{Roots!reality@the reality of the}%
\[
2a^{3} - 2b^{3} + 3a^{2}b - 3ab^{2}
\]
taken negatively; so that by changing the signs and
extracting the square root we shall have
\[
2a^{3} - 2b^{3} + 3a^{2}b - 3ab^{2} = \sqrt{-27q^{2} - 4p^{3}},
\]
whence it is easy to infer that the two roots $a$~and~$b$
cannot be real unless the quantity $27q^{2} + 4p^{3}$ be negative.
But I shall show that in that case, which is as
we know the irreducible case, the two roots $a$~and~$b$
are necessarily real. The quantity
\[
2a^{3} - 2b^{3} + 3a^{2}b - 3ab^{2}
\]
may be reduced to the form
\[
(a - b)(2a^{2} + 2b^{2} + 5ab),
\]
as multiplication will show. Now, we have already
seen that the two quantities $a + b$ and $ab$ are necessarily
real, whence it follows that
\[
2a^{2} + 2b^{2} + 5ab = 2(a + b)^2 + ab
\]
is also necessarily real. Hence the other factor~$a - b$
is also real when the radical $\sqrt{-27q^{2} - 4p^{3}}$ is real.
Therefore $a + b$ and $a - b$ being real quantities, it follows
that $a$~and~$b$ are real.

We have already derived the preceding theorems
from the form of the roots themselves. But the present
demonstration is in some respects more general
and more direct, being deduced from the fundamental
principles of the problem itself. We have made no
\PageSep{86}
suppositions, and the particular nature of the irreducible
case has introduced no imaginary quantities.

\MNote{Final solution on the new view.}
But the values of $a$~and~$b$ still remain to be found
from the preceding equations. And to this end I observe
that the left-hand side of the equation
\[
a^{3} - b^{3} + \frac{3}{2}(a^{2}b - ab^{2})
  = \frac{1}{2}\sqrt{-27q^{2} - 4p^{3}}
\]
can be made a perfect cube by adding the left-hand
side of the equation
\[
ab(a + b) = q,
\]
multiplied by $\dfrac{3\sqrt{-3}}{2}$, and that the root of this cube is
\[
\frac{1 - \sqrt{-3}}{2}b - \frac{1 + \sqrt{-3}}{2}a
\]
so that, extracting the cube root of both sides, we
shall have the expression
\[
\frac{1 - \sqrt{-3}}{2}b - \frac{1 + \sqrt{-3}}{2}a
\]
expressed in known quantities. And since the radical
$\sqrt{-3}$ may also be taken negatively, we shall also
have the expression
\[
\frac{1 + \sqrt{-3}}{2}b - \frac{1 - \sqrt{-3}}{2}a
\]
expressed in known quantities, from which the values
of $a$~and~$b$ can be deduced. And these values will
contain the imaginary quantity~$\sqrt{-3}$, which was introduced
by multiplication, and will be reducible to
the same form with the two roots
\PageSep{87}
\[
m\sqrt[3]{A} + n\sqrt[3]{B} \quad\text{and}\quad n\sqrt[3]{A} + m\sqrt[3]{B},
\]
which we found above. The third root
\MNote{Office of imaginary quantities.}
\[
c = -a - b
\]
will then be expressed by $\sqrt[3]{A} + \sqrt[3]{B}$.

By this method we see that the imaginary quantities
\index{Imaginary quantities, office of the}%
employed have simply served to facilitate the extraction
of the cube root without which we could not
determine separately the values of $a$~and~$b$. And since
it is apparently impossible to attain this object by a
different method, we may regard it as a demonstrated
truth that the general expression of the roots of an
equation of the third degree in the irreducible case
cannot be rendered independent of imaginary quantities.

Let us now pass to \emph{equations of the fourth degree}.
\index{Equations!fourth@of the fourth degree}%
We have already said that the artifice which was originally
employed for resolving these equations consisted
in so arranging them that the square root of
the two sides could be extracted, by which they were
reduced to equations of the second degree. The following
is the procedure employed. Let
\[
x^{4} + px^{2} + qx + r = 0
\]
be the general equation of the fourth degree deprived
of its second term, which can always be eliminated,
as you know, by increasing or diminishing the roots
by a suitable quantity. Let the equation be put in
the form
\[
x^{4} = -px^{2} - qx - r,
\]
\PageSep{88}
and to each side let there be added the terms $2x^{2}y + y^{2}$,
which contain a new undetermined quantity~$y$ but
\MNote{Biquadratic equations.}
\index{Biquadratic equations}%
\index{Equations!biquadratic}%
which still leave the left-hand side of the equation a
square. We shall then have
\[
(x^{2} + y)^{2} = (2y - p)x^{2} - qx + y^{2} - r.
\]
We must now make the right-hand side also a square.
To this end it is necessary that
\[
4(2y - p)(y^{2} - r) = q^{2},
\]
in which case the square root of the right-hand side
will have the form
\[
x\sqrt{2y - p} - \frac{q}{2\sqrt{2y - p}}.
\]
Supposing then that the quantity~$y$ satisfies the equation
\[
4(2y - p)(y^{2} - r) = q^{2},
\]
which developed becomes
\[
y^{3} - \frac{py^{2}}{2} - ry + \frac{pr}{2} - \frac{q^{2}}{8} = 0,
\]
and which, as we see, is an equation of the third degree,
the equation originally given may be reduced to
the following by extracting the square root of its two
members,~viz.:
\[
x^{2} + y = x\sqrt{2y - p} - \frac{q}{2\sqrt{2y - p}},
\]
where we may take either the plus or the positive
value for the radical $\sqrt{2y - p}$, and shall consequently
have two equations of the second degree to which the
given equation has been reduced and the roots of
which will give the four roots of the original equation.
\PageSep{89}
All of which furnishes us with our first instance of the
decomposition of equations into others of lower degree.

The method of Descartes which is commonly followed
\index{Descartes}%
in the elements of algebra is based upon the
\MNote{The method of Descartes.}
same principle and consists in assuming at the outset
that the proposed equation is produced by the multiplication
of two equations of the second degree, as
\[
x^{2} - ux + s = 0 \quad\text{and}\quad x^{2} + ux + t = 0,
\]
where $u$,~$s$, and~$t$ are indeterminate coefficients. Multiplying
\index{Coefficients!indeterminate}%
\index{Indeterminate coefficients}%
them together we have
\[
x^{4} + (s + t - u^{2})x + (s - t)ux + st = 0,
\]
comparison of which with the original equation gives
\[
s + t - u^{2} = p,\quad (s - t)u = q \quad\text{and}\quad st = r.
\]
The first two equations give
\[
2s = p + u^{2} + \frac{q}{u},\quad 2t = p + u^{2} - \frac{q}{u}.
\]
And if these values be substituted in the third equation
of condition $st = r$, we shall have an equation of
the sixth degree in~$u$, which owing to its containing
only even powers of~$u$ is resolvable by the rules for
cubic equations. And if we substitute in this equation
$2y - p$ for~$u^{2}$, we shall obtain in~$y$ the same reduced
equation that we found above by the old method.

Having the value of~$u^{2}$ we have also the values of
$s$~and~$t$, and our equation of the fourth degree will be
decomposed into two equations of the second degree
which will give the four roots sought. This method,
as well as the preceding, has been the occasion of some
\PageSep{90}
hesitancy as to which of the three roots of the reduced
cubic equation in $u^{2}$ or~$y$ should be employed.
\MNote{The determined character of the roots\Add{.}}
The difficulty has been well resolved in Clairaut's
\index{Clairaut}%
\textit{Algebra}, where we are led to see directly that we always
obtain the same four roots or values of~$x$ whatever
root of the reduced equation we employ. But
this generality is needless and prejudicial to the simplicity
which is to be desired in the expression of
the roots of the proposed equation, and we should
prefer the formulæ which you have learned in the
principal course and in which the three roots of the
reduced equation are contained in exactly the same
manner.

The following is another method of reaching the
same formulæ, less direct than that which has already
been expounded to you, but which, on the other hand
has the advantage of being analogous to the method
of Cardan for equations of the third degree.
\index{Cardan}%

I take up again the equation
\[
x^{4} + px^{2} + qx + r = 0,
\]
and I suppose
\[
x = y + z + t.
\]
Squaring I obtain
\[
x^{2} = y^{2} + z^{2} + t^{2} + 2(yz + yt + zt).
\]
Squaring again I have
\[
%[** TN: Set on two lines in original]
x^{4} = (y^{2} + z^{2} + t^{2})^{2} + 4(y^{2} + z^{2} + t^{2})(yz + yt + zt)
+ 4(yz + yt + zt)^{2};
\]
but
\begin{align*}
%[** TN: Re-broken]
(yz + yt + zt)^{2}
  &= y^{2}z^{2} + y^{2}t^{2} + z^{2}t^{2}
   + 2y^{2}zt + 2yz^{2}t + 2yzt^{2} \\
  &= y^{2}z^{2} + y^{2}t^{2} + z^{2}t^{2} + 2yzt(y + z + t).
\end{align*}
\PageSep{91}
Substituting these three values of $x$,~$x^{2}$, and~$x^{4}$ in the
original equation, and bringing together the terms
multiplied by~$y + z + t$ and the terms multiplied by~$yz + yt + zt$,
\MNote{A third method.}
I have the transformed equation
\begin{gather*}
%[** TN: Re-broken]
(y^{2} + z^{2} + t^{2})^{2} + p(y^{2} + z^{2} + t^{2}) \\
  + \bigl[4(y^{2} + z^{2} + t^{2}) + 2p\bigr](yz + yt + zt) \\
  + 4(y^{2}z^{2} + y^{2}t^{2} + z^{2}t^{2})
  + (8yzt + q)(y + z + t) + r = 0.
\end{gather*}
We now proceed as we did with equations of the third
degree, where we caused the terms containing $y + z$
to vanish, and in the same manner cause here the
terms containing $y + z + t$ and $yz + yt + zt$ to disappear,
which will give us the two equations of condition
\[
8yzt + q = 0 \quad\text{and}\quad 4(y^{2} + z^{2} + t^{2}) + 2p = 0.
\]

There remains the equation
\[
(y^{2} + z^{2} + t^{2})^{2} + p(y^{2} + z^{2} + t^{2})
  + 4(y^{2}z^{2} + y^{2}t^{2} + z^{2}t^{2}) + r = 0;
\]
and the three together will determine the quantities
$y$,~$z$, and~$t$. The second gives immediately
\[
y^{2} + z^{2} + t^{2} = -\frac{p}{2},
\]
which substituted in the third gives
\[
y^{2}z^{2} + y^{2}t^{2} + z^{2}t^{2} = \frac{p^{2}}{16} - \frac{r}{4}\Add{.}
\]
The first, raised to its square, gives
\[
y^{2}z^{2}t^{2} = \frac{q^{2}}{64}.
\]
Hence, by the general theory of equations the three
\PageSep{92}
quantities $y^{2}$,~$z^{2}$,~$t^{2}$ will be the roots of an equation of
the third degree having the form
\MNote{The reduced equation.}
\[
u^{3} + \frac{p}{2} u^{2}
  + \left(\frac{p^{2}}{16} - \frac{r}{4}\right)u
  - \frac{q^{2}}{64} = 0;
\]
so that if the three roots of this equation, which we
will call \emph{the reduced equation}, be designated by $a$,~$b$,~$c$,
we shall have
\[
y = \sqrta,\quad z = \sqrt{b},\quad t = \sqrtc,
\]
and the value of~$x$ will be expressed by
\[
\sqrta + \sqrt{b} + \sqrtc.
\]
Since the three radicals may each be taken with the
plus sign or the minus sign, we should have, if all
possible combinations were taken, eight different values
for~$x$. It is to be observed, however, that in the
preceding analysis we employed the equation $y^{2}z^{2}t^{2} = \dfrac{q^{2}}{64}$,
whereas the equation immediately given is $yzt = -\dfrac{q}{8}$.
Hence the product of the three quantities $y$,~$z$,~$t$,
that is to say of the three radicals
\[
\sqrta,\quad \sqrt{b}, \quad \sqrtc,
\]
must have the contrary sign to that of the quantity~$q$.
Therefore, if $q$~be a negative quantity, either three
positive radicals or one positive and two negative radicals
must be contained in the expression for~$x$. And
in this case we shall have the following four combinations
only:
\begin{alignat*}{2}
 &\sqrta + \sqrt{b} + \sqrtc,\qquad    && \sqrta - \sqrt{b} - \sqrtc,\\
-&\sqrta + \sqrt{b} - \sqrtc, &\Typo{}{-}&\sqrta - \sqrt{b} + \sqrtc,
\end{alignat*}
\PageSep{93}
which will be the four roots of the proposed equation
of the fourth degree. But if $q$~be a positive quantity,
either three negative radicals or one negative and two
\MNote{Euler's formulæ.}
positive radicals must be contained in the expression
for~$x$, which will give the following four other combinations
as the roots of the proposed equation:\footnote
  {These simple and elegant formulæ are due to Euler. But M.~Bret, Professor
  \index{Bret, M.|FN}%
  \index{Euler}%
  of Mathematics at Grenoble, has made the important observation (see
  the \textit{Correspondance sur l'\Typo{Ecole}{École} Polytechnique}, t.~II., 3\ieme~Cahier, p.~217) that
  they can give false values when imaginary quantities occur among the four
  roots.

  In order to remove all difficulty and ambiguity we have only to substitute
  for one of these radicals its value as derived from the equation $\sqrta\sqrt{b}\sqrtc = -\dfrac{q}{8}$.
  Then the formula
  \[
  \sqrta + \sqrt{b} - \frac{q}{8\sqrta\sqrt{b}}
  \]
  will give the four roots of the original equation by taking for $a$~and~$b$ any two
  of the three roots of the reduced equation, and by taking the two radicals
  successively positive and negative.

  The preceding remark should be added to article~777 of Euler's \textit{Algebra}
  and to article~37 of the author's Note~XIII of the \textit{Traité de la résolution des
  équations numériques}.}
\begin{alignat*}{2}
-&\sqrta - \sqrt{b} - \sqrtc,\qquad & -&\sqrta + \sqrt{b} + \sqrtc, \displaybreak[1] \\
 &\sqrta - \sqrt{b} + \sqrtc,         &&\sqrta + \sqrt{b} - \sqrtc.
\end{alignat*}

Now if the three roots $a$,~$b$,~$c$ of the reduced equation
\index{Reality of roots}%
\index{Roots!reality@the reality of the}%
\index{Three roots, reality of the}%
of the third degree are all real and positive, it is
evident that the four preceding roots will also all be
real. But if among the three real roots $a$,~$b$,~$c$, any
are negative, obviously the four roots of the given
biquadratic equation will be imaginary. Hence, besides
the condition for the reality of the three roots of
the reduced equation it is also requisite in the first
case, agreeably to the well-known rule of Descartes,
\index{Descartes}%
\PageSep{94}
that the coefficients of the terms of the reduced equation
should be alternatively positive and negative, and
\MNote{Roots of a biquadratic equation.}
\index{Biquadratic equations}%
\index{Roots!biquadratic@of a biquadratic equation}%
consequently that $p$~should be negative and $\dfrac{p^{2}}{16} - \dfrac{r}{4}$
positive, that is, $p^{2} > 4r$. If one of these conditions
is not realised the proposed biquadratic equation cannot
have four real roots. If the reduced equation have
but one real root, it will be observed, first, that by
reason of its last term being negative the one real root
of the equation must necessarily be positive. It is
then easy to see from the general expressions which
we gave for the roots of cubic equations deprived of
their second term,---a form to which the reduced equation
in~$u$ can easily be brought by simply increasing
all the roots by the quantity~$\dfrac{p}{6}$,---it is easy to see, I
say, that the two imaginary roots of this equation will
be of the form
\[
f + g\sqrt{-1} \quad\text{and}\quad f - g\sqrt{-1}.
\]
Therefore, supposing $a$~to be the real root and $b$,~$c$ the
two imaginary roots, $\sqrta$~will be a real quantity and
$\sqrt{b} + \sqrtc$ will also be real for reasons which we have
given above; while $\sqrt{b} - \sqrtc$ on the other hand will
be imaginary. Whence it follows that of the four
roots of the proposed biquadratic equation, the two
first will be real and the two others will be imaginary.

As for the rest, if we make $u = s - \dfrac{p}{6}$ in the reduced
equation in~$u$, so as to eliminate the second
term and to reduce it to the form which we have above
\PageSep{95}
examined, we shall have the following transformed
equation in~$s$:
\[
s^{3} - \left(\frac{p^{2}}{48} + \frac{r}{4}\right)s
  - \frac{p^{3}}{864} + \frac{pr}{24} - \frac{q^{2}}{64} = 0;
\]
and the condition for the reality of the three roots of
the reduced equation will be
\[
4\left(\frac{p^{2}}{48} + \frac{r}{4}\right)^{3}
  > 27\left(\frac{p^{3}}{864} - \frac{pr}{24} + \frac{q^{2}}{64}\right)^{2}.
\]
\PageSep{96}%XXXX


\Lecture{IV.}{On the Resolution of Numerical Equations.}
\index{Numerical equations!resolution of|(}%

\First{We} have seen how equations of the second, the
third, and the fourth degree can be resolved.
\MNote{Limits of the algebraical resolution of equations.}
\index{Algebraical resolution of equations!limits of the}%
\index{Equations!limits of the algebraical resolution of}%
The fifth degree constitutes a sort of barrier to analysts,
\index{Equations!fifth@of the fifth degree}%
\index{Fifth degree, equations of the}%
which by their greatest efforts they have never
yet been able to surmount, and the general resolution
of equations is one of the things that are still to be
desired in algebra. I say in algebra, for if with the
third degree the analytical expression of the roots is
insufficient for determining in all cases their numerical
value, \textit{a~fortiori} must it be so with equations of a
higher degree; and so we find ourselves constantly
under the necessity of having recourse to other means
for determining numerically the roots of a given equation,---for
to determine these roots is in the last resort
the object of the solution of all problems which
necessity or curiosity may offer.

I propose here to set forth the principal artifices
which have been devised for accomplishing this important
object. Let us consider any equation of the
\index{Equations!mth@of the $m$th degree}%
$m$th~degree, represented by the formula
\PageSep{97}
\[
x^{m} + px^{m-1} + qx^{m-2} + rx^{m-3} + \dots + u = 0,
\]
in which $x$~is the unknown quantity, $p$,~$q$,~$r$,~$\dots$ the
known positive or negative coefficients, and $u$~the
\MNote{Conditions of the resolution of numerical equations.}
\index{Numerical equations!conditions of the resolution of}%
last term, not containing~$x$ and consequently also a
known quantity. It is assumed that the values of
these coefficients are given either in numbers or in
lines; (it is indifferent which, seeing that by taking a
given line as the unit or common measure of the rest
we can assign to all the lines numerical values;) and it
is clear that this assumption is always permissible
when the equation is the result of a real and determinate
problem. The problem set us is to find the value,
or, if there be several, the values, of~$x$ which satisfy the
equation, i.e.\Add{,} which render the sum of all its terms
zero. Now any other value which may be given to~$x$
will render that sum equal to some positive or negative
quantity, for since only integral powers of~$x$ enter
the equation, it is plain that every real value of~$x$
will also give a real value for the quantity in question.
The more that value approaches to zero, the more
will the value of~$x$ which has produced it approach to
a root of the equation. And if we find two values of~$x$,
of which one renders the sum of the terms equal to
a positive quantity and the other to a negative quantity,
we may be assured in advance that between these
two values there will of necessity be at least one value
which will render the expression zero and will consequently
be a root of the equation.

Let $P$~stand for the sum of all the terms of the
\PageSep{98}
equation having the sign~$+$ and $Q$~for the sum of all
the terms having the sign~$-$; then the equation will
be represented by
\[
P - Q = 0.
\]
Let us suppose, for further simplicity, that the two
\MNote{Position of the roots of numerical equations.}
\index{Numerical equations!position of the roots of}%
values of~$x$ in question are positive, that $A$~is the
smaller, $B$~the greater, and that the substitution of~$A$
for~$x$ gives a negative result and the substitution of~$B$
for~$x$ a positive result; i.e., that the value of~$P - Q$
is negative when $x = A$, and positive when $x = B$.

Consequently, when $x = A$, $P$~will be less than~$Q$,
and when $x = B$, $P$~will be greater than~$Q$. Now,
from the very form of the quantities $P$~and~$Q$, which
contain only positive terms and whole positive powers
of~$x$, it is clear that these quantities augment continuously
as $x$~augments, and that by making $x$ augment by
insensible degrees through all values from $A$~to~$B$, they
also will augment by insensible degrees but in such
wise that $P$~will increase more than~$Q$, seeing that
from having been smaller than~$Q$ it will have become
greater. Therefore, there must of necessity be some
expression for the value of~$x$ between $A$~and~$B$ which
will make $P = Q$; just as two moving bodies which
\index{Moving bodies, two}%
we suppose to be travelling along the same straight
line and which having started simultaneously from
two different points arrive simultaneously at two other
points but in such wise that the body which was at first
in the rear is now in advance of the other,---just as
two such bodies, I say, must necessarily meet at some
\PageSep{99}
point in their path. That value of~$x$, therefore, which
will make $P = Q$ will be one of the roots of the equation,
and such a value will lie of necessity between $A$~and~$B$.

The same reasoning may be employed for the
\MNote{Position of the roots of numerical equations.}
other cases, and always with the same result.

The proposition in question is also demonstrable
by a direct consideration of the equation itself, which
may be regarded as made up of the product of the
factors,
\[
x - a,\quad x - b,\quad x - c,\dots,
\]
where $a$,~$b$,~$c$,~$\dots$ are the roots. For it is obvious
that this product cannot, by the substitution of two
different values for~$x$, be made to change its sign, unless
at least one of the factors changes its sign. And
it is likewise easy to see that if more than one of the
factors changes its sign, their number must be odd.
Thus, if $A$~and~$B$ are two values of~$x$ for which the
factor $x - b$, for example, has opposite signs, then if
$A$~be larger than~$b$, necessarily $B$~must be smaller
than~$b$, or \textit{vice versa}. Perforce, then, the root~$b$ will
fall between the two quantities $A$~and~$B$.

As for imaginary roots, if there be any in the equation,
\index{Imaginary roots, occur in pairs}%
since it has been demonstrated that they always
occur in pairs and are of the form
\[
f + g\sqrt{-1},\quad f - g\sqrt{-1},
\]
therefore if $a$~and~$b$ are imaginary, the product of the
factors $x - a$ and $x - b$ will be
\PageSep{100}
\[
(x - f - g\sqrt{-1})(x - f + g\sqrt{-1}) = (x - f)^{2} + g^{2},
\]
a quantity which is always positive whatever value be
given to~$x$. From this it follows that alterations in
the sign can be due only to real roots. But since the
theorem respecting the form of imaginary roots cannot
be rigorously demonstrated without employing the
other theorem that every equation of an odd degree
has necessarily one real root, a theorem of which the
general demonstration itself depends on the proposition
which we are concerned in proving, it follows
that that demonstration must be regarded as a sort of
vicious circle, and that it must be replaced by another
which is unassailable.

But there is a more general and simpler method
\MNote{Application of geometry to algebra.}
\index{Algebra!application of geometry to|EtSeq}%
\index{Geometry!application of to algebra|EtSeq}%
of considering equations, which enjoys the advantage
\index{Equations!constructions for solving|EtSeq}%
of affording direct demonstration to the eye of the
principal properties of equations. It is founded upon
a species of application of geometry to algebra which
is the more deserving of exposition as it finds extended
employment in all branches of mathematics.

Let us take up again the general equation proposed
above and let us represent by straight lines all
the successive values which are given to the unknown
quantity~$x$ and let us do the same for the corresponding
values which the left-hand side of the equation
assumes in this manner. To this end, instead of supposing
the right-hand side of the equation equal to
zero, we suppose it equal to an undetermined quantity~$y$.
We lay off the values of~$x$ upon an indefinite
\PageSep{101}
straight line~$AB$ (Fig.~1), starting from a fixed point~$O$
at which $x$~is zero and taking the positive values of~$x$
in the direction~$OB$ to the right of~$O$ and the negative
values of~$x$ in the opposite direction to the left of~$O$.
Then let~$OP$ be any value of~$x$. To represent
the corresponding value of~$y$ we erect at~$P$ a perpendicular
to the line~$OB$ and lay off on it the value of~$y$
in the direction~$PQ$ above the straight line~$OB$ if it is
positive, and on the same perpendicular below~$OB$ if
it is negative. We do the same for all the values of~$x$,
\MNote{Representation of equations by curves.}
\index{Curves!representation of equations by|EtSeq}%
\Figure{1}{0.8\textwidth}
positive as well as negative; that is, we lay off
corresponding values of~$y$ upon perpendiculars to the
straight line through all the points whose distance
from the point~$O$ is equal to~$x$. The extremities of all
these perpendiculars will together form a straight line
or a curve, which will furnish, so to speak, a picture
of the equation
\[
x^{m} + px^{m-1} + qx^{m-2} + \dots + u = y.
\]
The line~$AB$ is called the axis of the curve, $O$~the origin
of the abscissæ, $OP = x$ an abscissa, $PQ = y$ the corresponding
\PageSep{102}
ordinate, and the equations in $x$~and~$y$ the
\index{Equations!general remarks upon the roots of|EtSeq}%
equations of the curve. A curve such as that of Fig.~1
having been described in the manner indicated, it is
clear that its intersections with the axis~$AB$ will give
the roots of the proposed equation
\MNote{Graphic resolution of equations.}
\index{Equations!graphic resolution of}%
\index{Intersections, with the axis give roots|EtSeq}%
\[
x^{m} + px^{m-1} + qx^{m-2} + \dots + u = 0.
\]
For seeing that this equation is realised only when in
the equation of the curve $y$~becomes zero, therefore
those values of~$x$ which satisfy the equation in question
and which are its roots can only be the abscissæ
\ifthenelse{\not\boolean{ForPrinting}}{%
\Figure{1}{0.8\textwidth} %[** TN: [sic], figure repeated]
}{}% [Discard second copy if formatting for printing]
that correspond to the points at which the ordinates
are zero, that is, to the points at which the curve cuts
the axis~$AB$. Thus, supposing the curve of the equation
in $x$~and~$y$ is that represented in Fig.~1, the roots
of the proposed equation will be
\[
OM,\quad ON,\quad OR,\dots \quad\text{and}\quad -OI,\quad -OG,\dots.
\]
I give the sign~$-$ to the latter because the intersections
$I$,~$G$,~$\dots$ fall on the other side of the point~$O$.
The consideration of the curve in question gives rise
to the following general remarks upon equations:
\PageSep{103}

(1) Since the equation of the curve contains only
whole and positive powers of the unknown quantity~$x$
it is clear that to every value of~$x$ there must correspond
\MNote{The consequences of the graphic resolution.}
a determinate value of~$y$, and that the value in
question will be unique and finite so long as $x$~is finite.
But since there is nothing to limit the values of~$x$ they
may be supposed infinitely great, positive as well as
negative, and to them will correspond also values of~$y$
which are infinitely great. Whence it follows that
the curve will have a continuous and single course,
and that it may be extended to infinity on both sides
of the origin~$O$.

(2) It also follows that the curve cannot pass from
one side of the axis to the other without cutting it,
and that it cannot return to the same side without
having cut it twice. Consequently, between any two
points of the curve on the same side of the axis there
will necessarily be either no intersections or an even
number of intersections; for example, between the
points $H$~and~$Q$ we find two intersections $I$~and~$M$,
and between the points $H$~and~$S$ we find four, $I$, $M$
$N$, $R$, and so on. Contrariwise, between a point on
one side of the axis and a point on the other side, the
curve will have an odd number of intersections; for
example, between the points $L$~and~$Q$ there is one intersection~$M$,
and between the points $H$~and~$K$ there
are three intersections, $I$, $M$, $N$, and so on.

For the same reason there can be no simple intersection
unless on both sides of the point of intersection,
\PageSep{104}
above and below the axis, points of the curve are
situated as are the points $L$,~$Q$ with respect to the intersection~$M$.
\MNote{Intersections indicate the roots.}
But two intersections, such as $N$~and~$R$,
may approach each other so as ultimately to coincide
at~$T$. Then the branch~$QKS$ will take the form
of the dotted line~$QTS$ and touch the axis at~$T$, and
will consequently lie in its whole extent above the
axis; this is the case in which the two roots $ON$,~$OR$
are equal. If three intersections coincide at a point,---a
coincidence which occurs when there are three
equal roots,---then the curve will cut the axis in one
additional point only, as in the case of a single point
of intersection, and so on.

Consequently, if we have found for~$y$ two values
having the same sign, we may be assured that between
the two corresponding values of~$x$ there can fall only
an even number of roots of the proposed equation;
that is, that there will be none or there will be two, or
there will be four, etc. On the other hand, if we have
found for~$y$ two values having contrary signs, we may
be assured that between the corresponding values of~$x$
there will necessarily fall an odd number of roots of
the proposed equation; that is, there will be one, or
there will be three, or there will be five, etc.; so that,
in the case last mentioned, we may infer immediately
that there will be at least one root of the proposed
equation between the two values of~$x$.

Conversely, every value of~$x$ which is a root of the
equation will be found between some larger and some
\PageSep{105}
smaller value of~$x$ which on being substituted for~$x$ in
the equation will yield values of~$y$ with contrary signs.

This will not be the case, however, if the value of~$x$
is a double root; that is, if the equation contains
\MNote{Case of multiple roots.}
\index{Multiple roots}%
\index{Roots!multiple}%
two roots of the same value. On the other hand, if
the value of~$x$ is a triple root, there will again exist
a larger and a smaller value for~$x$ which will give to
the corresponding values of~$y$ contrary signs, and so
on with the rest.

If, now, we consider the equation of the curve, it
is plain in the first place, that by making $x = 0$ we
shall have $y = u$; and consequently that the sign of
the ordinate~$y$ will be the same as that of the quantity~$u$,
the last term of the proposed equation. It is also
easy to see that there can be given to~$x$ a positive or
negative value sufficiently great to make the first term~$x^{m}$
of the equation exceed the sum of all the other
terms which have the opposite sign to~$x^{m}$; with the
result that the corresponding value of~$y$ will have the
same sign as the first term~$x^{m}$. Now, if $m$~is odd $x^{m}$~will
be positive or negative according as $x$~is positive
or negative, and if $m$~is even, $x^{m}$~will always be positive
whether $x$~be positive or not.

Whence we may conclude:

(1) That every equation of an odd degree of which
\index{Equations!odd@of an odd degree, roots of}%
the last term is negative has an odd number of roots
between $x = 0$ and some very large positive value of~$x$,
and an even number of roots between $x = 0$ and
some very large negative value of~$x$, and consequently
\PageSep{106}
that it has at least one real positive root. That, contrariwise,
if the last term of the equation is positive it
\MNote{General conclusions as to the character of the roots.}
will have an odd number of roots between $x = 0$ and
some very large negative value of~$x$, and an even
number of roots between $x = 0$ and some very large
positive value of~$x$, and consequently that it will have
at least one real negative root.

(2) That every equation of an even degree, of
\index{Equations!even@of an even degree, roots of}%
which the last term is negative, has an odd number of
roots between $x = 0$ and some very large positive value
of~$x$, as well as an odd number of roots between $x = 0$
and some very large negative value of~$x$, and consequently
that it has at least one real positive root and
one real negative root. That, on the other hand, if
the last term is positive there will be an even number
of roots between $x = 0$ and some very large positive
value of~$x$, and also an even number of roots between
$x = 0$ and some very large negative value of~$x$; with
the result that in this case the equation may have no
real root, whether positive or negative.

We have said that there could always be given to~$x$
a value sufficiently great to make the first term~$x^{m}$ of
the equation exceed the sum of all the terms of contrary
sign. Although this proposition is not in need
of demonstration, seeing that, since the power~$x^{m}$ is
higher than any of the other powers of~$x$ which enter
the equation, it is bound, as $x$~increases, to increase
much more rapidly than these other powers; nevertheless,
in order to leave no doubts in the mind, we
\PageSep{107}
shall offer a very simple demonstration of it,---a demonstration
which will enjoy the collateral advantage
of furnishing a limit beyond which we may be certain
no root of the equation can be found.

To this end, let us first suppose that $x$~is positive,
\index{Limits of roots|(}%
and that $k$~is the greatest of the coefficients of the
\index{Coefficients!greatest negative|EtSeq}%
\MNote{Limits of the real roots of equations.}
\index{Equations!real roots of, limits of the|EtSeq}%
negative terms. If we make $x = k + 1$ we shall have
\[
x^{m} = (k + 1)^{m} = k(k + 1)^{m-1} + (k + 1)^{m-1}.
\]
Similarly,
\begin{align*}
(k + 1)^{m-1} &= k(k + 1)^{m-2} + (k + 1)^{m-2}, \\
(k + 1)^{m-2} &= k(k + 1)^{m-3} + (k + 1)^{m-3}
\end{align*}
and so on; so that we shall finally have
\[
(k + 1)^{m}
  = k(k + 1)^{m-1}
  + k(k + 1)^{m-2}
  + k(k + 1)^{m-3} + \dots + k + 1.
\]
Now this quantity is evidently greater than the sum
of all the negative terms of the equation taken positively,
on the supposition that $x = k + 1$. Therefore,
the supposition $x = k + 1$ necessarily renders the first
term~$x^{m}$ greater than the sum of all the negative terms.
Consequently, the value of~$y$ will have the same sign
as~$x$.

The same reasoning and the same result hold good
when $x$~is negative. We have here merely to change~$x$
into~$-x$ in the proposed equation, in order to change
the positive roots into negative roots, and \textit{vice versa}.

In the same way it may be proved that if any value
be given to~$x$ greater than~$k + 1$, the value of~$y$ will
still have the same sign. From this and from what
has been developed above, it follows immediately that
\PageSep{108}
the equation can have no root equal to or greater than~$k + 1$.

Therefore, in general, if $k$~is the greatest of the
\MNote{Limits of the positive and negative roots.}
coefficients of the negative terms of an equation, and
changing the unknown quantity~$x$ into~$-x$, $h$~is
the greatest of the coefficients of the negative terms
of the new equation,---the first term always being supposed
positive,---then all the real roots of the equation
will necessarily be comprised between the limits
\[
k + 1 \quad\text{and}\quad -h - 1.
\]

But if there are several positive terms in the equation
preceding the first negative term, we may take
for~$k$ a quantity less than the greatest negative coefficient.
In fact it is easy to see that the formula given
above can be put into the form
\[
(k + 1)^{m}
  = k(k + 1)(k + 1)^{m-2}
  + k(k + 1)(k + 1)^{m-3} + \dots + (k + 1)^{2}
\]
and similarly into the following
\[
(k + 1)^{m}
  = k(k + 1)^{2}(k + 1)^{m-3}
  + k(k + 1)^{2}(k + 1)^{m-4} + \dots + (k + 1)^{3}
\]
and so on.

Whence it is easy to infer that if $m - n$ is the exponent
of the first negative term of the proposed equation
of the $m$th~degree, and if $l$~is the largest coefficient
of the negative terms, it will be sufficient if $k$~is
so determined that
\[
k(k + 1)^{n-1} = l.
\]
And since we may take for~$k$ any larger value that we
please, it will be sufficient to take
\PageSep{109}
\[
k^{n} = l,\quad\text{or}\quad k = \sqrt[n]{l}.
\]
And the same will hold good for the quantity~$h$ as the
limit of the negative roots.
\index{Positive roots, superior and inferior limits of the}%
\index{Roots!superior and inferior limits of the positive}%

If, now, the unknown quantity~$x$ be changed into~$\dfrac{1}{z}$,
the largest roots of the equation in~$x$ will be converted
\MNote{Superior and inferior limits of the positive roots.}
into the smallest in the new equation in~$z$, and
conversely. Having effected this transformation, and
having so arranged the terms according to the powers
of~$z$ that the first term of the equation is~$z^{m}$, we may
then in the same manner seek for the limits $K + 1$ and
$-H - 1$ of the positive and negative roots of the
equation in~$z$.

Thus $K + 1$ being larger than the largest value of~$z$
or of~$\dfrac{1}{x}$, therefore, by the nature of fractions, $\dfrac{1}{K + 1}$
will be smaller than the smallest value of~$x$ and similarly
$\dfrac{1}{H + 1}$ will be smaller than the smallest negative
value of~$x$.

Whence it may be inferred that all the positive
real roots will necessarily be comprised between the
limits
\[
\frac{1}{K + 1} \quad\text{and}\quad k + 1,
\]
and that the negative real roots will fall between the
limits
\[
-\frac{1}{H + 1} \quad\text{and}\quad -h - 1.
\]

There are methods for finding still closer limits;
but since they require considerable labor, the preceding
\PageSep{110}
method is, in the majority of cases, preferable, as
being more simple and convenient.

For example, if in the proposed equation $l + z$ be
\MNote{A further method for finding the limits.}
\index{Roots!method for finding the limits of}%
substituted for~$x$, and if after having arranged the
terms according to the powers of~$z$, there be given to~$l$
a value such that the coefficients of all the terms
become positive, it is plain that there will then be no
positive value of~$z$ that can satisfy the equation. The
equation will have negative roots only, and consequently
$l$~will be a quantity greater than the greatest
value of~$x$. Now it is easy to see that these coefficients
will be expressed as follows:
\begin{gather*}
%[** TN: Re-broken]
p + ml, \\
q + (m - 1)pl + \frac{m(m - 1)}{2}\, l^{2}, \\
r + (m - 2)ql + \frac{(m - 1)(m - 2)}{2}\, pl^{2}
  + \frac{m(m - 1)(m - 2)}{2·3}\, l^{3},
\end{gather*}
and so on. Accordingly, it is only necessary to seek
by trial the smallest value of~$l$ which will render them
all positive.

But in the majority of cases it is not sufficient to
\index{Problems}%
know the limits of the roots of an equation; the thing
necessary is to know the values of those roots, at
least as approximately as the conditions of the problem
require. For every problem leads in its last analysis
to an equation which contains its solution; and
if it is not in our power to resolve this equation, all
\PageSep{111}
the pains expended upon its formulation are a sheer
loss. We may regard this point, therefore, as the
most important in all analysis, and for this reason I
\MNote{The real problem, the finding of the roots.}
have felt constrained to make it the principal subject
of the present lecture.

From the principles established above regarding
\index{Substitutions|EtSeq}%
the nature of the curve of which the ordinates~$y$ represent
all the values which the left-hand side of an
equation assumes, it follows that if we possessed
some means of describing this curve we should obtain
at once, by its intersections with the axis, all the roots
of the proposed equation. But for this purpose it is
not necessary to have all of the curve; it is sufficient
to know the parts which lie immediately above and
below each point of intersection. Now it is possible
to find as many points of a curve as we please, and as
near to one another as we please by successively substituting
for~$x$ numbers which are very little different
from one another, but which are still near enough for
our purpose, and by taking for~$y$ the results of these
substitutions in the left-hand side of the equation. If
among the results of these substitutions two be found
having contrary signs, we may be certain, by the principles
established above, that there will be between
these two values of~$x$ at least one real root. We can
then by new substitutions bring these two limits still
closer together and approach as nearly as we wish to
the roots sought.

Calling the smaller of the two values of~$x$ which
\PageSep{112}
have given results with contrary signs,~$A$, and the
larger~$B$, and supposing that we wish to find the
\MNote{Separation of the roots.}
\index{Roots!separation of the}%
\index{Roots!arithmetical@the arithmetical progression revealing the|EtSeq}%
value of the root within a degree of exactness denoted
by~$n$, where $n$~is a fraction of any degree of smallness
we please, we proceed to substitute successively for~$x$
the following numbers in arithmetical progression:
\index{Arithmetical progression revealing the roots|EtSeq}%
\[
A + n,\quad A + 2n,\quad A + 3n, \dots,
\]
or
\[
B - n,\quad B - 2n,\quad B - 3n, \dots,
\]
until a result is reached having the contrary sign to
that obtained by the substitution of~$A$ or of~$B$. Then
one of the two successive values of~$x$ which have given
results with contrary signs will necessarily be larger
than the root sought, and the other smaller; and since
by hypothesis these values differ from one another
only by the quantity~$n$, it follows that each of them
approaches to within less than~$n$ of the root sought,
and that the error is therefore less than~$n$.

But how are the initial values substituted for~$x$ to
be determined, so as on the one hand to avoid as
many useless trials as possible, and on the other to
make us confident that we have discovered by this
method all the real roots of this equation. If we examine
the curve of the equation it will be readily seen
that the question resolves itself into so selecting the
values of~$x$ that at least one of them shall fall between
two adjacent intersections, which will be necessarily
the case if the difference between two consecutive values
\PageSep{113}
is less than the smallest distance between two
adjacent intersections.

Thus, supposing that $D$~is a quantity smaller than
the smallest distance between two intersections immediately
\MNote{To find a quantity less than the difference between any two roots.}
\index{Roots!quantity less than the difference between any two}%
following each other, we form the arithmetical
progression
\[
0,\quad D,\quad 2D,\quad 3D,\quad 4D,\dots,
\]
and we select from this progression only the terms
which fall between the limits
\[
\frac{1}{K + 1} \quad\text{and}\quad k + 1,
\]
as determined by the method already given. We obtain,
in this manner, values which on being substituted
for~$x$ ultimately give us all the positive roots of
the equation, and at the same time give the initial
limits of each root. In the same manner, for obtaining
the negative roots we form the progression
\[
0,\quad -D,\quad -2D,\quad -3D,\quad -4D,\dots,
\]
from which we also take only the terms comprised
between the limits
\[
-\frac{1}{H + 1} \quad\text{and}\quad -h - 1.
\]

Thus this difficulty is resolved. But it still remains
to find the quantity~$D$,---that is, a quantity
smaller than the smallest interval between any two adjacent
intersections of the curve with the axis. Since
the abscissæ which correspond to the intersections are
\index{Intersections, with the axis give roots}%
the roots of the proposed equation, it is clear that the
question reduces itself to finding a quantity smaller
\PageSep{114}
than the smallest difference between two roots, neglecting
the signs. We have, therefore, to seek, by the
methods which were discussed in the lectures of the
principal course, the equation whose roots are the differences
between the roots of the proposed equation.
And we must then seek, by the methods expounded
above, a quantity smaller than the smallest root of
this last equation, and take that quantity for the value
of~$D$.

This method, as we see, leaves nothing to be desired
\MNote{The equation of differences.}
\index{Differences, the equation of|EtSeq}%
as regards the rigorous solution of the problem,
but it labors under great disadvantage in requiring
extremely long calculations, especially if the proposed
equation is at all high in degree. For example, if $m$~is
the degree of the original equation, that of the equation
of differences will be~$m(m - 1)$, because each root
can be subtracted from all the remaining roots, the
number of which is~$m - 1$,---which gives $m(m - 1)$
differences. But since each difference can be positive
or negative, it follows that the equation of differences
must have the same roots both in a positive and in a
negative form; that consequently the equation must
be wanting in all terms in which the unknown quantity
is raised to an odd power; so that by taking the
square of the differences as the unknown quantity, this
unknown quantity can occur only in the $\dfrac{m(m - 1)}{2}$th
degree. For an equation of the $m$th~degree, accordingly,
there is requisite at the start a transformed
\PageSep{115}
equation of the $\dfrac{m(m - 1)}{2}$th degree, which necessitates
an enormous amount of tedious labor, if $m$~is at all
large. For example, for an equation of the $10$th~degree,
\MNote{Impracticability of the method.}
the transformed equation would be of the~$45$th.
And since in the majority of cases this disadvantage
renders the method almost impracticable, it is of great
importance to find a means of remedying it.

To this end let us resume the proposed equation of
the $m$th~degree,
\[
x^{m} + px^{m-1} + qx^{m-2} + \dots + u = 0,
\]
of which the roots are $a$,~$b$,~$c$,~$\dots$. We shall have
then
\[
a^{m} + pa^{m-1} + qa^{m-2} + \dots + u = 0
\]
and also
\[
b^{m} + pb^{m-1} + qb^{m-2} + \dots + u = 0.
\]
Let $b - a = i$. Substitute this value of~$b$ in the second
equation, and after developing the different powers of~$a + i$
according to the well known binomial theorem,
\index{Binomial theorem}%
arrange the resulting equation according to the powers
of~$i$, beginning with the lowest. We shall have the
transformed equation
\[
P + Qi + Ri^{2} + \dots + i^{m} = 0,
\]
in which the coefficients $P$,~$Q$,~$R$,~$\dots$ have the following
values
\begin{align*}
P &= a^{m} + pa^{m-1} + qa^{m-2} + \dots + u, \displaybreak[1] \\
Q &= ma^{m-1} + (m - 1)pa^{m-2} + (m - 2)qa^{m-3} + \dots\Add{,} \displaybreak[1] \\
\PageSep{116}
R &= \begin{aligned}[t]
  \frac{m(m - 1)}{2}\, a^{m-2}
  &+ \frac{(m - 1)(m - 2)}{2}\, pa^{m-3} \\
  &+ \frac{(m - 2)(m - 3)}{2}\, qa^{m-4} + \dots\Add{,}
\end{aligned}
\end{align*}
\MNote{Attempt to remedy the method.}
and so on. The law of formation of these expressions
is evident.

Now, by the first equation in~$a$ we have~$P = 0$.
Rejecting, therefore, the term~$P$ of the equation in~$i$
and dividing all the remaining terms by~$i$, the equation
in question will be reduced to the $(m - 1)$th~degree,
and will have the form
\[
Q + Ri + Si^{2} + \dots + i^{m-1} = 0.
\]

This equation will have for its roots the $m - 1$~differences
between the root~$a$ and the remaining roots
$b$,~$c$,~$\dots$\Add{.} Similarly, if $b$~be substituted for~$a$ in the expressions
for the coefficients $Q$,~$R$,~$\dots$, we shall obtain
an equation of which the roots are the difference
between the root~$b$ and the remaining roots $a$,~$c$,~$\dots$,
and so on.

Accordingly, if a quantity can be found smaller
\index{Roots!smallest|EtSeq}%
than the smallest root of all these equations, it will
possess the property required and may be taken for
the quantity~$D$, the value of which we are seeking.

If, by means of the equation $P = 0$, $a$~be eliminated
from the equation in~$i$, we shall get a new equation in~$i$
which will contain all the other equations of which
we have just spoken, and of which it would only be
necessary to seek the smallest root. But this new
\PageSep{117}
equation in~$i$ is nothing else than the equation of differences
which we sought to dispense with.

\MNote{Further improvement.}
In the above equation in~$i$ let us put it $i = \dfrac{1}{z}$. We
shall have then the transformed equation in~$z$,
\[
z^{m-1}
  + \frac{R}{Q}\, z^{m-2}
  + \frac{S}{Q}\, z^{m-3} + \dots + \frac{1}{Q} = 0,
\]
and the greatest negative coefficient of this equation
will, from what has been demonstrated above, give a
value greater than its greatest root; so that calling~$L$
this greatest coefficient, $L + 1$~will be a quantity
greater than the greatest value of~$z$. Consequently,
$\dfrac{1}{L + 1}$ will be a quantity smaller than the smallest
positive value of~$i$; and in like manner we shall find
a quantity smaller than the smallest negative value
of~$i$. Accordingly, we may take for~$D$ the smallest of
these two quantities, or some quantity smaller than
either of them.

For a simpler result, and one which is independent
of signs, we may reduce the question to finding a
quantity~$L$ numerically greater than any of the coefficients
\index{Coefficients!greatest negative}%
of the equation in~$z$, and it is clear that if we
find a quantity~$N$ numerically smaller than the smallest
value of~$Q$ and a quantity~$M$ numerically greater
than the greatest value of any of the quantities $R$,
$S$,~$\dots$, we may put $L = \dfrac{M}{N}$.

Let us begin with finding the values of~$M$. It is
not difficult to demonstrate, by the principles established
above, that if $k + 1$~is the limit of the positive
\PageSep{118}
roots and $-h - 1$~the limit of the negative roots of
the proposed equation, and if for~$a$, $k + 1$~and~$-h - 1$
\MNote{Final resolution.}
be successively substituted in the expressions for $R$,
$S$,~$\dots$, considering only the terms which have the
same sign as the first,---it is easy to demonstrate that
we shall obtain in this manner quantities which are
greater than the greatest positive and negative values
of $R$, $S$,~$\dots$ corresponding to the roots $a$,~$b$, $c$\Add{,}~$\dots$ of
the proposed equation; so that we may take for~$M$
the quantity which is numerically the greatest of
these.

It accordingly only remains to find a value smaller
than the smallest value of~$Q$. Now it would seem
that we could arrive at this in no other way than by
employing the equation of which the different values
of~$Q$ are the roots,---an equation which can only be
reached by eliminating~$a$ from the following equations:
\begin{gather*}
a^{m} + pa^{m-1} + qa^{m-2} + \dots + u = 0, \\
ma^{m-1} + (m - 1)pa^{m-2} + (m - 2)qa^{m-3} + \dots = Q.
\end{gather*}

It can be easily demonstrated by the theory of
elimination that the resulting equation in~$Q$ will be of
the $m$th~degree, that is to say, of the same degree with
the proposed equation; and it can also be demonstrated
from the form of the roots of this equation
that its next to the last term will be missing. If, accordingly,
we seek by the method given above a quantity
numerically smaller than the smallest root of this
equation, the quantity found can be taken for~$N$. The
\PageSep{119}
problem is therefore resolved by means of an equation
of the same degree as the proposed equation.

The upshot of the whole is \Typo{a}{as} follows,---where for
\MNote{Recapitulation.}
the sake of simplicity I retain the letter~$x$ instead of
the letter~$a$.

Let the following be the proposed equation of the
$m$th~degree:
\[
x^{m} + px^{m-1} + qx^{m-2} + rx^{m-3} + \dots = 0;
\]
let $k$~be the largest coefficient of the negative terms,
and $m - n$~the exponent of~$x$ in the first negative term.
Similarly, let $h$ be the greatest coefficient of the terms
having a contrary sign to the first term after $x$~has
been changed into~$-x$; and let $m - n'$ be the exponent
of~$x$ in the first term having a contrary sign to
the first term of the equation as thus altered. Putting, then,
\[
f = \sqrt[n]{k} + 1 \quad\text{and}\quad g = \sqrt[n]{h} + 1,
\]
we shall have $f$~and~$-g$ for the limits of the positive
and negative roots. These limits are then substituted
\index{Roots!limits of the positive and negative}%
successively for~$x$ in the following formulæ, neglecting
the terms which have the same sign as the first
term:
\begin{gather*}
%[** TN: Re-broken]
\begin{aligned}
\frac{m(m - 1)}{2}\, x^{m-2}
  &+ \frac{(m - 1)(m - 2)}{2}\, px^{m-3} \\
  &+ \frac{(m - 2)(m - 3)}{2}\, qx^{m-4} + \dots,
\end{aligned} \\
\frac{m(m - 1)(m - 2)}{2·3}\, x^{m-3}
  + \frac{(m - 1)(m - 2)(m - 3)}{2·3}\, px^{m-4} + \dots,
\end{gather*}
\PageSep{120}
and so on. Of these formulæ there will be~$m - 2$. Let
the greatest of the numerical quantities obtained in
this manner be called M. We then take the equation
\MNote{The arithmetical progression revealing the roots.}
\index{Arithmetical progression revealing the roots}%
\index{Roots!arithmetical@the arithmetical progression revealing the}%
\[
mx^{m-1} + (m - 1)px^{m-2} + (m - 2)qx^{m-3} + (m - 3)rx^{m-4} + \dots = y
\]
and eliminate~$x$ from it by means of the proposed
equation,---which gives an equation in~$y$ of the $m$th~degree
with its next to the last term wanting. Let $V$~be
the last term of this equation in~$y$, and $T$~the largest
coefficient of the terms having the contrary sign
to~$V$, supposing $y$~positive as well as negative. Then
taking these two quantities $T$~and~$V$ positive, $N$~will
be determined by the equation
\[
\frac{N}{1 - N} = \sqrt[n]{\frac{V}{T}}
\]
where $n$~is equal to the exponent of the last term having
the contrary sign to~$V$. We then take $D$ equal to
or smaller than the quantity~$\dfrac{N}{M + N}$, and interpolate
the arithmetical progression:
\[
0,\quad D,\quad 2D,\quad 3D,\dots,\quad
-D,\quad -2D,\quad -3D, \dots
\]
between the limits $f$~and~$-g$. The terms of these
progressions being successively substituted for~$x$ in
the proposed equation will reveal all the real roots,
positive as well as negative, by the changes of sign
in the series of results produced by these substitutions,
and they will at the same time give the first
limits of these roots,---limits which can be narrowed
as much as we please, as we already know.
\index{Limits of roots|)}%
\PageSep{121}

If the last term~$V$ of the equation in~$y$ resulting
from the elimination of~$x$ is zero, then $N$~will be zero,
and consequently $D$~will be equal to zero. But in
\MNote{Method of elimination\Add{.}}
\index{Elimination!method of}%
this case it is clear that the equation in~$y$ will have
one root equal to zero and even two, because its next
to the last term is wanting. Consequently the equation
\[
mx^{m-1} + (m - 1)px^{m-2} + (m - 2)qx^{m-3} + \dots = 0\Typo{.}{}
\]
will hold good at the same time with the proposed
equation. These two equations will, accordingly, have
\index{Common divisor of two equations}%
\index{Equations!common divisor of two}%
a common divisor which can be found by the ordinary
method, and this divisor, put equal to zero, will give
one or several roots of the proposed equation, which
roots will be double or multiple, as is easily apparent
from the preceding theory; for if the last term~$Q$ of
the equation in~$i$ is zero, it follows that
\[
i = 0 \quad\text{and}\quad a = b.
\]
The equation in~$y$ is reduced, by the vanishing of its
last term, to the $(m - 2)$th~degree,---being divisible
by~$y^{2}$. If after this division its last term should still
be zero, this would be an indication that it had more
than two roots equal to zero, and so on. In such a
contingency we should divide it by~$y$ as many times
as possible, and then take its last term for~$V$, and the
greatest coefficient of the terms of contrary sign to~$V$
for~$T$, in order to obtain the value of~$D$, which will
enable us to find all the remaining roots of the proposed
equation. If the proposed equation is of the
third degree, as
\PageSep{122}
\[
x^{3} + qx + r = 0,
\]
we shall get for the equation in~$y$,
\[
y^{3} + 3qy^{2} - 4q^{3} - 27r^{2} = 0.
\]

If the proposed equation is
\[
x^{4} + qx^{2} + rx + s = 0
\]
we shall obtain for the equation in~$y$ the following
\begin{multline*}
%[** TN: Re-broken]
y^{4} + 8ry^{3} + (4q^{3} - 16qs + 18r^{2})y^{2} \\
  + 256s^{3} - 128s^{2}q^{2} + 16sq^{4} + 144r^{2}sq - 4r^{2}q^{3} - 27r^{4}
  = 0
\end{multline*}
and so on.

Since, however, the finding of the equation in~$y$ by
\MNote{General formulæ for elimination.}
\index{Elimination!general formulæ for}%
the ordinary methods of elimination may be fraught
with considerable difficulty, I here give the general
formulæ for the purpose, derived from the known
properties of equations. We form, first, from the coefficients
$p$,~$q$,~$r$ of the proposed equation, the quantities
$x_{1}$,~$x_{2}$,~$x_{3}$,~$\dots$, in the following manner:
\[
\begin{array}{r@{\,}l}
x_{1} &= -p, \\
x_{2} &= -px_{1} - 2q, \\
x_{3} &= -px_{2} - qx_{1} - 3r, \\
\hdotsfor{2}.
\end{array}
\]
We then substitute in the expressions for $y$,~$y^{2}$,~$y^{3}$,~$\dots$
up to~$y^{m}$, after the terms in~$x$ have been developed
the quantities $x_{1}$~for~$x$, $x_{2}$~for~$x^{}$, $x_{3}$~for~$x^{3}$, and so forth,
and designate by $y_{1}$,~$y_{2}$, $y_{3}$,~$\dots$ the values of $y$,~$y^{2}$, $y^{3}$,~$\dots$
resulting from these substitutions. We have then
simply to form the quantities $A$,~$B$,~$C$ from the formulæ
\PageSep{123}
\index{Differences, the equation of}%
\[
\begin{array}{r@{\,}l}
A &= y_{1}, \\
B &= \dfrac{Ay_{1} - y_{2}}{2}, \\
C &= \dfrac{By_{1} - Ay_{2} + y_{3}}{3}, \\
\hdotsfor{2},
\end{array}
\]
and we shall have the following equation in~$y$:
\[
y^{m} - Ay^{m-1} + By^{m-2} - Cy^{m-3} + \dots = 0.
\]

The value, or rather the limit of~$D$, which we find
by the method just expounded may often be much
\MNote{General result.}
smaller than is necessary for finding all the roots, but
there would be no further inconvenience in this than
to increase the number of successive substitutions for~$x$
\index{Substitutions}%
in the proposed equation. Furthermore, when there
are as many results found as there are units in the
highest exponent of the equation, we can continue
these results as far as we wish by the simple addition
of the first, second, third differences, etc., because
the differences of the order corresponding to the degree
of the equation are always constant.

We have seen above how the curve of the proposed
equation can be constructed by successively giving
different values to the abscissæ~$x$ and taking for the
ordinates~$y$ the values resulting from these substitutions
in the left-hand side of the equation. But these
values for~$y$ can also be found by another very simple
construction, which deserves to be brought to your
notice. Let us represent the proposed equation by
\[
a + bx + cx^{2} + dx^{3} + \dots = 0
\]
\PageSep{124}
where the terms are taken in the inverse order. The
equation of the curve will then be
\[
y = a + bx + cx^{2} + dx^{3} + \dots\Add{.}
\]
Drawing (Fig.~2) the straight line~$OX$, which we take
\MNote{A second construction for solving equations.}
\index{Equations!constructions for solving}%
\index{Machine for solving equations|(}%
as the axis of abscissæ with $O$~as origin, we lay off on
this line the segment~$OI$ equal to the unit in terms of
which we may suppose the quantities $a$,~$b$,~$c$\Add{,}~$\dots$, to
be expressed; and we erect at the points~$OI$ the perpendiculars
\Figure{2}{0.5\textwidth}
$OD$,~$IM$. We then lay off upon the line~$OD$
the segments
\[
OA = a,\quad AB = b,\quad BC = c,\quad CD = d, \dots,
\]
and so on. Let $OP = x$, and at the point~$P$ let the
perpendicular~$PT$\Typo{}{ }be erected. Suppose, for example,
that $d$~is the last of the coefficients $a$,~$b$,~$c$,~$\dots$, so that
the proposed equation is only of the third degree, and
that the problem is to find the value of
\[
y = a + bx + cx^{2} + dx^{3}.
\]
The point~$D$ being the last of the points determined
upon the perpendicular~$OD$, and the point~$C$ the next
\PageSep{125}
to the last, we draw through~$D$ the line~$DM$ parallel
to the axis~$OI$, and through the point~$M$ where this
line cuts the perpendicular~$IM$ we draw the straight
\MNote{The development and solution.}
line~$CM$ connecting $M$ with~$C$. Then through the
point~$S$ where this last straight line cuts the perpendicular~$PT$,
we draw $HSL$ parallel to~$OI$, and through
the point~$L$ where this parallel cuts the perpendicular~$IM$
we draw to the point~$B$ the straight line~$BL$.
Similarly, through the point~$R$, where this last line
cuts the perpendicular~$PT$, we draw $GRK$ parallel to~$OI$,
and through the point~$K$, where this parallel cuts
the perpendicular~$IM$ we draw to the first division
point~$A$ of the perpendicular~$DO$ the straight line~$AK$.
The point~$Q$ where this straight line cuts the perpendicular~$PT$
will give the segment $PQ = y$.

Through $Q$ draw the line $FQ$ parallel to the axis~$OP$.
The two similar triangles $CDM$~and~$CHS$ give
\[
DM(1) : DC(d) = HS(x) : CH(= dx).
\]
Adding $CB(c)$ we have
\[
BH = c + dx.
\]
Also the two similar triangles $BHL$~and~$BGR$ give
\[
HL(1) : HB(c + dx)= GR(x) : BG(= cx + dx^{2}).
\]
Adding $AB(b)$ we have
\[
AG = b + cx + dx^{2}.
\]
Finally the similar triangles $AGK$~and~$AFQ$ give
\[
%[** TN: Set on two lines in original]
GK(1) : GA(b + cx + dx^{2}) = FQ(x) : FA(= bx + cx^{2} + dx^{3}),
\]
and we obtain by adding $OA(a)$
\[
OF = PQ = a + bx + cx^{2} + dx^{3} = y.
\]
\PageSep{126}

The same construction and the same demonstration
hold, whatever be the number of terms in the
proposed equation. When negative coefficients occur
among $a$,~$b$, $c$,~$\dots$, it is simply necessary to take
them in the opposite direction to that of the positive
coefficients. For example, if $a$~were negative we
should have to lay off the segment~$OA$ below the axis~$OI$.
Then we should start from the point~$A$ and add
to it the segment $AB = b$. If $b$~were positive, $AB$~would
be taken in the direction of~$OD$; but if $b$~were
negative, $AB$~would be taken in the opposite direction,
and so on with the rest.

With regard to~$x$, $OP$~is taken in the direction of~$OI$,
which is supposed to be equal to positive unity,
when $x$~is positive; but in the opposite direction when
$x$~is negative.

It would not be difficult to construct, on the foregoing
\MNote{A machine for solving equations.}
\index{Equations!machine@a machine for solving}%
model, an instrument which would be applicable
to all values of the coefficients $a$,~$b$, $c$,~$\dots$, and which
by means of a number of movable and properly jointed
rulers would give for every point~$P$ of the straight
line~$OP$ the corresponding point~$Q$, and which could
be even made by a continuous movement to describe
the curve. Such an instrument might be used for
solving equations of all degrees; at least it could be
used for finding the first approximate values of the
roots, by means of which afterwards more exact values
could be reached.
\index{Machine for solving equations|)}%
\index{Numerical equations!resolution of|)}%
\PageSep{127}


\Lecture[The Employment of Curves.]
{V.}{On the Employment of Curves in the Solution
of Problems.}
\index{Curves!employment of in the solution of problems|(}%
\index{Problems!employment of curves in the solution of|(}%

\First{As long} as algebra and geometry travelled separate
\index{Algebra!application of geometry to|EtSeq}%
\index{Geometry!application of to algebra|EtSeq}%
paths their advance was slow and their
\MNote{Geometry applied to algebra.}
applications limited. But when these two sciences
joined company, they drew from each other fresh vitality
and thenceforward marched on at a rapid pace
towards perfection. It is to Descartes that we owe
\index{Descartes}%
the application of algebra to geometry,---an application
which has furnished the key to the greatest discoveries
in all branches of mathematics. The method
which I last expounded to you for finding and demonstrating
divers general properties of equations by considering
the curves which represent them, is, properly
speaking, a species of application of geometry to algebra,
and since this method has extended \Typo{applicacations}{applications},
and is capable of readily solving problems
whose direct solution would be extremely difficult or
even impossible, I deem it proper to engage your attention
in this lecture with a further view of this subject,---especially
\PageSep{128}
since it is not ordinarily found in
elementary works on algebra.

You have seen how an equation of any degree
\MNote{Method of resolution by curves.}
whatsoever can be resolved by means of a curve, of
which the abscissæ represent the unknown quantity
of the equation, and the ordinates the values which
the left-hand member assumes for every value of the
unknown quantity. It is clear that this method can be
applied generally to all equations, whatever their form,
and that it only requires them to be developed and
arranged according to the different powers of the unknown
quantity. It is simply necessary to bring all
the terms of the equation to one side, so that the other
side shall be equal to zero. Then taking the unknown
quantity for the abscissa~$x$, and the function of the
unknown quantity, or the quantity compounded of
that quantity and the known quantities, which forms
one side of the equation, for the ordinate~$y$, the curve
described by these co-ordinates $x$~and~$y$ will give by
its intersections with the axis those values of~$x$ which
are the required roots of the equation. And since
most frequently it is not necessary to know all possible
values of the unknown quantity but only such as
solve the problem in hand, it will be sufficient to describe
that portion of the curve which corresponds to
these roots, thus saving much unnecessary calculation.
We can even determine in this manner, from the shape
of the curve itself, whether the problem has possible
solutions satisfying the proposed conditions.
\PageSep{129}

Suppose, for instance, that it is required to find on
\index{Light, law of the intensity of}%
\index{Lights, problem of the two|EtSeq}%
the line joining two luminous points of given intensity,
the point which receives a given quantity of light,---the
\MNote{Problem of the two lights.}
law of physics being that the intensity of light decreases
with the square of the distance.

Let $a$~be the distance between the two lights and
$x$~the distance between the point sought and one of
the lights, the intensity' of which at unit distance is~$M$,
the intensity of the other at that distance being~$N$.
The expressions $\dfrac{M}{x^{2}}$ and $\dfrac{N}{(a - x)^{2}}$, accordingly,
give the intensity of the two lights at the point in
question, so that, designating the total given effect by~$A$,
we have the equation
\[
\frac{M}{x^{2}} + \frac{N}{(a - x)^{2}} = A\Add{,}
\]
or
\[
\frac{M}{x^{2}} + \frac{N}{(a - x)^{2}} - A = 0.
\]

We will now consider the curve having the equation
\[
\frac{M}{x^{2}} + \frac{N}{(a - x)^{2}} - A = y
\]
in which it will be seen at once that by giving to~$x$ a
very small value, positive or negative, the term~$\dfrac{M}{x^{2}}$,
while continuing positive, will grow very large, because
a fraction increases in proportion as its denominator
decreases, and it will be infinite when $x = 0$.
Further, if $x$~be made to increase, the expression~$\dfrac{M}{x^{2}}$
will constantly diminish; but the other expression~$\dfrac{N}{(a - x)^{2}}$,
\PageSep{130}
which was $\dfrac{N}{a^{2}}$ when $x = 0$, will constantly increase
until it becomes very large or infinite when $x$
has a value very near to or equal to~$a$.

Accordingly, if, by giving to~$x$ values from zero to~$a$,
\MNote{Various solutions.}
the sum of these two expressions can be made to
become less than the given quantity~$A$, then the value
of~$y$, which at first was very large and positive, will
become negative, and afterwards again become very
large and positive. Consequently, the curve will cut
the axis twice between the two lights, and the problem
will have two solutions. These two solutions will
be reduced to a single solution if the smallest value of
\[
\frac{M}{x^{2}} + \frac{N}{(a - x)^{2}}
\]
is exactly equal to~$A$, and they will become imaginary
if that value is greater than~$A$, because then the value
of~$y$ will always be positive from $x = 0$ to $x = a$.
Whence it is plain that if one of the conditions of the
problem be that the required point shall fall between
the two lights it is possible that the problem has no
solution. But if the point be allowed to fall on the
prolongation of the line joining the two lights, we
shall see that the problem is always resolvable in two
ways. In fact, supposing $x$~negative, it is plain that
the term~$\dfrac{M}{x^{2}}$ will always remain positive and from being
very large when $x$~is near to zero, it will commence
and keep decreasing as $x$~increases until it grows very
small or becomes zero when $x$~is very great or infinite.
\PageSep{131}
The other term~$\dfrac{N}{(a - x)^{2}}$, which at first was equal to~$\dfrac{N}{a^{2}}$,
also goes on diminishing until it becomes zero
when $x$~is negative infinity. It will be the same if $x$~is
positive and greater than~$a$; for when $x = a$, the
expression $\dfrac{N}{(a - x)^{2}}$ will be infinitely great; afterwards
it will keep on decreasing until it becomes zero when $x$~is
infinite, while the other expression $\dfrac{M}{x^{2}}$ will first be
equal to $\dfrac{M}{a^{2}}$ and will also go on diminishing towards
zero as $x$~increases.

Hence, whatever be the value of the quantity~$A$,
it is plain that the values of~$y$ will necessarily pass
\MNote{General solution.}
from positive to negative, both for $x$~negative and for
$x$~positive and greater than~$a$. Accordingly, there
will be a negative value of~$x$ and a positive value of~$x$
greater than~$a$ which will resolve the problem in all
cases. These values may be found by the general
method by successively causing the values of~$x$ which
give values of~$y$ with contrary signs, to approach
nearer and nearer to each other.

With regard to the values of~$x$ which are less than~$a$
we have seen that the reality of these values depends
on the smallest value of the quantity
\[
\frac{M}{x^{2}} + \frac{N}{(a - x)^{2}}.
\]
Directions for finding the smallest and greatest values
of variable quantities are given in the Differential Calculus.
\index{Differential Calculus}%
We shall here content ourselves with remarking
\PageSep{132}
that the quantity in question will be a minimum
when
\MNote{Minimal values.}
\index{Minimal values}%
\index{Values!minimal}%
\[
\frac{x}{a - x} = \sqrt[3]{\frac{M}{N}};
\]
so that we shall have
\[
x = \frac{a\sqrt[3]{M}}{\sqrt[3]{M} + \sqrt[3]{N}},
\]
from which we get, as the smallest value of the expression
\[
\frac{M}{x^{2}} + \frac{N}{(a - x)^{2}},
\]
the quantity
\[
\frac{(\sqrt[3]{M} + \sqrt[3]{N})^{3}}{a^{2}}.
\]
Hence there will be two real values for~$x$ if this quantity
is less than~$A$; but these values will be imaginary
if it is greater. The case of equality will give two
equal values for~$x$.

I have dwelt at considerable length on the analysis
of this problem, (though in itself it is of slight importance,)
for the reason that it can be made to serve
as a type for all analogous cases.

The equation of the foregoing problem, having
been freed from fractions, will assume the following
form:
\[
Ax^{2}(a - x)^{2} - M(a - x)^{2} - Nx^{2} = 0.
\]
With its terms developed and properly arranged it
will be found to be of the fourth degree, and will consequently
have four roots. Now by the analysis which
we have just given, we can recognise at once the character
\PageSep{133}
of these roots. And since a method may spring
from this consideration applicable to all equations of
\index{Equations!fourth@of the fourth degree}%
\index{Fourth degree, equations of the}%
the fourth degree, we shall make a few brief remarks
\MNote{Preceding analysis applied to bi-quadratic equations.}
\index{Biquadratic equations}%
upon it in passing. Let the general equation be
\[
x^{4} + px^{2} + qx + r = 0.
\]
We have already seen that if the last term of this
equation be negative it will necessarily have two real
roots, one positive and one negative; but that if the
last term be positive we can in general infer nothing
as to the character of its roots. If we give to this
equation the following form
\[
(x^{2} - a^{2})^{2} + b(x + a)^{2} + c(x - a)^{2} = 0,
\]
a form which developed becomes
\[
x^{4} + (b + c - 2a^{2})x^{2} + 2a(b - c)x + a^{4} + a^{2}(b + c) = 0,
\]
and from this by comparison derive the following
equations of condition
\[
b + c - 2a^{2} = p,\quad 2a(b - c) = q,\quad a^{4} + a^{2}(b + c) = r,
\]
and from these, again, the following,
\[
b + c = p + 2a^{2},\quad b - c = \frac{q}{2a},\quad 3a^{4} + pa^{2} = r,
\]
we shall obtain, by resolving the last equation,
\[
a^{2} = -\frac{p}{6} + \sqrt{\frac{r}{3} + \frac{p^{2}}{36}}.
\]
If $r$~be supposed positive, $a^{2}$~will be positive and real,
and consequently $a$~will be real, and therefore, also,
$b$~and~$c$ will be real.

Having determined in this manner the three quantities
$a$,~$b$,~$c$, we obtain the transformed equation
\[
(x^{2} - a^{2})^{2} + b(x + a)^{2} + c(x - a)^{2} = 0.
\]
\PageSep{134}

Putting the right-hand side of this equation equal
to~$y$, and considering the curve having for abscissæ
\MNote{Consideration of equations of the fourth degree.}
the different values of~$y$, it is plain, that when $b$~and~$c$
are positive quantities this curve will lie wholly
above the axis and that consequently the equation
will have no real root. Secondly, suppose that $b$~is a
negative quantity and $c$~a positive quantity; then $x = a$
will give $y = 4ba^{2}$,---a negative quantity. A very
large positive or negative~$x$ will then give a very large
positive~$y$,---whence it is easy to conclude that the
equation will have two real roots, one larger than~$a$
and one less than~$a$. We shall likewise find that if
$b$~is positive and $c$~is negative, the equation will have
two real roots, one greater and one less than~$-a$.
Finally, if $b$~and~$c$ are both negative, then $y$~will become
negative by making
\[
x = a \quad\text{and}\quad x = -a
\]
and it will be positive and very large for a very large
positive or negative value of~$x$,---whence it follows
that the equation will have two real roots, one greater
than~$a$ and one less than~$-a$. The preceding considerations
might be greatly extended, but at present we
must forego their pursuit.

It will be seen from the preceding example that
the consideration of the curve does not require the
equation to be freed from fractional expressions. The
\index{Fractional expressions in equations}%
\index{Radical expressions in equations}%
same may be said of radical expressions. There is
an advantage even in retaining these expressions in
\PageSep{135}
the form given by the analysis of the problem; the
advantage being that we may in this way restrict our
attention to those signs of the radicals which answer
\MNote{Advantages of the  method of curves.}
\index{Curves!advantages of the method of}%
to the special exigencies of each problem, instead of
causing the fractions and the radicals to disappear
and obtaining an equation arranged according to the
different whole powers of the unknown quantity in
which frequently roots are introduced which are entirely
foreign to the question proposed. It is true that
these roots are always part of the question viewed in
its entire extent; but this wealth of algebraical analysis,
although in itself and from a general point of view
extremely valuable, may be inconvenient and burdensome
in particular cases where the solution of which
we are in need cannot by direct methods be found independently
of all other possible solutions. When
the equation which immediately flows from the conditions
of the problem contains radicals which are essentially
ambiguous in sign, the curve of that equation
(constructed by making the side which is equal to
zero, equal to the ordinate~$y$) will necessarily have as
many branches as there are possible different combinations
of these signs, and for the complete solution it
would be necessary to consider each of these branches.
But this generality may be restricted by the particular
conditions of the problem which determine the branch
on which the solution is to be sought; the result being
that we are spared much needless calculation,---an
advantage which is not the least of those offered by
\PageSep{136}
the method of solving equations from the consideration
of curves.

But this method can be still further generalised
\MNote{The curve of errors.}
\index{Errors, curve of|EtSeq}%
and even rendered independent of the equation of the
problem. It is sufficient in applying it to consider
the conditions of the problem in and for themselves,
to give to the unknown quantity different arbitrary
values, and to determine by calculation or construction
the errors which result from such suppositions
according to the original conditions. Taking these
errors as the ordinates~$y$ of a curve having for abscissæ
the corresponding values of the unknown quantity,
we obtain a continuous curve called \emph{the curve of errors},
which by its intersections with the axis also gives all
solutions of the problem. Thus, if two successive errors
be found, one of which is an excess, and another
a defect, that is, one positive and one negative, we
may conclude at once that between these two corresponding
values of the unknown quantity there will
be one for which the error is zero, and to which we
can approach as near as we please by successive substitutions,
or by the mechanical description of the
curve.

This mode of resolving questions by curves of errors
\index{Astronomy, mechanics, and physics, curves of errors in}%
\index{Mechanics, astronomy, and physics, curves of errors in}%
\index{Physics, astronomy, and mechanics, curves of errors in}%
is one of the most useful that have been devised.
It is constantly employed in astronomy when direct
solutions are difficult or impossible. It can be employed
for resolving important problems of geometry
and mechanics and even of physics. It is properly
\PageSep{137}
speaking the \textit{regula falsi}, taken in its most general
\index{False, rule of}%
\index{Regula@\textit{Regula falsi}}%
\index{Rule!false@of false}%
sense and rendered applicable to all questions where
there is an unknown quantity to be determined. It
\MNote{Solution of a problem by the curve of errors.}
can also be applied to problems that depend on two
or several unknown quantities by successively giving
to these unknown quantities different arbitrary values
and calculating the errors which result therefrom, afterwards
linking them together by different curves, or
reducing them to tables; the result being that we may
\index{Tables}%
by this method obtain directly the solution sought
\Figure{3}{0.4\textwidth}
without preliminary elimination of the unknown quantities.

We shall illustrate its use by a few examples.

\textit{Required a circle in which a polygon of given sides can
be inscribed.}

This problem gives an equation which is proportionate
in degree to the number of sides of the polygon.
To solve it by the method just expounded we
describe any circle~$ABCD$ (Fig.~3) and lay off in this
circle the given sides $AB$,~$BC$, $CD$, $DE$,~$EF$ of the
\PageSep{138}
polygon, which for the sake of simplicity I here suppose
to be pentagonal. If the extremity of the last
\MNote{Problem of the circle and inscribed polygon.}
\index{Circle!and inscribed polygon, problem of the}%
\index{Polygon, problem of the circle and inscribed}%
side falls on~$A$, the problem is solved. But since it
is very improbable that this should happen at the first
trial we lay off on the straight line~$PR$ (Fig.~4) the
radius~$PA$ of the circle, and erect on it at the point~$A$
the perpendicular~$AF$ equal to the chord~$AF$ of the
arc~$AF$ which represents the error in the supposition
\index{Supposition, rule of}%
\index{Trial and error, rule of}%
made regarding the length of the radius~$PA$. Since
this error is an excess, it will be necessary to describe
\Figure{4}{0.3\textwidth}
a circle having a larger radius and to perform the
same operation as before, and so on, trying circles of
various sizes. Thus, the circle having the radius~$PA$
gives the error~$F'A'$ which, since it falls on the hither
side of the point~$A'$, should be accounted negative. It
will consequently be necessary in Fig.~4 in applying
the ordinate~$A'F'$ to the abscissa~$PA'$ to draw that
ordinate below the axis. In this manner we shall obtain
several points $F$,~$F'$,~$\dots$, which will lie on a
curve of which the intersection~$R$ with the axis~$PA$
\PageSep{139}
will give the true radius~$PR$ of the circle satisfying
the problem, and we shall find this intersection by
successively causing the points of the curve lying on
\MNote{Solution of a second problem by the curve of errors.}
the two sides of the axis as $F$,~$F'$,~$\dots$ to approach
nearer and nearer to one another.

\textit{From a point, the position of which is unknown, three
\index{Point in space, position of a}%
objects are observed, the distances of which from one another
are known. The three angles formed by the rays of
light from these three objects to the eye of the observer are
also known. Required the position of the observer with
respect to the three objects.}

If the three objects be joined by three straight
lines, it is plain that these three lines will form with
the visual rays from the eye of the observer a triangular
pyramid of which the base and the three face angles
forming the solid angle at the vertex are given.
And since the observer is supposed to be stationed at
the vertex, the question is accordingly reduced to determining
the dimensions of this pyramid.

Since the position of a point in space is completely
determined by its three distances from three given
points, it is clear that the problem will be resolved, if
the distances of the point at which the observer is
stationed from each of the three objects can be determined.
Taking these three distances as the unknown
quantities we shall have three equations of the second
degree, which after elimination will give a resultant
equation of the eighth degree; but taking only one of
these distances and the relations of the two others to it
\PageSep{140}
for the unknown quantities, the final equation will be
only of the fourth degree. We can accordingly rigorously
\MNote{Problem of the observer and three objects.}
solve this problem by the known methods; but
the direct solution, which is complicated and inconvenient
in practice, may be replaced by the following
which is reached by the curve of errors.

Let the three successive angles $APB$, $BPC$, $CPD$
\index{Observer, problem of the, and three objects}%
(Fig.~5) be constructed, having the vertex~$P$ and
respectively equal to the angles observed between the
first object and the second, the second and the third,
\Figure{5}{0.4\textwidth}
the third and the first; and let the straight line~$PA$
be taken at random to represent the distance from the
observer to the first object. Since the distance of
that object to the second is supposed to be known,
let it be denoted by~$AB$, and let it be laid off on the
line~$AB$. We shall in this way obtain the distance~$BP$
of the second object to the observer. In like manner,
let $BC$, the distance of the second object to the
third, be laid off on~$BC$, and we shall have the distance~$PC$
of that object to the observer. If, now, the
\PageSep{141}
distance of the third object to the first be laid off on
the line~$CD$, we shall obtain~$PD$ as the distance of
the first object to the observer. Consequently, if the
\MNote{Employment of the curve of errors.}
distance first assumed is exact, the two lines $PA$~and~$PD$
will necessarily coincide. Making, therefore, on
the line~$PA$, prolonged if necessary, the segment
$PE = PD$, if the point~$E$ does not fall upon the point~$A$,
the difference will be the error of the first assumption~$PA$.
Having drawn the straight line~$PR$ (Fig.~6)
we lay off upon it from the fixed point~$P$, the abscissa~$PA$,
and apply to it at right angles the ordinate~$EA$;
we shall have the point~$E$ of the curve of errors~$ERS$.
\Figure{6}{0.4\textwidth}
Taking other distances for~$PA$, and making the same
construction, we shall obtain other errors which can be
similarly applied to the line~$PR$, and which will give
other points in the same curve.

We can thus trace this curve through several
points, and the point~$R$ where it cuts the axis~$PR$ will
give the distance~$PR$, of which the error is zero, and
which will consequently represent the exact distance
of the observer from the first object. This distance
being known, the others may be obtained by the same
construction.

It is well to remark that the construction we have
been considering gives for each point~$A$ of the line~$PA$,
\PageSep{142}
two points $B$~and~$B'$ of the line~$PB$; for, since
the distance~$AB$ is given, to find the point~$B$ it is only
\MNote{Eight possible solutions of the preceding problem.}
necessary to describe from the point~$A$ as centre and
with radius~$AB$ an arc of a circle cutting the straight
line~$PB$ at the two points $B$~and~$B'$,---both of which
points satisfy the conditions of the problem. In the
same manner, each of these last-mentioned points will
give two more upon the straight line~$PC$, and each of
the last will give two more on the straight line~$PD$.
Whence it follows that every point~$A$ taken upon the
straight line~$PA$ will in general give eight upon the
straight line~$PD$, all of which must be separately and
successively considered to obtain all the possible solutions.
I have said, \emph{in general}, because it is possible
(1)~for the two points $B$~and~$B'$ to coincide at a single
point, which will happen when the circle described
with the centre~$A$ and radius~$AB$ touches the straight
line~$PB$; and (2)~that the circle may not cut the
straight line~$PB$ at all, in which case the rest of the
construction is impossible, and the same is also to be
said regarding the points $C$,~$D$. Accordingly, drawing
the line~$GF$ parallel to~$BP$ and at a distance from it
equal to the given line~$AB$, the point~$F$ at which this
line cuts the line~$PE$, prolonged if necessary, will be
the limit beyond which the points~$A$ must not be taken
if we desire to obtain possible solutions. There exist
also limits for the points $B$~and~$C$, which may be employed
in restricting the primitive suppositions made
with respect to the distance~$PA$.
\PageSep{143}

The eight points~$D$, which depend in general on
each point~$A$, answer to the eight solutions of which
the problem is susceptible, and when one has no special
\MNote{Reduction of the possible solutions in practice.}
datum by means of which it can be determined
which of these solutions answer best to the case proposed,
it is indispensable to ascertain them all by employing
for each one of the eight combinations a special
curve of errors. But if it be known, for example,
that the distance of the observer to the second object
is greater or less than his distance to the first, it will
then be necessary to take on the line~$PB$ only the
point~$B$ in the first case and the point~$B'$ in the second,---a
course which will reduce the eight combinations
one-half. If we had the same datum with regard
to the third object relatively to the second, and with
regard to the first object relatively to the third, then
the points $C$~and~$D$ would be determined, and we
should have but a single solution.

These two examples may suffice to illustrate the
uses to which the method of curves can be put in solving
\index{Curves!method of, submitted to analysis|EtSeq}%
problems. But this method, which we have presented,
so to speak, in a mechanical manner, can also
be submitted to analysis.

The entire question in fact is reducible to the description
of a curve which shall pass through a certain
number of points, whether these points be given by
calculation or construction, or whether they be given
by observation or single experiences entirely independent
of one another. The problem is in truth indeterminate,
\PageSep{144}
for strictly speaking there can be made
to pass through a given number of points an infinite
\MNote{General conclusion on the method of curves.}
\index{Curves!advantages of the method of}%
number of different curves, regular or irregular, that
is, subject to equations or arbitrarily drawn by the
hand. But the question is not to find any solutions
whatever but the simplest and easiest in practice.

Thus if there are only two points given, the simplest
solution is a straight line between the two points.
\index{Straight line}%
If there are three points given, the arc of a circle is
\index{Circle}%
drawn through these points, for the arc of a circle
after the straight line is the simplest line that can be
described.

But if the circle is the simplest curve with respect
to description, it is not so with respect to the equation
between its abscissæ and rectangular ordinates.
In this latter point of view, those curves may be regarded
as the simplest of which the ordinates are expressed
by an integral rational function of the abscissæ,
as in the following equation
\[
y = a + bx + cx^{2} + dx^{3} + \dots,
\]
where $y$~is the ordinate and $x$~the abscissa. Curves
of this class are called in general \emph{parabolic}, because
\index{Parabolic@\textit{Parabolic} curves|EtSeq}%
they may be regarded as a generalisation of the parabola,---a
curve represented by the foregoing equation
when it has only the first three terms. We have already
illustrated their employment in resolving equations,
and their consideration is always useful in the
approximate description of curves, for the reason that
a curve of this kind can always be made to pass
\PageSep{145}
through as many points of a given curve as we please,---it
being only necessary to take as many undetermined
coefficients $a$,~$b$,~$c$,~$\dots$ as there are points given,
\MNote{Parabolic curves.}
and to determine these coefficients so as to obtain the
abscissæ and ordinates for these points. Now it is
clear that whatever be the curve proposed, the parabolic
curve so described will always differ from it by
less and less according as the number of the different
points is larger and larger and their distance from
one another smaller and smaller.

Newton was the first to propose this problem. The
\index{Newton, his problem}%
following is the solution which he gave of it:

Let $P$,~$Q$, $R$,~$S$,~$\dots$ be the values of the ordinates~$y$
corresponding to the values $p$,~$q$, $r$,~$s$,~$\dots$ of
the abscissæ~$x$; we shall have the following equations
\[
\begin{array}{r@{\,}*{3}{l@{\,}}l}
P &= a + bp &+ cp^{2} &+ dp^{3} &+ \dots, \\
Q &= a + bq &+ cq^{2} &+ dq^{3} &+ \dots, \\
R &= a + br &+ cr^{2} &+ dr^{3} &+ \dots, \\
\hdotsfor{5}\Add{.}
\end{array}
\]
The number of these equations must be equal to the
number of the undetermined coefficients $a$,~$b$,~$c$,~$\dots$.
Subtracting these equations from one another, the remainders
will be divisible by $q - p$, $r - q$,~$\dots$, and
we shall have after such division
\[
\begin{array}{r@{\,}*{2}{l@{\,}}l}
\dfrac{Q - P}{q - p} &= b + c(q + p) &= d(q^{2} + qp + p^{2}) &+ \dots, \\[8pt]
\dfrac{R - Q}{r - q} &= b + c(r + q) &= d(r^{2} + rq + q^{2}) &+ \dots, \\
\hdotsfor{4}\Add{.}
\end{array}
\]
\PageSep{146}

Let
\[
\frac{Q - P}{q - p} = Q_{1},\quad
\frac{R - Q}{r - q} = R_{1},\quad
\frac{S - R}{s - r} = S_{1},\dots\Add{.}
\]
\MNote{Newton's problem.}
We shall find in like manner, by subtraction and division,
the following:
\[
\begin{array}{r@{\,}l@{\,}l}
\dfrac{R_{1} - Q_{1}}{r - p} &= c + d(r + q + p) &+ \dots, \\[8pt]
\dfrac{S_{1} - R_{1}}{s - q} &= c + d(s + r + q) &+ \dots, \\
\hdotsfor{3}\Add{.}
\end{array}
\]

Further let
\[
\frac{R_{1} - Q_{1}}{r - p} = R_{2},\quad
\frac{S_{1} - R_{1}}{s - q} = S_{2},\dots.
\]
We shall have
\[
\frac{S_{2} - R_{2}}{s - p} = d + \dots,
\]
and so on.

In this manner we shall find the value of the coefficients
$a$,~$b$,~$c$,~$\dots$ commencing with the last; and,
substituting them in the general equation
\[
y = a + bx + cx^{2} + dx^{3} + \dots,
\]
we shall obtain, after the appropriate reductions have
been made, the formula
\[
y = P
  + Q_{1}(x - p)
  + R_{2}(x - p)(x - q)
  + S_{3}(x - p)(x - q)(x - r) + \dots,
\Tag{(1)}
\]
which can be carried as far as we please.

But this solution may be simplified by the following
consideration.

Since $y$~necessarily becomes $P$,~$Q$,~$R$\Add{,}~$\dots$, when $x$~becomes
\PageSep{147}
$p$,~$q$,~$r$, it is easy to see that the expression
for~$y$ will be of the form
\MNote{Simplification of Newton's solution.}
\[
y = AP + BQ + CR + DS + \dots
\Tag{(2)}
\]
where the quantities $A$,~$B$, $C$,~$\dots$ are so expressed in
terms of~$x$ that by making $x = p$ we shall have
\[
A = 1,\quad B = 0,\quad C = 0,\dots,
\]
and by making $x = q$ we shall have
\[
A = 0,\quad B = 1,\quad C = 0,\quad D = 0,\dots,
\]
and by making $x = r$ we shall similarly have
\[
A = 0,\quad B = 0,\quad C = 1,\quad D = 0,\dots\ \text{etc.}
\]
Whence it is easy to conclude that the values of $A$,
$B$, $C$,~$\dots$ must be of the form
\begin{align*}
A &= \frac{(x - q)(x - r)(x - s)\dots}{(p - q)(p - r)(p - s)\dots}, \\
B &= \frac{(x - p)(x - r)(x - s)\dots}{(q - p)(q - r)(q - s)\dots}, \\
C &= \frac{(x - p)(x - q)(x - s)\dots}{(r - p)(r - q)(r - s)\dots},
\end{align*}
where there are as many factors in the numerators
and denominators as there are points given of the
curve less one.

The last expression for~$y$ (see equation~2), although
different in form, is the same as equation~1. To show
this, the values of the quantities $Q_{1}$,~$R_{2}$, $S_{3}$,~$\dots$ need
only be developed and substituted in equation~1 and
the terms arranged with respect to the quantities $P$,
$Q$, $R$,~$\dots$\Add{.} But the last expression for~$y$ (equation~2)
is preferable, partly because of the simplicity of the
\PageSep{148}
analysis from which it is derived, and also because of
its form, which is more convenient for computation.

\MNote{Possible uses of Newton's problem.}
Now, by means of this formula, which it is not
difficult to reduce to a geometrical construction, we
are able to find the value of the ordinate~$y$ for any abscissa~$x$,
because the ordinates $P$,~$Q$, $R$,~$\dots$ for the
given abscissæ $p$,~$q$, $r$,~$\dots$ are known. Thus, if we
have several of the terms of any series, we can find
any intermediate term that we wish,---an expedient
which is extremely valuable for supplying lacunæ
which may arise in a series of observations or experiments,
\index{Experiments!expedient@an expedient for supplying lacunæ in a series of}%
\index{Observations, expedient for supplying lacunæ in series of}%
or in tables calculated by formulæ or in given
\index{Tables!expedient for supplying lacunæ in}%
constructions.

If this theory now be applied to the two examples
\index{Regula@\textit{Regula falsi}}%
\index{Supposition, rule of}%
\index{Trial and error, rule of}%
discussed above and to similar examples in which we
have errors corresponding to different suppositions, we
can directly find the error~$y$ which corresponds to any
intermediate supposition~$x$ by taking the quantities
$P$,~$Q$, $R$,~$\dots$, for the errors found, and $p$,~$q$, $r$,~$\dots$ for
the suppositions from which they result. But since
in these examples the question is to find not the error
which corresponds to a given supposition, but the
supposition for which the error is zero, it is clear that
the present question is the opposite of the preceding
and that it can also be resolved by the same formula
by reciprocally taking the quantities $p$,~$q$, $r$,~$\dots$ for
the errors, and the quantities $P$,~$Q$, $R$,~$\dots$ for the
corresponding suppositions. Then $x$~will be the error
for the supposition~$y$; and consequently, by making
\PageSep{149}
$x = 0$, the value of~$y$ will be that of the supposition
for which the error is zero.

Let $P$,~$Q$, $R$,~$\dots$ be the values of the unknown
quantity in the different suppositions, and $p$,~$q$, $r$\Add{,}~$\dots$
\MNote{Application of Newton's problem to the preceding examples.}
the errors resulting from these suppositions, to which
the appropriate signs are given. We shall then have
for the value of the unknown quantity of which the
error is zero, the expression
\[
AP + BQ + CR + \dots,
\]
in which the values of $A$,~$B$,~$C$\Add{,}~$\dots$ are
\begin{align*}
A &= \frac{q}{q - r} × \frac{r}{r - p} × \dots, \displaybreak[1] \\
B &= \frac{P}{p - q} × \frac{r}{r - q} × \dots, \displaybreak[1] \\
C &= \frac{p}{p - r} × \frac{q}{q - r} × \dots,
\end{align*}
where as many factors are taken as there are suppositions
less one.
\index{Curves!employment of in the solution of problems|)}%
\index{Problems!employment of curves in the solution of|)}%
\PageSep{150}
%[Blank page]
\PageSep{151}


\Appendix{Note on the Origin of Algebra.}
\PgLabel{151}
\index{Algebra!history of}%

\First{The} impression (\PgRef{54}) that Diophantus was the
\index{Diophantus}%
``inventor'' of algebra, which sprang, in its Diophantine
form, full-fledged from his brain, was a widespread
one in the eighteenth and in the beginning of
the nineteenth century. But, apart from the intrinsic
improbability of this view which is at variance with
the truth that science is nearly always gradual and
organic in growth, modern historical researches have
traced the germs and beginnings of algebra to a much
remoter date, even in the line of European historical
continuity. The Egyptian book of Ahmes contains
\index{Ahmes}%
examples of equations of the first degree. The early
Greek mathematicians performed the partial resolution
\index{Greeks, mathematics of the}%
of equations of the second and third degree
by geometrical methods. According to Tannery, an
\index{Tannery, M. Paul}%
embryonic indeterminate analysis existed in Pre-Christian
times (Archimedes, Hero, Hypsicles). But
\index{Archimedes}%
\index{Hero}%
\index{Hypsicles}%
the merit of Diophantus as organiser and inaugurator
of a more systematic short-hand notation, at
least in the European line, remains; he enriched
whatever was handed down to him with the most
manifold extensions and applications, betokening his
\PageSep{152}
originality and genius, and carried the science of algebra
\index{Algebra!among the Arabs}%
\index{Algebra!India@in India}%
to its highest pitch of perfection among the
\PgLabel{152}
Greeks. (See Cantor, \textit{Geschichte der Mathematik}, second
\index{Cantor}%
edition, Vol.~I., p.~438, et~seq.; Ball, \textit{Short Account
\index{Ball}%
of the History of Mathematics}, second edition, p.~104
et~seq.; Fink, \textit{A Brief History of Mathematics}, pp.~63
\index{Fink}%
et~seq., 77~et~seq. (Chicago: The Open Court
Publishing~Co.)

The development of Hindu algebra is also to be
noted in connexion with the text of \PgRange{59}{60}. The
Arabs, who had considerable commerce with India,
\index{Arabs!Algebra among the}%
drew not a little of their early knowledge from the
works of the Hindus. Their algebra rested on both
that of the Hindus and the Greeks. (See Ball, \textit{op.~cit.},
p.~150 et~seq.; Cantor, \textit{op.~cit.}, Vol.~I., p.~651 et~seq.).---\textit{Trans.}
\PageSep{153}
\BackMatter
\printindex
\iffalse
INDEX.

Academies, rise of 62, 63

Ahmes 151

Algebra
  definition of 2
  history of|EtSeq#history 54 % et seq.,
  history of 151
  essence of 55
  name@the name of 59
  among the Arabs|EtSeq 59 % et seq,
  among the Arabs 152
  Europe@in Europe 60
  Italy@in Italy 64
  India@in India 152
  generality@the generality of 69
  hand-writing of 69
  application of geometry to|EtSeq 100, 127 % et seq.

Algebraical resolution of equations
  limits of the 96

Alligation
  generally|EtSeq 44 % et seq.;
  alternate 47

Analysis
  indeterminate|EtSeq 47 % et seq.,
  indeterminate 55

Angle, trisection of an 62, 81

Angular sections, theory of 80

Annuities 16

Apollonius 54, 59

Arabs
  Algebra among the|EtSeq 59 % et seq.,
  Algebra among the 152

Archimedes 54, 151

Archimedes|FN 58 % footnote

Arithmetic
  universal|EtSeq 2 % et seq.;
  operations of|EtSeq 24 % et seq.

Arithmetical progression revealing the roots 120

Arithmetical progression revealing the roots|EtSeq 112 % et seq.

Arithmetical proportion 12

Astronomy, mechanics, and physics, curves of errors in 136

Average life|EtSeq 45 % et seq.

Bachet de Méziriac 58

Ball 152

Binomial theorem 115

Binomials, extraction of the square roots of two imaginary 77

Biquadratic equations 63, 88, 94, 133

Bombelli 63, 64

Bret, M.|FN 93 % footnote.

Briggs 20

Buteo 61

Cantor|FN 54, 60 % footnote,

Cantor 152

Cardan 60, 61, 68, 82, 90

Checks on multiplication and division 39

Circle 144
  squaring of the 62
  and inscribed polygon, problem of the 138

Clairaut 69, 90

Coefficients
  indeterminate 89
  greatest negative|EtSeq 107 % et seq.,
  greatest negative 117

Common divisor of two equations 121

Complements, subtraction by 26

Constantinople 58

Continued fractions, solution of alligation by|EtSeq 50 % et seq.

Convergents 7

Cube, duplication of the 62

Cube roots of a quantity, the three 70

Cubic radicals 75

Curves
  representation of equations by|EtSeq 101 % et seq;
  employment of in the solution of problems 127-149
  method of, submitted to analysis|EtSeq 143 % et seq.;
  advantages of the method of 135, 144

Decimal
  fractions 9
  numbers|EtSeq 27 % et seq.

Decimals
  multiplication of 30
  division of 31
\PageSep{154}

DeMorgan@{\Typo{DeMorgan}{De Morgan}} v

Descartes viii, 60, 65, 89, 93, 127

Differences, the equation of|EtSeq 114 % et seq.,

Differences, the equation of 123

Differential Calculus 131

Diophantine problems 55

Diophantus|EtSeq 54 % et seq

Diophantus 151

Division
  nine@by \textit{nine} 34
  eight@by \textit{eight} 34
  seven@by \textit{seven}|EtSeq 34 % et seq.;
  decimals@of decimals 31

Divisor, greatest common|EtSeq 2 % et seq.

Duhring@{Dühring, E.} v

Duodecimal system 32

Ecole@{\Typo{Ecole}{École} Normale} v, xi, 12

Economy of thought vii

Efflux, law of 42

Eleven, the number, test of divisibility by 37

Elimination
  method of 121
  general formulæ for 122

Equations
  second@of the second degree 56
  third@of the third degree 60, 66, 82
  fourth@of the fourth degree 63, 87, 133
  fifth@of the fifth degree 64
  theory of 65, 84
  biquadratic 88
  limits of the algebraical resolution of 96
  fifth@of the fifth degree 96
  mth@of the $m$th degree 96
  general remarks upon the roots of|EtSeq 102 % et seq.;
  graphic resolution of 102
  odd@of an odd degree, roots of 105
  even@of an even degree, roots of 106
  real roots of, limits of the|EtSeq 107 % et seq.;
  common divisor of two 121
  constructions for solving|EtSeq 100 % et seq.
  constructions for solving 124
  machine@a machine for solving 126

Equi-different numbers 13

Errors, curve of|EtSeq 136 % et seq.

Euclid 2, 57

Euler viii, x, 93

Europe, algebra in 60

Evolution 11, 40

Experiments
  average of 46
  expedient@an expedient for supplying lacunæ in a series of 148

Falling stone, spaces traversed by a 42

False, rule of 137

Fermat 58

Ferrari, Louis 64

Ferrous, Scipio|EtSeq 60 % et seq.

Fifth degree, equations of the 96

Fink 152

Fourth degree, equations of the 133

Fractional expressions in equations 134

Fractions|EtSeq 2 % et seq.;

Fractions
  continued|EtSeq 3 % et seq.;
  converging 6
  decimal 9
  origin of continued 10

France 58, 61

Galileo ix

Geometers, ancient|EtSeq 54 % et seq.

Geometers, ancient 58, 59

Geometrical
  proportion 13
  calculus 24

Geometry 24, 60
  application of to algebra|EtSeq 100, 127 % et seq.

Germany 61

Girard, Albert 62

Grain, of different prices 44

Greeks, mathematics of the vii, 151

Greeks, mathematics of the|EtSeq 54 % et seq.

Hand-writing of algebra 69

Harriot 65

Hero 59, 151

Horses 43

Hudde 65, 82

Huygens ix, 10

Hypsicles 151

Imaginary binomials, square roots of 77

Imaginary expressions|EtSeq 79 % et seq.

Imaginary expressions 83

Imaginary quantities, office of the 87

Imaginary roots, occur in pairs 99

Indeterminate analysis|EtSeq 47 % et seq.

Indeterminate analysis 55

Indeterminate coefficients 89

Indeterminates, the method of 83

Ingredients 48

Interest 15

Intersections, with the axis give roots|EtSeq 102 % et seq ,

Intersections, with the axis give roots 113

Inventors, great 22

Involution and evolution 11

Irreducible case 61, 65, 69, 73, 82

Italy, cradle of algebra in Europe 61, 64

Laborers, work of 41

Lagrange, J. L.#Lagrange v

Lagrange, J. L.|EtSeq#Lagrange vii % et seq.
\PageSep{155}

Laplace v, xi

Lavoisier xii

Leibnitz viii

Life insurance|EtSeq 45 % et seq.

Life, probability of 46

Light, law of the intensity of 129

Lights, problem of the two|EtSeq 129 % et seq.

Limits of roots 107-120

Logarithms|EtSeq 16 % et seq.

Logarithms 40
  advantages in calculating by 28
  origin of 19
  tables of 20

Machine for solving equations 124-126

Mathematics
  wings of 24
  exactness of 43
  evolution of vii

Mean values|EtSeq 45 % et seq.

Mechanics, astronomy, and physics, curves of errors in 136

Metals, mingling of, by fusion 44

Meziriac@Méziriac, Bachet de 58

Minimal values 132

Mixtures, rule of|EtSeq 44 % et seq.

Mixtures, rule of 49

Monge v, xi

Mortality, tables of 45

Moving bodies, two 98

Multiple roots 105

Multiplication
  abridged methods of|EtSeq 26 % et seq.;
  inverted 28
  approximate 29
  decimals@of decimals 30

Music 22

Napier|EtSeq 17 % et seq.

Napoleon xii

Negative roots 60

Newton, his problem 145, viii

Nine
  property of the number|EtSeq 31 % et seq.;
  property of the number generalised 33

Nizze|FN 58 % footnote.

Numeration, systems of 1

Numerical equations |See Equations 0

Numerical equations
  resolution of 96-126
  conditions of the resolution of 97
  position of the roots of 98

Observations, expedient for supplying lacunæ in series of 148

Observer, problem of the, and three objects 140

Oughtred 30

Paciolus, Lucas 59, 60

Pappus 59

Parabolic@\textit{Parabolic} curves|EtSeq 144 % et seq.

Peletier 61

Peyrard 58

Physics, astronomy, and mechanics, curves of errors in 136

Planetarium 9

Point in space, position of a 139

Polygon, problem of the circle and inscribed 138

Polytechnic School v, xi

Positive roots, superior and inferior limits of the 109

Powers|EtSeq 10 % et seq.

Practice, theory and 43

Present value 15

Printing, invention of 59

Probabilities, calculus of|EtSeq 45 % et seq.

Problems 110
  solution@for solution 62
  employment of curves in the solution of 127-149

Proclus 59

Progressions, theory of 12, 14

Proportion|EtSeq 11 % et seq.

Ptolemy 59

Radical expressions in equations 134

Radicals, cubic 75

Ratios, constant 42

Ratios, constant|EtSeq 2, 11 % et seq.

Reality of roots 76, 83, 85, 93

Regula@\textit{Regula falsi} 137, 148

Remainders
  theory of|EtSeq 34 % et seq.
  theory of 38
  negative|EtSeq 35 % et seq.

Romans, mathematics of the 54

Roots
  negative 60
  equations@of equations of the third degree 71
  reality@the reality of the 74, 76, 79, 83, 85, 93
  biquadratic@of a biquadratic equation 94
  multiple 105
  superior and inferior limits of the positive 109
  method for finding the limits of 110
  separation of the 112
  arithmetical@the arithmetical progression revealing the|EtSeq 112 % et seq.
  arithmetical@the arithmetical progression revealing the 120
  quantity less than the difference between any two 113
  smallest|EtSeq 116 % et seq.;
  limits of the positive and negative 119

Rule
  Cardan's 68
  false@of false 137
  mixtures@of mixtures|EtSeq 44 % et seq.;
  three@of three|EtSeq 11, 40 % et seq.
\PageSep{156}

Science
  history of 22
  development of|EtSeq vii % et seq.

Seven, tests of divisibility by 35

Short-mind symbols|EtSeq vii % et seq.

Signs $+$ and $-$ 57

Squaring of the circle 62

Stenophrenic symbols|EtSeq vii % et seq.

Straight line 144

Substitutions|EtSeq 111 % et seq.

Substitutions 123

Subtraction, new method of|EtSeq 25 % et seq.

Sum and difference, of two numbers 56

Supposition, rule of 137, 148

Symbols|EtSeq vii % et seq.

Tables 137
  expedient for supplying lacunæ in 148

Tannery, M. Paul|FN 58 % footnote

Tannery, M. Paul 151

Tartaglia 60, 61

Temperament, theory of 23

Theon 59

Theory and practice 43

Theory of remainders, utility of the 38

Third degree, equations of the 71, 82

Three roots, reality of the 93

Trial and error, rule of 137, 148

Trisection of an angle 62, 81

Turks 58

Undetermined quantities 82

Unity, three cubic roots of 72

Unknown quantity 55

Values
  mean|EtSeq 45 % et seq.;
  minimal 132

Variations, calculus of x

Vatican library 58

Vieta viii, 62, 65

Vlacq 20

Wallis viii

Wertheim, G.|FN 58 % footnote.

Woodhouse x

Xylander 58
\fi
\PageSep{157}

\Catalog
%[** TN: Macro prints the following text]
% Catalogue of Publications
% of the
% Open Court Publishing Co.

\begin{Author}{COPE, E. D.}
\Title{THE PRIMARY FACTORS OF ORGANIC EVOLUTION.}
{121~cuts. Pp.~xvi,~547. Cloth,~\$2.00 (10s.).}
\end{Author}

\begin{Author}{MÜLLER, F. MAX.}
\Title{THREE INTRODUCTORY LECTURES ON THE SCIENCE OF
THOUGHT.}
{128~pages. Cloth,~75c (3s.\ 6d.).}

\Title{THREE LECTURES ON THE SCIENCE OF LANGUAGE.}
{112~pages. 2nd~Edition. Cloth,~75c (3s.\ 6d.).}
\end{Author}

\begin{Author}{ROMANES, GEORGE JOHN.}
\Title{DARWIN AND AFTER DARWIN.}
{Three Vols., \$4.00. Singly, as follows:}{}

%[** TN: Next three extries get a bit less hanging indentation]
\Title[3\parindent]{}{1.~\textsc{The Darwinian Theory.} 460~pages. 125~illustrations. Cloth, \$2.00\Add{.}}

\Title[3\parindent]{}{2.~\textsc{Post-Darwinian Questions.} Heredity and Utility. Pp.~338. \$1.50\Add{.}}

\Title[3\parindent]{}{3.~\textsc{Post-Darwinian Questions.} Isolation and Physiological Selection
Pp.~181. \$1.00.}

\Title{AN EXAMINATION OF WEISMANNISM.}
{236~pages. Cloth, \$1.00.}

\Title{THOUGHTS ON RELIGION.}
{Third Edition, Pages,~184. Cloth, gilt top, \$1.25.}
\end{Author}

\begin{Author}{SHUTE, DR. D. KERFOOT.}
\Title{FIRST BOOK IN ORGANIC EVOLUTION.}
{9~colored plates, 39~cuts. Pp.~xvi+285. Price, \$2.00 (7s.\ 6d.).}
\end{Author}

\begin{Author}{MACH, ERNST.}
\Title{THE SCIENCE OF MECHANICS.}
{Translated by \textsc{T. J. McCormack.} 250~cuts. 534~pages. \$2.50 (12s.\ 6d.)}

\Title{POPULAR SCIENTIFIC LECTURES.}
{Third Edition. 415~pages. 59~cuts. Cloth, gilt top. \$1.50 (7s.\ 6d.).}

\Title{THE ANALYSIS OF THE SENSATIONS.}
{Pp.~208. 37~cuts. Cloth, \$1.25 (6s.\ 6d.).}
\end{Author}

\begin{Author}{LAGRANGE, JOSEPH LOUIS.}
\Title{LECTURES ON ELEMENTARY MATHEMATICS.}
{With portrait of the author. Pp.~172. Price, \$1.00 (5s.).}
\end{Author}

\begin{Author}{DE MORGAN, AUGUSTUS.}
\Title{ON THE STUDY AND DIFFICULTIES OF MATHEMATICS.}
{New Reprint edition with notes. Pp.~viii+288. Cloth, \$1.25 (5s.).}

\Title{ELEMENTARY ILLUSTRATIONS OF THE DIFFERENTIAL AND
INTEGRAL CALCULUS.}
{New reprint edition. Price, \$1.00 (5s.).}
\end{Author}

\begin{Author}{FINK, KARL.}
\Title{A BRIEF HISTORY OF MATHEMATICS.}
{Trans.\ by W. W. Beman and D. E. Smith. Pp.\Typo{,}{}~333. Cloth, \$1.50 (5s.\ 6d.)}
\end{Author}

\begin{Author}{SCHUBERT, HERMANN.}
\Title{MATHEMATICAL ESSAYS AND RECREATIONS.}
{Pp.~149. Cuts,~37. Cloth, 75c (33.\ 6d.).}
\end{Author}

\begin{Author}{HUC AND GABET, MM.}
\Title{TRAVELS IN TARTARY, THIBET AND CHINA.}
{100~engravings. Pp\Add{.}~28+660. 2~vols. \$2.00 (10s.). One vol., \$1.25 (5s.)}
\end{Author}
\PageSep{158}

\begin{Author}{CARUS, PAUL.}
\Title{THE HISTORY OF THE DEVIL, AND THE IDEA OF EVIL.}
{311~Illustrations. Pages,~500. Price, \$6.00 (30s.).}

\Title{EROS AND PSYCHE.}
{Retold after Apuleius. With Illustrations by Paul Thumann. Pp.~125.
Price, \$1.50 (6s.).}

\Title{WHENCE AND WHITHER?}
{An Inquiry into the Nature of the Soul. 196~pages. Cloth, 75c (3s.\ 6d.)}

\Title{THE ETHICAL PROBLEM.}
{Second edition, revised and enlarged. 351~pages. Cloth, \$1.25 (6s.\ 6d.)}

\Title{FUNDAMENTAL PROBLEMS.}
{Second edition, revised and enlarged. 372~pp.\ Cl., \$1.50 (7s.\ 6d.).}

\Title{HOMILIES OF SCIENCE.}
{317~pages. Cloth, Gilt Top, \$1.50 (7s.\ 6d.).}

\Title{THE IDEA OF GOD.}
{Fourth edition. 32~pages. Paper, 15c (9d.).}

\Title{THE SOUL OF MAN.}
{2nd~ed. 182~cuts. 482~pages. Cloth, \$1.50 (6s.).}

\Title{TRUTH IN FICTION. \textsc{Twelve Tales with a Moral.}}
{White and gold binding, gilt edges. Pp.~111. \$1.00 (5s.).}

\Title{THE RELIGION OF SCIENCE.}
{Second, extra edition. Pp.~103. Price, 50c (2s.\ 6d.).}

\Title{PRIMER OF PHILOSOPHY.}
{240~pages. Second Edition. Cloth, \$1.00 (5s.).}

\Title{THE GOSPEL OF BUDDHA. According to Old Records.}
{Fifth Edition. Pp.~275. Cloth, \$1.00 (5s.). In German, \$1.25 (6s.\ 6d.)\Add{.}}

\Title{BUDDHISM AND ITS CHRISTIAN CRITICS.}
{Pages,~311. Cloth, \$1.25 (6s.\ 6d.).}

\Title{KARMA. \textsc{A Story of Early Buddhism.}}
{Illustrated by Japanese artists. Crêpe paper, 75c (3s.\ 6d.).}

\Title{NIRVANA: \textsc{A Story of Buddhist Psychology.}}
{Japanese edition, like \textit{Karma}. \$1.00 (4s.\ 6d.).}

\Title{LAO-TZE'S TAO-TEH-KING.}
{Chinese-English. Pp.~360. Cloth, \$3.00 (15s.).}
\end{Author}

\begin{Author}{CORNILL, CARL HEINRICH.}
\Title{THE PROPHETS OF ISRAEL.}
{Pp.,~200\Add{.} Cloth, \$1.00 (5s.).}

\Title{HISTORY OF THE PEOPLE OF ISRAEL.}
{Pp.~vi+325. Cloth, \$1.50 (7s. 6d.).}
\end{Author}

\begin{Author}{POWELL, J. W.}
\Title{TRUTH AND ERROR; or, the Science of Intellection.}
{Pp.~423. Cloth, \$1.75 (7s. 6d.).}
\end{Author}

\begin{Author}{RIBOT, TH.}
\Title{THE PSYCHOLOGY OF ATTENTION.}{}

\Title{THE DISEASES OF PERSONALITY.}{}

\Title{THE DISEASES OF THE WILL.}
{Cloth, 75~cents each (3s.\ 6d.). \textit{Full set, cloth, \$1.75} (9s.).}

\Title{EVOLUTION OF GENERAL IDEAS.}
{Pp.~231. Cloth, \$1.25 (5s.).}
\end{Author}

\begin{Author}{WAGNER, RICHARD.}
\Title{A PILGRIMAGE TO BEETHOVEN.}
{A Story. With portrait of Beethoven. Pp.~40. Boards, 50c (2s.\ 6d.).}
\end{Author}

\begin{Author}{HUTCHINSON, WOODS.}
\Title{THE GOSPEL ACCORDING TO DARWIN.}
{Pp.~xii+241. Price, \$1.50 (6s.).}
\end{Author}

\begin{Author}{FREYTAG, GUSTAV.}
\Title{THE LOST MANUSCRIPT. A Novel.}
{2~vols. 953~pages. Extra cloth, \$4.00 (21s\Add{.}). One vol., cl., \$1.00 (5s.)\Add{.}}

\Title{MARTIN LUTHER.}
{Illustrated. Pp.~130. Cloth, \$1.00 (5s.).}
\end{Author}
\PageSep{159}

\begin{Author}{AÇVAGHOSHA.}
\Title{DISCOURSE ON THE AWAKENING OF FAITH in the Mahâyâna.}
{Translated for the first time from the Chinese version by Tietaro
Suzuki. Pages,~176. Price, cloth, \$1.25 (5s.\ 6d.).}
\end{Author}

\begin{Author}{TRUMBULL, M. M.}
\Title{THE FREE TRADE STRUGGLE IN ENGLAND.}
{Second Edition. 296~pages. Cloth,~75c (3s.\ 6d.).}

\Title{WHEELBARROW: \textsc{Articles and Discussions on the Labor Question.}}
{With portrait of the author. 303~pages. Cloth, \$1.00 (5s.).}
\end{Author}

\begin{Author}{GOETHE AND SCHILLER'S XENIONS.}
\Title{Translated by Paul Carus. Album form. Pp.~162. Cl., \$1.00 (5s.).}{}
\end{Author}

\begin{Author}{OLDENBERG, H.}
\Title{ANCIENT INDIA: ITS LANGUAGE AND RELIGIONS.}
{Pp.~100. Cloth, 50c (2s. 6d.).}
\end{Author}

\begin{Author}{CONWAY, DR. MONCURE DANIEL.}
\Title{SOLOMON, AND SOLOMONIC LITERATURE.}
{Pp.~243. Cloth, \$1.50 (6s.).}
\end{Author}

\begin{Author}{GARBE, RICHARD.}
\Title{THE REDEMPTION OF THE BRAHMAN. \textsc{A Tale of Hindu Life.}}
{Laid paper. Gilt top. 96~pages. Price, 75c (3s.\ 6d.).}

\Title{THE PHILOSOPHY OF ANCIENT INDIA.}
{Pp.~89. Cloth, 50c (2s.\ 6d.).}
\end{Author}

\begin{Author}{HUEPPE, FERDINAND.}
\Title{THE PRINCIPLES OF BACTERIOLOGY.}
{28~Woodcuts. Pp.~x+467. Price, \$1.75 (9s.).}
\end{Author}

\begin{Author}{LÉVY-BRUHL, PROF. L.}
\Title{HISTORY OF MODERN PHILOSOPHY IN FRANCE.}
{23 Portraits. Handsomely bound. Pp. 500. Price, \$3.00 (12s.).}
\end{Author}

\begin{Author}{TOPINARD, DR. PAUL.}
\Title{SCIENCE AND FAITH, \textsc{or Man as an Animal and Man as a Member
of Society.}}
{Pp.~374. Cloth, \$1.50 (6s.\ 6d.).}
\end{Author}

\begin{Author}{BINET, ALFRED.}
\Title{THE PSYCHOLOGY OF REASONING.}
{Pp.~193. Cloth, 75c (3s.\ 6d.).}

\Title{THE PSYCHIC LIFE OF MICRO-ORGANISMS.}
{Pp.~135. Cloth, 75 cents.}

\Title{ON DOUBLE CONSCIOUSNESS.}
{See No.~8, Religion of Science Library.}
\end{Author}

\begin{Author}{THE OPEN COURT.}
\Title{A Monthly Magazine Devoted to the Science of Religion, the Religion of
Science, and the Extension of the Religious Parliament Idea.}
{Terms: \$1.00 a year; 5s.\ 6d.\ to foreign countries in the Postal Union.
Single Copies, 10~cents (6d.).}
\end{Author}

\begin{Author}{THE MONIST.}
\Title{A Quarterly Magazine of Philosophy and Science.}
{Per copy, 50~cents; Yearly, \$2.00. In England and all countries in
U.P.U. per copy, 2s.~6d.: Yearly, 9s.~6d.}
\end{Author}

\tb
\vfill
\begin{center}
CHICAGO: \\
\large THE OPEN COURT PUBLISHING CO. \\
\footnotesize Monon Building, 324 Dearborn St. \\
LONDON: Kegan Paul, Trench, Trübner \&~Company, Ltd.
\end{center}
\PageSep{160}
\newpage
\begin{center}
\makebox[0.9\textwidth][s]{\LARGE\itshape The Religion of Science Library.}
\tb
\end{center}

\CatalogSmallFont
A collection of bi-monthly publications, most of which are reprints of
books published by The Open Court Publishing Company. Yearly, \$1.50.
Separate copies according to prices quoted. The books are printed upon
good paper, from large type.

The Religion of Science Library, by its extraordinarily reasonable price
will place a large number of valuable books within the reach of all readers.

The following have already appeared in the series:

\Item{No.\ 1.} \textit{The Religion of Science.} By \textsc{Paul Carus.} 25c (1s.\ 6d.).

\Item{2.} \textit{Three Introductory Lectures on the Science of Thought.} By \textsc{F. Max
Müller.} 25c (1s.\ 6d.).

\Item{3.} \textit{Three Lectures on the Science of Language.} \textsc{F. Max Müller.} 25 (1s.\ 6d.)

\Item{4.} \textit{The Diseases of Personality.} By \textsc{Th.\ Ribot.} 25c (1s.\ 6d.).

\Item{5.} \textit{The Psychology of Attention.} By \textsc{Th.\ Ribot.} 25c (1s.\ 6d.).

\Item{6.} \textit{The Psychic Life of Micro-Organisms.} By \textsc{Alfred Binet.} 25c (1s.\ 6d.)

\Item{7.} \textit{The Nature of the State.} By \textsc{Paul Carus.} 15c (9d.).

\Item{8.} \textit{On Double Consciousness.} By \textsc{Alfred Binet.} 15c (9d.).

\Item{9.} \textit{Fundamental Problems.} By \textsc{Paul Carus.} 50c (2s. 6d.).

\Item{10.} \textit{The Diseases of the Will.} By \textsc{Th.\ Ribot.} 25c (1s.\ 6d.).

\Item{11.} \textit{The Origin of Language.} By \textsc{Ludwig Noire.} 15c (9d.).

\Item{12.} \textit{The Free Trade Struggle in England.} By \textsc{M. M. Trumbull.} 25c (1s.\ 6d.)

\Item{13.} \textit{Wheelbarrow on the Labor Question.} By \textsc{M. M. Trumbull.} 35c (2s.).

\Item{14.} \textit{The Gospel of Buddha.} By \textsc{Paul Carus.} 35c (2s.).

\Item{15.} \textit{The Primer of Philosophy.} By \textsc{Paul Carus.} 25c (1s.\ 6d.).

\Item{16.} \textit{On Memory, and The Specific Energies of the Nervous System.} By \textsc{Prof.\
Ewald Hering.} 15c (9d.).

\Item{17.} \textit{The Redemption of the Brahman. Tale of Hindu Life.} By \textsc{Richard
Garbe.} 25c (1s.\ 6d.).

\Item{18.} \textit{An Examination of Weismannism.} By \textsc{G. J. Romanes.} 35c (2s.).

\Item{19.} \textit{On Germinal Selection.} By \textsc{August Weismann.} 25c (1s.\ 6d.).

\Item{20.} \textit{Lovers Three Thousand Years Ago.} By \textsc{T. A. Goodwin.} (Out of print.)

\Item{21.} \textit{Popular Scientific Lectures.} By \textsc{Ernst Mach.} 50c (2s.\ 6d.).

\Item{22.} \textit{Ancient India: Its Language and Religions.} By \textsc{H. Oldenberg.} 25c
(1s.\ 6d.).

\Item{23.} \textit{The Prophets of Israel.} By \textsc{Prof.\ C. H. Cornill.} 25c (1\Add{s}.\ 6d.).

\Item{24.} \textit{Homilies of Science.} By \textsc{Paul Carus.} 35c (2s.).

\Item{25.} \textit{Thoughts on Religion.} By \textsc{G. J. Romanes.} 50c (2s.\ 6d.).

\Item{26.} \textit{The Philosophy of Ancient India.} By \textsc{Prof.\ Richard Garbe.} 25c (1s.\ 6d.)

\Item{27.} \textit{Martin Luther.} By \textsc{Gustav Freytag.} 25c (1s.\ 6d.).

\Item{28.} \textit{English Secularism.} By \textsc{George Jacob Holyoake.} 25c (1s.\ 6d.).

\Item{29.} \textit{On Orthogenesis.} By \textsc{Th.\ Eimer.} 25c (1s.\ 6d.).

\Item{30.} \textit{Chinese Philosophy.} By \textsc{Paul Carus.} 25c (1s.\ 6d.).

\Item{31.} \textit{The Lost Manuscript.} By \textsc{Gustav Freytag.} 60c (35.).

\Item{32.} \textit{A Mechanico-Physiological Theory of Organic Evolution.} By \textsc{Carl von
Naegeli.} 15c (9d.).

\Item{33.} \textit{Chinese Fiction.} By \textsc{Dr.\ George T. Candlin.} 15c (9d.).

\Item{34.} \textit{Mathematical Essays and Recreations.} By \textsc{H. Schubert.} 25c (1s.\ 6d.)

\Item{35.} \textit{The Ethical Problem.} By \textsc{Paul Carus.} 50c (2s.\ 6d.).

\Item{36.} \textit{Buddhism and Its Christian Critics.} By \textsc{Paul Carus.} 50c (2s.\ 6d.).

\Item{37.} \textit{Psychology for Beginners.} By \textsc{Hiram M. Stanley.} 20c (1s.).

\Item{38.} \textit{Discourse on Method.} By \textsc{Descartes.} 25c (1s.\ 6d.).

\Item{39.} \textit{The Dawn of a New Era.} By \textsc{Paul Carus.} 15c (9d.).

\Item{40.} \textit{Kant and Spencer.} By \textsc{Paul Carus.} 20c (1s.).

\Item{41.} \textit{The Soul of Man.} By \textsc{Paul Carus.} 75c (3s.\ 6d.).

\Item{42.} \textit{World' s Congress Addresses.} By \textsc{C. C. Bonney.} 15c (9d.).

\Item{43.} \textit{The Gospel According to Darwin.} By \textsc{Woods Hutchinson.} 50c (2s.\ 6d.)

\Item{44.} \textit{Whence and Whither.} By \textsc{Paul Carus.} 25c (1s.\ 6d.).

\Item{45.} \textit{Enquiry Concerning Human Understanding.} By \textsc{David Hume.} 25c
(1s.\ 6d.).

\Item{46.} \textit{Enquiry Concerning the Principles of Morals.} By \textsc{David Hume.}
25c (1s.\ 6d.)

\normalsize
\tb
\vfill
\begin{center}
\makebox[\textwidth][s]{\Large THE OPEN COURT PUBLISHING CO.,} \\[4pt]
\normalsize CHICAGO: 324 \textsc{Dearborn Street.} \\[4pt]
\footnotesize \textsc{London}: Kegan Paul, Trench, Trübner \&~Company, Ltd.
\end{center}
%%%%%%%%%%%%%%%%%%%%%%%%% GUTENBERG LICENSE %%%%%%%%%%%%%%%%%%%%%%%%%%
\PGLicense
\begin{PGtext}
End of the Project Gutenberg EBook of Lectures on Elementary Mathematics, by 
Joseph Louis Lagrange

*** END OF THIS PROJECT GUTENBERG EBOOK LECTURES ON ELEMENTARY MATHEMATICS ***

***** This file should be named 36640-pdf.pdf or 36640-pdf.zip *****
This and all associated files of various formats will be found in:
        http://www.gutenberg.org/3/6/6/4/36640/

Produced by Andrew D. Hwang.

Updated editions will replace the previous one--the old editions
will be renamed.

Creating the works from public domain print editions means that no
one owns a United States copyright in these works, so the Foundation
(and you!) can copy and distribute it in the United States without
permission and without paying copyright royalties.  Special rules,
set forth in the General Terms of Use part of this license, apply to
copying and distributing Project Gutenberg-tm electronic works to
protect the PROJECT GUTENBERG-tm concept and trademark.  Project
Gutenberg is a registered trademark, and may not be used if you
charge for the eBooks, unless you receive specific permission.  If you
do not charge anything for copies of this eBook, complying with the
rules is very easy.  You may use this eBook for nearly any purpose
such as creation of derivative works, reports, performances and
research.  They may be modified and printed and given away--you may do
practically ANYTHING with public domain eBooks.  Redistribution is
subject to the trademark license, especially commercial
redistribution.



*** START: FULL LICENSE ***

THE FULL PROJECT GUTENBERG LICENSE
PLEASE READ THIS BEFORE YOU DISTRIBUTE OR USE THIS WORK

To protect the Project Gutenberg-tm mission of promoting the free
distribution of electronic works, by using or distributing this work
(or any other work associated in any way with the phrase "Project
Gutenberg"), you agree to comply with all the terms of the Full Project
Gutenberg-tm License (available with this file or online at
http://gutenberg.org/license).


Section 1.  General Terms of Use and Redistributing Project Gutenberg-tm
electronic works

1.A.  By reading or using any part of this Project Gutenberg-tm
electronic work, you indicate that you have read, understand, agree to
and accept all the terms of this license and intellectual property
(trademark/copyright) agreement.  If you do not agree to abide by all
the terms of this agreement, you must cease using and return or destroy
all copies of Project Gutenberg-tm electronic works in your possession.
If you paid a fee for obtaining a copy of or access to a Project
Gutenberg-tm electronic work and you do not agree to be bound by the
terms of this agreement, you may obtain a refund from the person or
entity to whom you paid the fee as set forth in paragraph 1.E.8.

1.B.  "Project Gutenberg" is a registered trademark.  It may only be
used on or associated in any way with an electronic work by people who
agree to be bound by the terms of this agreement.  There are a few
things that you can do with most Project Gutenberg-tm electronic works
even without complying with the full terms of this agreement.  See
paragraph 1.C below.  There are a lot of things you can do with Project
Gutenberg-tm electronic works if you follow the terms of this agreement
and help preserve free future access to Project Gutenberg-tm electronic
works.  See paragraph 1.E below.

1.C.  The Project Gutenberg Literary Archive Foundation ("the Foundation"
or PGLAF), owns a compilation copyright in the collection of Project
Gutenberg-tm electronic works.  Nearly all the individual works in the
collection are in the public domain in the United States.  If an
individual work is in the public domain in the United States and you are
located in the United States, we do not claim a right to prevent you from
copying, distributing, performing, displaying or creating derivative
works based on the work as long as all references to Project Gutenberg
are removed.  Of course, we hope that you will support the Project
Gutenberg-tm mission of promoting free access to electronic works by
freely sharing Project Gutenberg-tm works in compliance with the terms of
this agreement for keeping the Project Gutenberg-tm name associated with
the work.  You can easily comply with the terms of this agreement by
keeping this work in the same format with its attached full Project
Gutenberg-tm License when you share it without charge with others.

1.D.  The copyright laws of the place where you are located also govern
what you can do with this work.  Copyright laws in most countries are in
a constant state of change.  If you are outside the United States, check
the laws of your country in addition to the terms of this agreement
before downloading, copying, displaying, performing, distributing or
creating derivative works based on this work or any other Project
Gutenberg-tm work.  The Foundation makes no representations concerning
the copyright status of any work in any country outside the United
States.

1.E.  Unless you have removed all references to Project Gutenberg:

1.E.1.  The following sentence, with active links to, or other immediate
access to, the full Project Gutenberg-tm License must appear prominently
whenever any copy of a Project Gutenberg-tm work (any work on which the
phrase "Project Gutenberg" appears, or with which the phrase "Project
Gutenberg" is associated) is accessed, displayed, performed, viewed,
copied or distributed:

This eBook is for the use of anyone anywhere at no cost and with
almost no restrictions whatsoever.  You may copy it, give it away or
re-use it under the terms of the Project Gutenberg License included
with this eBook or online at www.gutenberg.org

1.E.2.  If an individual Project Gutenberg-tm electronic work is derived
from the public domain (does not contain a notice indicating that it is
posted with permission of the copyright holder), the work can be copied
and distributed to anyone in the United States without paying any fees
or charges.  If you are redistributing or providing access to a work
with the phrase "Project Gutenberg" associated with or appearing on the
work, you must comply either with the requirements of paragraphs 1.E.1
through 1.E.7 or obtain permission for the use of the work and the
Project Gutenberg-tm trademark as set forth in paragraphs 1.E.8 or
1.E.9.

1.E.3.  If an individual Project Gutenberg-tm electronic work is posted
with the permission of the copyright holder, your use and distribution
must comply with both paragraphs 1.E.1 through 1.E.7 and any additional
terms imposed by the copyright holder.  Additional terms will be linked
to the Project Gutenberg-tm License for all works posted with the
permission of the copyright holder found at the beginning of this work.

1.E.4.  Do not unlink or detach or remove the full Project Gutenberg-tm
License terms from this work, or any files containing a part of this
work or any other work associated with Project Gutenberg-tm.

1.E.5.  Do not copy, display, perform, distribute or redistribute this
electronic work, or any part of this electronic work, without
prominently displaying the sentence set forth in paragraph 1.E.1 with
active links or immediate access to the full terms of the Project
Gutenberg-tm License.

1.E.6.  You may convert to and distribute this work in any binary,
compressed, marked up, nonproprietary or proprietary form, including any
word processing or hypertext form.  However, if you provide access to or
distribute copies of a Project Gutenberg-tm work in a format other than
"Plain Vanilla ASCII" or other format used in the official version
posted on the official Project Gutenberg-tm web site (www.gutenberg.org),
you must, at no additional cost, fee or expense to the user, provide a
copy, a means of exporting a copy, or a means of obtaining a copy upon
request, of the work in its original "Plain Vanilla ASCII" or other
form.  Any alternate format must include the full Project Gutenberg-tm
License as specified in paragraph 1.E.1.

1.E.7.  Do not charge a fee for access to, viewing, displaying,
performing, copying or distributing any Project Gutenberg-tm works
unless you comply with paragraph 1.E.8 or 1.E.9.

1.E.8.  You may charge a reasonable fee for copies of or providing
access to or distributing Project Gutenberg-tm electronic works provided
that

- You pay a royalty fee of 20% of the gross profits you derive from
     the use of Project Gutenberg-tm works calculated using the method
     you already use to calculate your applicable taxes.  The fee is
     owed to the owner of the Project Gutenberg-tm trademark, but he
     has agreed to donate royalties under this paragraph to the
     Project Gutenberg Literary Archive Foundation.  Royalty payments
     must be paid within 60 days following each date on which you
     prepare (or are legally required to prepare) your periodic tax
     returns.  Royalty payments should be clearly marked as such and
     sent to the Project Gutenberg Literary Archive Foundation at the
     address specified in Section 4, "Information about donations to
     the Project Gutenberg Literary Archive Foundation."

- You provide a full refund of any money paid by a user who notifies
     you in writing (or by e-mail) within 30 days of receipt that s/he
     does not agree to the terms of the full Project Gutenberg-tm
     License.  You must require such a user to return or
     destroy all copies of the works possessed in a physical medium
     and discontinue all use of and all access to other copies of
     Project Gutenberg-tm works.

- You provide, in accordance with paragraph 1.F.3, a full refund of any
     money paid for a work or a replacement copy, if a defect in the
     electronic work is discovered and reported to you within 90 days
     of receipt of the work.

- You comply with all other terms of this agreement for free
     distribution of Project Gutenberg-tm works.

1.E.9.  If you wish to charge a fee or distribute a Project Gutenberg-tm
electronic work or group of works on different terms than are set
forth in this agreement, you must obtain permission in writing from
both the Project Gutenberg Literary Archive Foundation and Michael
Hart, the owner of the Project Gutenberg-tm trademark.  Contact the
Foundation as set forth in Section 3 below.

1.F.

1.F.1.  Project Gutenberg volunteers and employees expend considerable
effort to identify, do copyright research on, transcribe and proofread
public domain works in creating the Project Gutenberg-tm
collection.  Despite these efforts, Project Gutenberg-tm electronic
works, and the medium on which they may be stored, may contain
"Defects," such as, but not limited to, incomplete, inaccurate or
corrupt data, transcription errors, a copyright or other intellectual
property infringement, a defective or damaged disk or other medium, a
computer virus, or computer codes that damage or cannot be read by
your equipment.

1.F.2.  LIMITED WARRANTY, DISCLAIMER OF DAMAGES - Except for the "Right
of Replacement or Refund" described in paragraph 1.F.3, the Project
Gutenberg Literary Archive Foundation, the owner of the Project
Gutenberg-tm trademark, and any other party distributing a Project
Gutenberg-tm electronic work under this agreement, disclaim all
liability to you for damages, costs and expenses, including legal
fees.  YOU AGREE THAT YOU HAVE NO REMEDIES FOR NEGLIGENCE, STRICT
LIABILITY, BREACH OF WARRANTY OR BREACH OF CONTRACT EXCEPT THOSE
PROVIDED IN PARAGRAPH 1.F.3.  YOU AGREE THAT THE FOUNDATION, THE
TRADEMARK OWNER, AND ANY DISTRIBUTOR UNDER THIS AGREEMENT WILL NOT BE
LIABLE TO YOU FOR ACTUAL, DIRECT, INDIRECT, CONSEQUENTIAL, PUNITIVE OR
INCIDENTAL DAMAGES EVEN IF YOU GIVE NOTICE OF THE POSSIBILITY OF SUCH
DAMAGE.

1.F.3.  LIMITED RIGHT OF REPLACEMENT OR REFUND - If you discover a
defect in this electronic work within 90 days of receiving it, you can
receive a refund of the money (if any) you paid for it by sending a
written explanation to the person you received the work from.  If you
received the work on a physical medium, you must return the medium with
your written explanation.  The person or entity that provided you with
the defective work may elect to provide a replacement copy in lieu of a
refund.  If you received the work electronically, the person or entity
providing it to you may choose to give you a second opportunity to
receive the work electronically in lieu of a refund.  If the second copy
is also defective, you may demand a refund in writing without further
opportunities to fix the problem.

1.F.4.  Except for the limited right of replacement or refund set forth
in paragraph 1.F.3, this work is provided to you 'AS-IS' WITH NO OTHER
WARRANTIES OF ANY KIND, EXPRESS OR IMPLIED, INCLUDING BUT NOT LIMITED TO
WARRANTIES OF MERCHANTIBILITY OR FITNESS FOR ANY PURPOSE.

1.F.5.  Some states do not allow disclaimers of certain implied
warranties or the exclusion or limitation of certain types of damages.
If any disclaimer or limitation set forth in this agreement violates the
law of the state applicable to this agreement, the agreement shall be
interpreted to make the maximum disclaimer or limitation permitted by
the applicable state law.  The invalidity or unenforceability of any
provision of this agreement shall not void the remaining provisions.

1.F.6.  INDEMNITY - You agree to indemnify and hold the Foundation, the
trademark owner, any agent or employee of the Foundation, anyone
providing copies of Project Gutenberg-tm electronic works in accordance
with this agreement, and any volunteers associated with the production,
promotion and distribution of Project Gutenberg-tm electronic works,
harmless from all liability, costs and expenses, including legal fees,
that arise directly or indirectly from any of the following which you do
or cause to occur: (a) distribution of this or any Project Gutenberg-tm
work, (b) alteration, modification, or additions or deletions to any
Project Gutenberg-tm work, and (c) any Defect you cause.


Section  2.  Information about the Mission of Project Gutenberg-tm

Project Gutenberg-tm is synonymous with the free distribution of
electronic works in formats readable by the widest variety of computers
including obsolete, old, middle-aged and new computers.  It exists
because of the efforts of hundreds of volunteers and donations from
people in all walks of life.

Volunteers and financial support to provide volunteers with the
assistance they need, are critical to reaching Project Gutenberg-tm's
goals and ensuring that the Project Gutenberg-tm collection will
remain freely available for generations to come.  In 2001, the Project
Gutenberg Literary Archive Foundation was created to provide a secure
and permanent future for Project Gutenberg-tm and future generations.
To learn more about the Project Gutenberg Literary Archive Foundation
and how your efforts and donations can help, see Sections 3 and 4
and the Foundation web page at http://www.pglaf.org.


Section 3.  Information about the Project Gutenberg Literary Archive
Foundation

The Project Gutenberg Literary Archive Foundation is a non profit
501(c)(3) educational corporation organized under the laws of the
state of Mississippi and granted tax exempt status by the Internal
Revenue Service.  The Foundation's EIN or federal tax identification
number is 64-6221541.  Its 501(c)(3) letter is posted at
http://pglaf.org/fundraising.  Contributions to the Project Gutenberg
Literary Archive Foundation are tax deductible to the full extent
permitted by U.S. federal laws and your state's laws.

The Foundation's principal office is located at 4557 Melan Dr. S.
Fairbanks, AK, 99712., but its volunteers and employees are scattered
throughout numerous locations.  Its business office is located at
809 North 1500 West, Salt Lake City, UT 84116, (801) 596-1887, email
business@pglaf.org.  Email contact links and up to date contact
information can be found at the Foundation's web site and official
page at http://pglaf.org

For additional contact information:
     Dr. Gregory B. Newby
     Chief Executive and Director
     gbnewby@pglaf.org


Section 4.  Information about Donations to the Project Gutenberg
Literary Archive Foundation

Project Gutenberg-tm depends upon and cannot survive without wide
spread public support and donations to carry out its mission of
increasing the number of public domain and licensed works that can be
freely distributed in machine readable form accessible by the widest
array of equipment including outdated equipment.  Many small donations
($1 to $5,000) are particularly important to maintaining tax exempt
status with the IRS.

The Foundation is committed to complying with the laws regulating
charities and charitable donations in all 50 states of the United
States.  Compliance requirements are not uniform and it takes a
considerable effort, much paperwork and many fees to meet and keep up
with these requirements.  We do not solicit donations in locations
where we have not received written confirmation of compliance.  To
SEND DONATIONS or determine the status of compliance for any
particular state visit http://pglaf.org

While we cannot and do not solicit contributions from states where we
have not met the solicitation requirements, we know of no prohibition
against accepting unsolicited donations from donors in such states who
approach us with offers to donate.

International donations are gratefully accepted, but we cannot make
any statements concerning tax treatment of donations received from
outside the United States.  U.S. laws alone swamp our small staff.

Please check the Project Gutenberg Web pages for current donation
methods and addresses.  Donations are accepted in a number of other
ways including checks, online payments and credit card donations.
To donate, please visit: http://pglaf.org/donate


Section 5.  General Information About Project Gutenberg-tm electronic
works.

Professor Michael S. Hart is the originator of the Project Gutenberg-tm
concept of a library of electronic works that could be freely shared
with anyone.  For thirty years, he produced and distributed Project
Gutenberg-tm eBooks with only a loose network of volunteer support.


Project Gutenberg-tm eBooks are often created from several printed
editions, all of which are confirmed as Public Domain in the U.S.
unless a copyright notice is included.  Thus, we do not necessarily
keep eBooks in compliance with any particular paper edition.


Most people start at our Web site which has the main PG search facility:

     http://www.gutenberg.org

This Web site includes information about Project Gutenberg-tm,
including how to make donations to the Project Gutenberg Literary
Archive Foundation, how to help produce our new eBooks, and how to
subscribe to our email newsletter to hear about new eBooks.
\end{PGtext}

% %%%%%%%%%%%%%%%%%%%%%%%%%%%%%%%%%%%%%%%%%%%%%%%%%%%%%%%%%%%%%%%%%%%%%%% %
%                                                                         %
% End of the Project Gutenberg EBook of Lectures on Elementary Mathematics, by 
% Joseph Louis Lagrange                                                   %
%                                                                         %
% *** END OF THIS PROJECT GUTENBERG EBOOK LECTURES ON ELEMENTARY MATHEMATICS ***
%                                                                         %
% ***** This file should be named 36640-t.tex or 36640-t.zip *****        %
% This and all associated files of various formats will be found in:      %
%         http://www.gutenberg.org/3/6/6/4/36640/                         %
%                                                                         %
% %%%%%%%%%%%%%%%%%%%%%%%%%%%%%%%%%%%%%%%%%%%%%%%%%%%%%%%%%%%%%%%%%%%%%%% %

\end{document}
###
@ControlwordReplace = (
  ['\\Preface', 'Preface'],
  ['\\Frontispiece', '<Frontispiece>'],
  ['\\Catalog', 'Catalogue of Publications\\nof the\\nOpen Court Publishing Co.'],
  ['\\end{Author}', ''],
  ['\\tb', '-----'],
  ['\\stars', '*    *    *'],
  ['\\ieme', '^{me}'],
  );

@ControlwordArguments = (
  ['\\SetRunningHeads', 1, 0, '', ''],
  ['\\BookMark', 1, 0, '', '', 1, 0, '', ''],
  ['\\Lecture', 0, 0, '', '', 1, 1, 'Lecture ', '', 1, 1, ' ', ''],
  ['\\SectTitle', 1, 1, '', ''],
  ['\\MNote', 1, 0, '', ''],
  ['\\index', 1, 0, '', ''],
  ['\\Appendix', 1, 1, '', ''],
  ['\\BioSketch', 1, 1, '', '', 1, 1, ' ', ''],
  ['\\Signature', 1, 1, '', ' ', 1, 1, '', ''],
  ['\\FrontCatalog', 1, 1, '', ''],
  ['\\Book', 1, 1, '', ''],
  ['\\Title', 0, 0, '', '', 1, 1, '', ' ', 1, 1, '', ''],
  ['\\begin{Author}', 1, 1, '', ''],
  ['\\Item', 1, 1, '', ''],
  ['\\Typo', 1, 0, '', '', 1, 1, '', ''],
  ['\\Add', 1, 1, '', ''],
  ['\\PageSep', 1, 1, '%%-- Page [', ']'],
  ['\\Figure', 1, 1, '<Figure ', '>', 1, 0, '', ''],
  ['\\First', 1, 1, '', '']
  );
###
This is pdfTeXk, Version 3.141592-1.40.3 (Web2C 7.5.6) (format=pdflatex 2010.5.6)  6 JUL 2011 08:11
entering extended mode
 %&-line parsing enabled.
**36640-t.tex
(./36640-t.tex
LaTeX2e <2005/12/01>
Babel <v3.8h> and hyphenation patterns for english, usenglishmax, dumylang, noh
yphenation, arabic, farsi, croatian, ukrainian, russian, bulgarian, czech, slov
ak, danish, dutch, finnish, basque, french, german, ngerman, ibycus, greek, mon
ogreek, ancientgreek, hungarian, italian, latin, mongolian, norsk, icelandic, i
nterlingua, turkish, coptic, romanian, welsh, serbian, slovenian, estonian, esp
eranto, uppersorbian, indonesian, polish, portuguese, spanish, catalan, galicia
n, swedish, ukenglish, pinyin, loaded.
(/usr/share/texmf-texlive/tex/latex/base/book.cls
Document Class: book 2005/09/16 v1.4f Standard LaTeX document class
(/usr/share/texmf-texlive/tex/latex/base/bk12.clo
File: bk12.clo 2005/09/16 v1.4f Standard LaTeX file (size option)
)
\c@part=\count79
\c@chapter=\count80
\c@section=\count81
\c@subsection=\count82
\c@subsubsection=\count83
\c@paragraph=\count84
\c@subparagraph=\count85
\c@figure=\count86
\c@table=\count87
\abovecaptionskip=\skip41
\belowcaptionskip=\skip42
\bibindent=\dimen102
) (/usr/share/texmf-texlive/tex/latex/base/inputenc.sty
Package: inputenc 2006/05/05 v1.1b Input encoding file
\inpenc@prehook=\toks14
\inpenc@posthook=\toks15
(/usr/share/texmf-texlive/tex/latex/base/latin1.def
File: latin1.def 2006/05/05 v1.1b Input encoding file
)) (/usr/share/texmf-texlive/tex/generic/babel/babel.sty
Package: babel 2005/11/23 v3.8h The Babel package
(/usr/share/texmf-texlive/tex/generic/babel/greek.ldf
Language: greek 2005/03/30 v1.3l Greek support from the babel system
(/usr/share/texmf-texlive/tex/generic/babel/babel.def
File: babel.def 2005/11/23 v3.8h Babel common definitions
\babel@savecnt=\count88
\U@D=\dimen103
) Loading the definitions for the Greek font encoding (/usr/share/texmf-texlive
/tex/generic/babel/lgrenc.def
File: lgrenc.def 2001/01/30 v2.2e Greek Encoding
)) (/usr/share/texmf-texlive/tex/generic/babel/english.ldf
Language: english 2005/03/30 v3.3o English support from the babel system
\l@british = a dialect from \language\l@english 
\l@UKenglish = a dialect from \language\l@english 
\l@canadian = a dialect from \language\l@american 
\l@australian = a dialect from \language\l@british 
\l@newzealand = a dialect from \language\l@british 
)) (/usr/share/texmf-texlive/tex/latex/base/ifthen.sty
Package: ifthen 2001/05/26 v1.1c Standard LaTeX ifthen package (DPC)
) (/usr/share/texmf-texlive/tex/latex/amsmath/amsmath.sty
Package: amsmath 2000/07/18 v2.13 AMS math features
\@mathmargin=\skip43
For additional information on amsmath, use the `?' option.
(/usr/share/texmf-texlive/tex/latex/amsmath/amstext.sty
Package: amstext 2000/06/29 v2.01
(/usr/share/texmf-texlive/tex/latex/amsmath/amsgen.sty
File: amsgen.sty 1999/11/30 v2.0
\@emptytoks=\toks16
\ex@=\dimen104
)) (/usr/share/texmf-texlive/tex/latex/amsmath/amsbsy.sty
Package: amsbsy 1999/11/29 v1.2d
\pmbraise@=\dimen105
) (/usr/share/texmf-texlive/tex/latex/amsmath/amsopn.sty
Package: amsopn 1999/12/14 v2.01 operator names
)
\inf@bad=\count89
LaTeX Info: Redefining \frac on input line 211.
\uproot@=\count90
\leftroot@=\count91
LaTeX Info: Redefining \overline on input line 307.
\classnum@=\count92
\DOTSCASE@=\count93
LaTeX Info: Redefining \ldots on input line 379.
LaTeX Info: Redefining \dots on input line 382.
LaTeX Info: Redefining \cdots on input line 467.
\Mathstrutbox@=\box26
\strutbox@=\box27
\big@size=\dimen106
LaTeX Font Info:    Redeclaring font encoding OML on input line 567.
LaTeX Font Info:    Redeclaring font encoding OMS on input line 568.
\macc@depth=\count94
\c@MaxMatrixCols=\count95
\dotsspace@=\muskip10
\c@parentequation=\count96
\dspbrk@lvl=\count97
\tag@help=\toks17
\row@=\count98
\column@=\count99
\maxfields@=\count100
\andhelp@=\toks18
\eqnshift@=\dimen107
\alignsep@=\dimen108
\tagshift@=\dimen109
\tagwidth@=\dimen110
\totwidth@=\dimen111
\lineht@=\dimen112
\@envbody=\toks19
\multlinegap=\skip44
\multlinetaggap=\skip45
\mathdisplay@stack=\toks20
LaTeX Info: Redefining \[ on input line 2666.
LaTeX Info: Redefining \] on input line 2667.
) (/usr/share/texmf-texlive/tex/latex/amsfonts/amssymb.sty
Package: amssymb 2002/01/22 v2.2d
(/usr/share/texmf-texlive/tex/latex/amsfonts/amsfonts.sty
Package: amsfonts 2001/10/25 v2.2f
\symAMSa=\mathgroup4
\symAMSb=\mathgroup5
LaTeX Font Info:    Overwriting math alphabet `\mathfrak' in version `bold'
(Font)                  U/euf/m/n --> U/euf/b/n on input line 132.
)) (/usr/share/texmf-texlive/tex/latex/base/alltt.sty
Package: alltt 1997/06/16 v2.0g defines alltt environment
) (/usr/share/texmf-texlive/tex/latex/tools/array.sty
Package: array 2005/08/23 v2.4b Tabular extension package (FMi)
\col@sep=\dimen113
\extrarowheight=\dimen114
\NC@list=\toks21
\extratabsurround=\skip46
\backup@length=\skip47
) (/usr/share/texmf-texlive/tex/latex/footmisc/footmisc.sty
Package: footmisc 2005/03/17 v5.3d a miscellany of footnote facilities
\FN@temptoken=\toks22
\footnotemargin=\dimen115
\c@pp@next@reset=\count101
\c@@fnserial=\count102
Package footmisc Info: Declaring symbol style bringhurst on input line 817.
Package footmisc Info: Declaring symbol style chicago on input line 818.
Package footmisc Info: Declaring symbol style wiley on input line 819.
Package footmisc Info: Declaring symbol style lamport-robust on input line 823.

Package footmisc Info: Declaring symbol style lamport* on input line 831.
Package footmisc Info: Declaring symbol style lamport*-robust on input line 840
.
) (/usr/share/texmf-texlive/tex/latex/tools/multicol.sty
Package: multicol 2006/05/18 v1.6g multicolumn formatting (FMi)
\c@tracingmulticols=\count103
\mult@box=\box28
\multicol@leftmargin=\dimen116
\c@unbalance=\count104
\c@collectmore=\count105
\doublecol@number=\count106
\multicoltolerance=\count107
\multicolpretolerance=\count108
\full@width=\dimen117
\page@free=\dimen118
\premulticols=\dimen119
\postmulticols=\dimen120
\multicolsep=\skip48
\multicolbaselineskip=\skip49
\partial@page=\box29
\last@line=\box30
\mult@rightbox=\box31
\mult@grightbox=\box32
\mult@gfirstbox=\box33
\mult@firstbox=\box34
\@tempa=\box35
\@tempa=\box36
\@tempa=\box37
\@tempa=\box38
\@tempa=\box39
\@tempa=\box40
\@tempa=\box41
\@tempa=\box42
\@tempa=\box43
\@tempa=\box44
\@tempa=\box45
\@tempa=\box46
\@tempa=\box47
\@tempa=\box48
\@tempa=\box49
\@tempa=\box50
\@tempa=\box51
\c@columnbadness=\count109
\c@finalcolumnbadness=\count110
\last@try=\dimen121
\multicolovershoot=\dimen122
\multicolundershoot=\dimen123
\mult@nat@firstbox=\box52
\colbreak@box=\box53
) (/usr/share/texmf-texlive/tex/latex/base/makeidx.sty
Package: makeidx 2000/03/29 v1.0m Standard LaTeX package
) (/usr/share/texmf-texlive/tex/latex/caption/caption.sty
Package: caption 2007/01/07 v3.0k Customising captions (AR)
(/usr/share/texmf-texlive/tex/latex/caption/caption3.sty
Package: caption3 2007/01/07 v3.0k caption3 kernel (AR)
(/usr/share/texmf-texlive/tex/latex/graphics/keyval.sty
Package: keyval 1999/03/16 v1.13 key=value parser (DPC)
\KV@toks@=\toks23
)
\captionmargin=\dimen124
\captionmarginx=\dimen125
\captionwidth=\dimen126
\captionindent=\dimen127
\captionparindent=\dimen128
\captionhangindent=\dimen129
)) (/usr/share/texmf-texlive/tex/latex/graphics/graphicx.sty
Package: graphicx 1999/02/16 v1.0f Enhanced LaTeX Graphics (DPC,SPQR)
(/usr/share/texmf-texlive/tex/latex/graphics/graphics.sty
Package: graphics 2006/02/20 v1.0o Standard LaTeX Graphics (DPC,SPQR)
(/usr/share/texmf-texlive/tex/latex/graphics/trig.sty
Package: trig 1999/03/16 v1.09 sin cos tan (DPC)
) (/etc/texmf/tex/latex/config/graphics.cfg
File: graphics.cfg 2007/01/18 v1.5 graphics configuration of teTeX/TeXLive
)
Package graphics Info: Driver file: pdftex.def on input line 90.
(/usr/share/texmf-texlive/tex/latex/pdftex-def/pdftex.def
File: pdftex.def 2007/01/08 v0.04d Graphics/color for pdfTeX
\Gread@gobject=\count111
))
\Gin@req@height=\dimen130
\Gin@req@width=\dimen131
) (/usr/share/texmf-texlive/tex/latex/tools/calc.sty
Package: calc 2005/08/06 v4.2 Infix arithmetic (KKT,FJ)
\calc@Acount=\count112
\calc@Bcount=\count113
\calc@Adimen=\dimen132
\calc@Bdimen=\dimen133
\calc@Askip=\skip50
\calc@Bskip=\skip51
LaTeX Info: Redefining \setlength on input line 75.
LaTeX Info: Redefining \addtolength on input line 76.
\calc@Ccount=\count114
\calc@Cskip=\skip52
) (/usr/share/texmf-texlive/tex/latex/fancyhdr/fancyhdr.sty
\fancy@headwidth=\skip53
\f@ncyO@elh=\skip54
\f@ncyO@erh=\skip55
\f@ncyO@olh=\skip56
\f@ncyO@orh=\skip57
\f@ncyO@elf=\skip58
\f@ncyO@erf=\skip59
\f@ncyO@olf=\skip60
\f@ncyO@orf=\skip61
) (/usr/share/texmf-texlive/tex/latex/geometry/geometry.sty
Package: geometry 2002/07/08 v3.2 Page Geometry
\Gm@cnth=\count115
\Gm@cntv=\count116
\c@Gm@tempcnt=\count117
\Gm@bindingoffset=\dimen134
\Gm@wd@mp=\dimen135
\Gm@odd@mp=\dimen136
\Gm@even@mp=\dimen137
\Gm@dimlist=\toks24
(/usr/share/texmf-texlive/tex/xelatex/xetexconfig/geometry.cfg)) (/usr/share/te
xmf-texlive/tex/latex/hyperref/hyperref.sty
Package: hyperref 2007/02/07 v6.75r Hypertext links for LaTeX
\@linkdim=\dimen138
\Hy@linkcounter=\count118
\Hy@pagecounter=\count119
(/usr/share/texmf-texlive/tex/latex/hyperref/pd1enc.def
File: pd1enc.def 2007/02/07 v6.75r Hyperref: PDFDocEncoding definition (HO)
) (/etc/texmf/tex/latex/config/hyperref.cfg
File: hyperref.cfg 2002/06/06 v1.2 hyperref configuration of TeXLive
) (/usr/share/texmf-texlive/tex/latex/oberdiek/kvoptions.sty
Package: kvoptions 2006/08/22 v2.4 Connects package keyval with LaTeX options (
HO)
)
Package hyperref Info: Option `hyperfootnotes' set `false' on input line 2238.
Package hyperref Info: Option `bookmarks' set `true' on input line 2238.
Package hyperref Info: Option `linktocpage' set `false' on input line 2238.
Package hyperref Info: Option `pdfdisplaydoctitle' set `true' on input line 223
8.
Package hyperref Info: Option `pdfpagelabels' set `true' on input line 2238.
Package hyperref Info: Option `bookmarksopen' set `true' on input line 2238.
Package hyperref Info: Option `colorlinks' set `true' on input line 2238.
Package hyperref Info: Hyper figures OFF on input line 2288.
Package hyperref Info: Link nesting OFF on input line 2293.
Package hyperref Info: Hyper index ON on input line 2296.
Package hyperref Info: Plain pages OFF on input line 2303.
Package hyperref Info: Backreferencing OFF on input line 2308.
Implicit mode ON; LaTeX internals redefined
Package hyperref Info: Bookmarks ON on input line 2444.
(/usr/share/texmf-texlive/tex/latex/ltxmisc/url.sty
\Urlmuskip=\muskip11
Package: url 2005/06/27  ver 3.2  Verb mode for urls, etc.
)
LaTeX Info: Redefining \url on input line 2599.
\Fld@menulength=\count120
\Field@Width=\dimen139
\Fld@charsize=\dimen140
\Choice@toks=\toks25
\Field@toks=\toks26
Package hyperref Info: Hyper figures OFF on input line 3102.
Package hyperref Info: Link nesting OFF on input line 3107.
Package hyperref Info: Hyper index ON on input line 3110.
Package hyperref Info: backreferencing OFF on input line 3117.
Package hyperref Info: Link coloring ON on input line 3120.
\Hy@abspage=\count121
\c@Item=\count122
)
*hyperref using driver hpdftex*
(/usr/share/texmf-texlive/tex/latex/hyperref/hpdftex.def
File: hpdftex.def 2007/02/07 v6.75r Hyperref driver for pdfTeX
\Fld@listcount=\count123
)
\TmpLen=\skip62
\@indexfile=\write3
\openout3 = `36640-t.idx'.

Writing index file 36640-t.idx
\c@MNote=\count124
(./36640-t.aux)
\openout1 = `36640-t.aux'.

LaTeX Font Info:    Checking defaults for OML/cmm/m/it on input line 597.
LaTeX Font Info:    ... okay on input line 597.
LaTeX Font Info:    Checking defaults for T1/cmr/m/n on input line 597.
LaTeX Font Info:    ... okay on input line 597.
LaTeX Font Info:    Checking defaults for OT1/cmr/m/n on input line 597.
LaTeX Font Info:    ... okay on input line 597.
LaTeX Font Info:    Checking defaults for OMS/cmsy/m/n on input line 597.
LaTeX Font Info:    ... okay on input line 597.
LaTeX Font Info:    Checking defaults for OMX/cmex/m/n on input line 597.
LaTeX Font Info:    ... okay on input line 597.
LaTeX Font Info:    Checking defaults for U/cmr/m/n on input line 597.
LaTeX Font Info:    ... okay on input line 597.
LaTeX Font Info:    Checking defaults for LGR/cmr/m/n on input line 597.
LaTeX Font Info:    Try loading font information for LGR+cmr on input line 597.

(/usr/share/texmf-texlive/tex/generic/babel/lgrcmr.fd
File: lgrcmr.fd 2001/01/30 v2.2e Greek Computer Modern
)
LaTeX Font Info:    ... okay on input line 597.
LaTeX Font Info:    Checking defaults for PD1/pdf/m/n on input line 597.
LaTeX Font Info:    ... okay on input line 597.
(/usr/share/texmf-texlive/tex/latex/ragged2e/ragged2e.sty
Package: ragged2e 2003/03/25 v2.04 ragged2e Package (MS)
(/usr/share/texmf-texlive/tex/latex/everysel/everysel.sty
Package: everysel 1999/06/08 v1.03 EverySelectfont Package (MS)
LaTeX Info: Redefining \selectfont on input line 125.
)
\CenteringLeftskip=\skip63
\RaggedLeftLeftskip=\skip64
\RaggedRightLeftskip=\skip65
\CenteringRightskip=\skip66
\RaggedLeftRightskip=\skip67
\RaggedRightRightskip=\skip68
\CenteringParfillskip=\skip69
\RaggedLeftParfillskip=\skip70
\RaggedRightParfillskip=\skip71
\JustifyingParfillskip=\skip72
\CenteringParindent=\skip73
\RaggedLeftParindent=\skip74
\RaggedRightParindent=\skip75
\JustifyingParindent=\skip76
)
Package caption Info: hyperref package v6.74m (or newer) detected on input line
 597.
(/usr/share/texmf/tex/context/base/supp-pdf.tex
[Loading MPS to PDF converter (version 2006.09.02).]
\scratchcounter=\count125
\scratchdimen=\dimen141
\scratchbox=\box54
\nofMPsegments=\count126
\nofMParguments=\count127
\everyMPshowfont=\toks27
\MPscratchCnt=\count128
\MPscratchDim=\dimen142
\MPnumerator=\count129
\everyMPtoPDFconversion=\toks28
)
-------------------- Geometry parameters
paper: class default
landscape: --
twocolumn: --
twoside: true
asymmetric: --
h-parts: 9.03374pt, 325.215pt, 9.03375pt
v-parts: 4.15848pt, 495.49379pt, 6.23773pt
hmarginratio: 1:1
vmarginratio: 2:3
lines: --
heightrounded: --
bindingoffset: 0.0pt
truedimen: --
includehead: true
includefoot: true
includemp: --
driver: pdftex
-------------------- Page layout dimensions and switches
\paperwidth  343.28249pt
\paperheight 505.89pt
\textwidth  325.215pt
\textheight 433.62pt
\oddsidemargin  -63.23625pt
\evensidemargin -63.23624pt
\topmargin  -68.11151pt
\headheight 12.0pt
\headsep    19.8738pt
\footskip   30.0pt
\marginparwidth 98.0pt
\marginparsep   7.0pt
\columnsep  10.0pt
\skip\footins  10.8pt plus 4.0pt minus 2.0pt
\hoffset 0.0pt
\voffset 0.0pt
\mag 1000
\@twosidetrue \@mparswitchtrue 
(1in=72.27pt, 1cm=28.45pt)
-----------------------
(/usr/share/texmf-texlive/tex/latex/graphics/color.sty
Package: color 2005/11/14 v1.0j Standard LaTeX Color (DPC)
(/etc/texmf/tex/latex/config/color.cfg
File: color.cfg 2007/01/18 v1.5 color configuration of teTeX/TeXLive
)
Package color Info: Driver file: pdftex.def on input line 130.
)
Package hyperref Info: Link coloring ON on input line 597.
(/usr/share/texmf-texlive/tex/latex/hyperref/nameref.sty
Package: nameref 2006/12/27 v2.28 Cross-referencing by name of section
(/usr/share/texmf-texlive/tex/latex/oberdiek/refcount.sty
Package: refcount 2006/02/20 v3.0 Data extraction from references (HO)
)
\c@section@level=\count130
)
LaTeX Info: Redefining \ref on input line 597.
LaTeX Info: Redefining \pageref on input line 597.
(./36640-t.out) (./36640-t.out)
\@outlinefile=\write4
\openout4 = `36640-t.out'.


Overfull \hbox (14.78989pt too wide) in paragraph at lines 625--625
[]\OT1/cmtt/m/n/8 *** START OF THIS PROJECT GUTENBERG EBOOK LECTURES ON ELEMENT
ARY MATHEMATICS ***[] 
 []

LaTeX Font Info:    Try loading font information for U+msa on input line 627.
(/usr/share/texmf-texlive/tex/latex/amsfonts/umsa.fd
File: umsa.fd 2002/01/19 v2.2g AMS font definitions
)
LaTeX Font Info:    Try loading font information for U+msb on input line 627.
(/usr/share/texmf-texlive/tex/latex/amsfonts/umsb.fd
File: umsb.fd 2002/01/19 v2.2g AMS font definitions
) [1

{/var/lib/texmf/fonts/map/pdftex/updmap/pdftex.map}] [2] [1


]
Underfull \hbox (badness 1097) detected at line 700
\OT1/cmr/m/n/14.4 THE OPEN COURT PUBLISHING COMPANY
 []

<./images/lagrange.jpg, id=103, 104.3097pt x 154.176pt>
File: ./images/lagrange.jpg Graphic file (type jpg)
<use ./images/lagrange.jpg> [2] [3 <./images/lagrange.jpg>] [4

] [5] [6


] [7] [8


] [9]
Overfull \hbox (0.8094pt too wide) in paragraph at lines 886--900
[]\OT1/cmr/m/n/12 But it should never be for-got-ten that the mighty stenophren
ic
 []

[10] [11] [12] [13] [14] [15] (./36640-t.toc [16



] [17] [18] [19])
\tf@toc=\write5
\openout5 = `36640-t.toc'.

[20] [1





] [2] [3] [4] [5] [6] [7] [8] [9] [10] [11] [12] [13] [14] [15] [16] [17] [18] 
[19] [20


] [21] [22] [23] [24] [25] [26] [27] [28] [29] [30] [31] [32] [33] [34] [35] [3
6] [37] [38] [39] [40] [41] [42] [43] [44] [45] [46


] [47] [48] [49] [50] [51] [52] [53] [54] [55] [56] [57] [58] [59] [60] [61] [6
2] [63] [64] [65] [66] [67] [68] [69] [70] [71] [72] [73] [74] [75] [76] [77] [
78] [79] [80] [81] [82] [83] [84] [85] [86] [87


] [88] [89] <./images/fig1.png, id=1073, 334.851pt x 172.9662pt>
File: ./images/fig1.png Graphic file (type png)
<use ./images/fig1.png> [90] [91 <./images/fig1.png (PNG copy)>]
File: ./images/fig1.png Graphic file (type png)
<use ./images/fig1.png> [92] [93] [94] [95] [96] [97] [98] [99] [100] [101] [10
2] [103] [104] [105] [106] [107] [108] [109] [110] <./images/fig2.png, id=1187,
 226.9278pt x 201.8742pt>
File: ./images/fig2.png Graphic file (type png)
<use ./images/fig2.png> [111] [112 <./images/fig2.png (PNG copy)>] [113] [114] 
[115


] [116] [117] [118] [119] [120] [121] [122] <./images/fig3.png, id=1254, 169.59
36pt x 167.6664pt>
File: ./images/fig3.png Graphic file (type png)
<use ./images/fig3.png> [123] <./images/fig4.png, id=1262, 151.767pt x 179.2296
pt>
File: ./images/fig4.png Graphic file (type png)
<use ./images/fig4.png> [124 <./images/fig3.png (PNG copy)>] [125 <./images/fig
4.png (PNG copy)>] <./images/fig5.png, id=1275, 204.765pt x 182.6022pt>
File: ./images/fig5.png Graphic file (type png)
<use ./images/fig5.png> <./images/fig6.png, id=1276, 187.902pt x 71.3064pt>
File: ./images/fig6.png Graphic file (type png)
<use ./images/fig6.png> [126] [127 <./images/fig5.png (PNG copy)>] [128 <./imag
es/fig6.png (PNG copy)>] [129] [130] [131] [132] [133] [134] [135] [136


] [137] (./36640-t.ind [138



] [139] [140] [141] [142] [143] [144]) [145




] [146] [147] [148] [149] [150]
Underfull \hbox (badness 2726) detected at line 7750
\OT1/cmr/m/n/17.28 THE OPEN COURT PUBLISHING CO.,
 []

[151]
Overfull \hbox (6.28976pt too wide) in paragraph at lines 7760--7760
[]\OT1/cmtt/m/n/8 *** END OF THIS PROJECT GUTENBERG EBOOK LECTURES ON ELEMENTAR
Y MATHEMATICS ***[] 
 []

[1


] [2] [3] [4] [5] [6] [7] [8] (./36640-t.aux)

 *File List*
    book.cls    2005/09/16 v1.4f Standard LaTeX document class
    bk12.clo    2005/09/16 v1.4f Standard LaTeX file (size option)
inputenc.sty    2006/05/05 v1.1b Input encoding file
  latin1.def    2006/05/05 v1.1b Input encoding file
   babel.sty    2005/11/23 v3.8h The Babel package
   greek.ldf    2005/03/30 v1.3l Greek support from the babel system
  lgrenc.def    2001/01/30 v2.2e Greek Encoding
 english.ldf    2005/03/30 v3.3o English support from the babel system
  ifthen.sty    2001/05/26 v1.1c Standard LaTeX ifthen package (DPC)
 amsmath.sty    2000/07/18 v2.13 AMS math features
 amstext.sty    2000/06/29 v2.01
  amsgen.sty    1999/11/30 v2.0
  amsbsy.sty    1999/11/29 v1.2d
  amsopn.sty    1999/12/14 v2.01 operator names
 amssymb.sty    2002/01/22 v2.2d
amsfonts.sty    2001/10/25 v2.2f
   alltt.sty    1997/06/16 v2.0g defines alltt environment
   array.sty    2005/08/23 v2.4b Tabular extension package (FMi)
footmisc.sty    2005/03/17 v5.3d a miscellany of footnote facilities
multicol.sty    2006/05/18 v1.6g multicolumn formatting (FMi)
 makeidx.sty    2000/03/29 v1.0m Standard LaTeX package
 caption.sty    2007/01/07 v3.0k Customising captions (AR)
caption3.sty    2007/01/07 v3.0k caption3 kernel (AR)
  keyval.sty    1999/03/16 v1.13 key=value parser (DPC)
graphicx.sty    1999/02/16 v1.0f Enhanced LaTeX Graphics (DPC,SPQR)
graphics.sty    2006/02/20 v1.0o Standard LaTeX Graphics (DPC,SPQR)
    trig.sty    1999/03/16 v1.09 sin cos tan (DPC)
graphics.cfg    2007/01/18 v1.5 graphics configuration of teTeX/TeXLive
  pdftex.def    2007/01/08 v0.04d Graphics/color for pdfTeX
    calc.sty    2005/08/06 v4.2 Infix arithmetic (KKT,FJ)
fancyhdr.sty    
geometry.sty    2002/07/08 v3.2 Page Geometry
geometry.cfg
hyperref.sty    2007/02/07 v6.75r Hypertext links for LaTeX
  pd1enc.def    2007/02/07 v6.75r Hyperref: PDFDocEncoding definition (HO)
hyperref.cfg    2002/06/06 v1.2 hyperref configuration of TeXLive
kvoptions.sty    2006/08/22 v2.4 Connects package keyval with LaTeX options (HO
)
     url.sty    2005/06/27  ver 3.2  Verb mode for urls, etc.
 hpdftex.def    2007/02/07 v6.75r Hyperref driver for pdfTeX
  lgrcmr.fd    2001/01/30 v2.2e Greek Computer Modern
ragged2e.sty    2003/03/25 v2.04 ragged2e Package (MS)
everysel.sty    1999/06/08 v1.03 EverySelectfont Package (MS)
supp-pdf.tex
   color.sty    2005/11/14 v1.0j Standard LaTeX Color (DPC)
   color.cfg    2007/01/18 v1.5 color configuration of teTeX/TeXLive
 nameref.sty    2006/12/27 v2.28 Cross-referencing by name of section
refcount.sty    2006/02/20 v3.0 Data extraction from references (HO)
 36640-t.out
 36640-t.out
    umsa.fd    2002/01/19 v2.2g AMS font definitions
    umsb.fd    2002/01/19 v2.2g AMS font definitions
./images/lagrange.jpg
./images/fig1.png
./images/fig1.png
./images/fig2.png
./images/fig3.png
./images/fig4.png
./images/fig5.png
./images/fig6.png
 36640-t.ind
 ***********

 ) 
Here is how much of TeX's memory you used:
 6322 strings out of 94074
 84976 string characters out of 1165154
 165256 words of memory out of 1500000
 8938 multiletter control sequences out of 10000+50000
 17173 words of font info for 63 fonts, out of 1200000 for 2000
 645 hyphenation exceptions out of 8191
 34i,14n,44p,366b,766s stack positions out of 5000i,500n,6000p,200000b,5000s
</usr/share/texmf-texlive/fonts/type1/bluesky/cm/cmcsc10.pfb></usr/share/texm
f-texlive/fonts/type1/bluesky/cm/cmex10.pfb></usr/share/texmf-texlive/fonts/typ
e1/bluesky/cm/cmmi10.pfb></usr/share/texmf-texlive/fonts/type1/bluesky/cm/cmmi7
.pfb></usr/share/texmf-texlive/fonts/type1/bluesky/cm/cmmi8.pfb></usr/share/tex
mf-texlive/fonts/type1/bluesky/cm/cmr10.pfb></usr/share/texmf-texlive/fonts/typ
e1/bluesky/cm/cmr12.pfb></usr/share/texmf-texlive/fonts/type1/bluesky/cm/cmr17.
pfb></usr/share/texmf-texlive/fonts/type1/bluesky/cm/cmr7.pfb></usr/share/texmf
-texlive/fonts/type1/bluesky/cm/cmr8.pfb></usr/share/texmf-texlive/fonts/type1/
bluesky/cm/cmsl8.pfb></usr/share/texmf-texlive/fonts/type1/bluesky/cm/cmsy10.pf
b></usr/share/texmf-texlive/fonts/type1/bluesky/cm/cmsy7.pfb></usr/share/texmf-
texlive/fonts/type1/bluesky/cm/cmsy8.pfb></usr/share/texmf-texlive/fonts/type1/
bluesky/cm/cmti10.pfb></usr/share/texmf-texlive/fonts/type1/bluesky/cm/cmti12.p
fb></usr/share/texmf-texlive/fonts/type1/bluesky/cm/cmti8.pfb></usr/share/texmf
-texlive/fonts/type1/bluesky/cm/cmtt10.pfb></usr/share/texmf-texlive/fonts/type
1/bluesky/cm/cmtt8.pfb></usr/share/texmf-texlive/fonts/type1/public/cb/grmn1000
.pfb>
Output written on 36640-t.pdf (181 pages, 892352 bytes).
PDF statistics:
 2007 PDF objects out of 2073 (max. 8388607)
 548 named destinations out of 1000 (max. 131072)
 196 words of extra memory for PDF output out of 10000 (max. 10000000)

